\chapter{Parent Hamiltonians}\label{ch:parent_hamiltonian}

This chapter introduces the local ground-space construction and the parent
interaction for translation-invariant periodic MPS.
We follow the standard setup from
\cite{Cirac2021Matrix, PerezGarcia2007Matrix}, with the $N$-site Hilbert space
identified with functions on the computational basis,
\[
    \{0,\ldots,d{-}1\}^N \to \C,
\]
rather than as an explicit tensor product.

\begin{definition}[$N$-site Hilbert space]\label{def:nsite_space_parent}
    \lean{MPSTensor.NSiteSpace}
    \leanok
    For local dimension $d$ and chain length $N$, the $N$-site Hilbert space is
    \[
        \mathcal H_{N}^{(d)} := \{0,\ldots,d{-}1\}^{N} \to \C.
    \]
    This is the computational-basis realization used throughout the chapter for
    periodic-chain states.
\end{definition}

\section{\texorpdfstring{Local ground space $G_L(A)$}{Local ground space G\_L(A)}}\label{sec:ground_space_parent}

Fix an MPS tensor $A = \{A^i\}_{i=0}^{d-1}$ with $A^i \in \MN{D}$ and a block
length $L \ge 1$.
For a word $w = (i_1,\ldots,i_L)$, write
$A^w := A^{i_1}\cdots A^{i_L}$.
The boundary-condition parametrization is
\[
    \Gamma_L : \MN{D} \to (\{0,\ldots,d{-}1\}^L \to \C),
    \qquad
    \Gamma_L(X)(\sigma) := \tr\!(A^{\sigma} X),
\]
where $A^{\sigma}$ abbreviates
$A^{\sigma(0)}\cdots A^{\sigma(L-1)}$.

\begin{definition}[Ground-space map]\label{def:ground_space_map}
    \lean{MPSTensor.groundSpaceMap}
    \leanok
    \uses{def:eval_word}
    The map $\Gamma_L$ above is a $\C$-linear map
    \[
        \Gamma_L : \MN{D} \to (\{0,\ldots,d{-}1\}^L \to \C).
    \]
    In tensor-network notation,
    \begin{center}
        \TNGroundSpaceMap{A}{\sigma_1}{\sigma_L}{L}{X}
    \end{center}
    the upper chain denotes the word tensor $A^\sigma$ with open physical
    indices $\sigma_1,\ldots,\sigma_L$, and the lower red node denotes the boundary
    matrix $X$ inserted on the closing virtual bond before taking the trace.
    The chain length~$L$ is fixed by the surrounding formula, not by an extra
    label inside the diagram.
\end{definition}

\begin{lemma}[Evaluation formula]\label{lem:ground_space_map_apply}
    \lean{MPSTensor.groundSpaceMap_apply}
    \leanok
    \uses{def:ground_space_map}
    For every boundary matrix $X \in \MN{D}$ and every length-$L$ word $\sigma$,
    \[
        \Gamma_{L}(X)(\sigma) = \tr(A^{\sigma} X).
    \]
\end{lemma}

\begin{definition}[Local ground space]\label{def:ground_space_parent}
    \lean{MPSTensor.groundSpace}
    \leanok
    \uses{def:ground_space_map}
    The \emph{local ground space} is the image
    \[
        G_L(A) := \mathrm{im}(\Gamma_L)
        \subseteq (\{0,\ldots,d{-}1\}^L \to \C).
    \]
\end{definition}

\begin{lemma}[Dimension bound]\label{lem:ground_space_finrank_le}
    \lean{MPSTensor.groundSpace_finrank_le}
    \leanok
    \uses{def:ground_space_parent}
    For every $L$,
    \[
        \dim G_L(A) \le D^2.
    \]
\end{lemma}

\begin{proof}
    \leanok
    Since $G_L(A)$ is the range of a linear map from $\MN{D}$,
    \[
        \dim G_L(A) \le \dim \MN{D} = D^2.
    \]
\end{proof}

\begin{lemma}[Ambient-space dimension]\label{lem:nsite_space_finrank}
    \lean{MPSTensor.nSiteSpace_finrank}
    \leanok
    The above Hilbert-space realization satisfies
    \[
        \dim\!(\{0,\ldots,d{-}1\}^L \to \C)=d^L.
    \]
\end{lemma}

\begin{lemma}[Nontriviality criterion]\label{lem:ground_space_ne_top}
    \lean{MPSTensor.groundSpace_ne_top}
    \leanok
    \uses{lem:ground_space_finrank_le, lem:nsite_space_finrank}
    If $d^L > D^2$, then $G_L(A)$ is a proper subspace of the local Hilbert
    space:
    \[
        G_L(A) \neq (\{0,\ldots,d{-}1\}^L \to \C).
    \]
\end{lemma}

\begin{proof}
    \leanok
    If $G_L(A)$ were the whole space, its dimension would be $d^L$ by
    Lemma~\ref{lem:nsite_space_finrank};
    Lemma~\ref{lem:ground_space_finrank_le} gives the contradiction
    $d^L \le D^2$.
\end{proof}

\section{Parent interaction and chain Hamiltonian}\label{sec:parent_interaction_defs}

The canonical parent interaction at range $L$ is
\[
    h_L(A)=1-\Pi_{G_L(A)}=\Pi_{G_L(A)^\perp},
    \qquad
    \ker h_L(A)=G_L(A).
\]
The sources also allow any positive local interaction with this kernel; the
orthogonal projector is the canonical representative used here.

\begin{definition}[Hilbert-space copy of the local ground space]\label{def:ground_space_es}
    \lean{MPSTensor.groundSpaceES}
    \leanok
    \uses{def:nsite_space_parent, def:ground_space_parent}
    Let
    \[
        \iota_L :
        (\{0,\ldots,d{-}1\}^{L}\to\C)
        \simeq \ell^{2}(\{0,\ldots,d{-}1\}^{L})
    \]
    be the canonical computational-basis identification. The Hilbert-space copy
    of the local ground space is
    \[
        G_{L}^{\mathrm{ES}}(A):=\iota_L(G_L(A))
        \subseteq \ell^{2}(\{0,\ldots,d{-}1\}^{L}).
    \]
\end{definition}

\begin{lemma}[Membership in the Hilbert-space copy]\label{lem:mem_ground_space_es}
    \lean{MPSTensor.mem_groundSpaceES_iff}
    \leanok
    \uses{def:ground_space_es, def:ground_space_parent}
    For $v\in\ell^{2}(\{0,\ldots,d{-}1\}^{L})$,
    \[
        v\in G_L^{\mathrm{ES}}(A)
        \Longleftrightarrow
        \iota_L^{-1}v\in G_L(A).
    \]
\end{lemma}

\begin{proof}
    \leanok
    This is the defining equation
    $G_L^{\mathrm{ES}}(A)=\iota_L(G_L(A))$.
\end{proof}

\begin{definition}[Canonical parent interaction]\label{def:parent_interaction}
    \lean{MPSTensor.parentInteraction}
    \leanok
    \uses{def:ground_space_parent}
    The parent interaction at range $L$ is the canonical orthogonal projector
    \[
        h_L(A) := \Pi_{G_L(A)^\perp}=1-\Pi_{G_L(A)},
        \qquad
        \ker h_L(A)=G_L(A),
    \]
    on the local $L$-site space.
\end{definition}

\begin{lemma}[Parent interaction annihilates ground-space elements]%
    \label{lem:parent_interaction_apply_mem_ground_space}
    \lean{MPSTensor.parentInteraction_apply_mem_groundSpace}
    \leanok
    \uses{def:parent_interaction, def:ground_space_parent}
    For any $v \in G_L(A)$, we have $h_L(A)\,v = 0$.
\end{lemma}

\begin{proof}
    \leanok
    Since $h_L(A)$ is the orthogonal projector onto $G_L(A)^\perp$, it
    annihilates every element of $G_L(A)$.
\end{proof}

\begin{definition}[Restriction to a cyclic interval]\label{def:extract_window}
    \lean{MPSTensor.extractWindow}
    \leanok
    \uses{def:nsite_space_parent}
    For $\sigma \in \{0,\ldots,d{-}1\}^{N}$ and a starting site $i$, write
    \[
        \sigma|_{[i,i+L-1]}
        \in \{0,\ldots,d{-}1\}^{L}
    \]
    for the word defined by
    \[
        \bigl(\sigma|_{[i,i+L-1]}\bigr)(r)
        =\sigma(i+r \bmod N),
        \qquad 0\le r<L.
    \]
\end{definition}

\begin{definition}[Prescribing a cyclic interval]\label{def:replace_window}
    \lean{MPSTensor.replaceWindow}
    \leanok
    \uses{def:extract_window}
    If $L \le N$ and $\tau\in\{0,\ldots,d{-}1\}^{L}$, write
    \[
        \sigma^{[i,i+L-1]\leftarrow\tau}
    \]
    for the configuration defined by
    \[
        \sigma^{[i,i+L-1]\leftarrow\tau}(k)
        =
        \begin{cases}
            \tau(r),
                & k\equiv i+r \pmod N\text{ for some }0\le r<L,\\
            \sigma(k), & \text{otherwise}.
        \end{cases}
    \]
    The index $r$ in the first case is unique because $L\le N$.
\end{definition}

\begin{lemma}[Extracting a replaced window]\label{lem:extract_window_replace_window}
    \lean{MPSTensor.extractWindow_replaceWindow}
    \leanok
    \uses{def:extract_window, def:replace_window}
    If $L \le N$, then restricting to the same cyclic interval after prescribing
    its values recovers the inserted word:
    \[
        \left(\sigma^{[i,i+L-1]\leftarrow\tau}\right)|_{[i,i+L-1]}=\tau.
    \]
\end{lemma}

\begin{lemma}[Replacing an extracted window]\label{lem:replace_window_extract_window}
    \lean{MPSTensor.replaceWindow_extractWindow}
    \leanok
    \uses{def:extract_window, def:replace_window}
    If $L \le N$, then prescribing on a cyclic interval the values already
    carried by $\sigma$ leaves $\sigma$ unchanged:
    \[
        \sigma^{[i,i+L-1]\leftarrow\sigma|_{[i,i+L-1]}}=\sigma.
    \]
\end{lemma}
\begin{proof}\leanok
    Case-split on whether the offset $(k - i + N) \bmod N$ lies in $[0,L)$:
    inside the cyclic interval the prescribed value is the original value of
    $\sigma$; outside the interval both sides already agree with $\sigma$.
\end{proof}

For a periodic chain of length $N$, the parent Hamiltonian is the sum of
translated copies of the local interaction.

\begin{definition}[Translated local term]\label{def:local_term_parent}
    \lean{MPSTensor.localTerm}
    \leanok
    \uses{def:parent_interaction, def:replace_window, def:extract_window}
    For each site index $i \in \{0,\ldots,N{-}1\}$, the translated local term is
    determined by
    \[
        (h_i\psi)(\sigma)
        =
        \begin{cases}
        \left[
            h_L(A)
            \bigl(\tau\mapsto
                \psi(\sigma^{[i,i+L-1]\leftarrow\tau})\bigr)
        \right]\!\bigl(\sigma|_{[i,i+L-1]}\bigr), & L \le N,\\
        0, & L > N.
        \end{cases}
    \]
\end{definition}

\begin{definition}[Parent Hamiltonian]\label{def:parent_hamiltonian}
    \lean{MPSTensor.parentHamiltonian}
    \leanok
    \uses{def:local_term_parent}
    The chain Hamiltonian is
    \[
        H_N(A,L) := \sum_{i=0}^{N-1} h_i.
    \]
\end{definition}

\begin{definition}[Frustration-free condition]\label{def:is_frustration_free}
    \lean{MPSTensor.IsFrustrationFree}
    \leanok
    \uses{def:local_term_parent}
    A vector $\psi$ is a frustration-free ground state of the parent Hamiltonian
    when each local interaction annihilates it:
    \[
        \forall\, i \in \{0,\ldots,N{-}1\},\quad h_i \psi = 0.
    \]
\end{definition}

\begin{lemma}[MPV window membership]\label{lem:mpv_window_mem_ground_space}
    \lean{MPSTensor.mpv_window_mem_groundSpace}
    \leanok
    \uses{def:ground_space_parent, def:mpv}
    For any window of $L$ consecutive sites starting at position~$i$ in an
    $N$-site periodic chain ($L \le N$), the restriction of the MPV state to
    that window lies in $G_L(A)$.
\end{lemma}

\begin{proof}
    \leanok
    The boundary matrix is the product of $A$-matrices on the complement sites.
    Trace cyclicity rotates the full $N$-site product so that window
    indices come first, matching the ground-space map definition.
\end{proof}

\begin{lemma}[Local term annihilates the MPV]\label{lem:local_term_annihilates_mpv}
    \lean{MPSTensor.localTerm_annihilates_mpv}
    \leanok
    \uses{def:local_term_parent, def:mpv, lem:mpv_window_mem_ground_space,
        lem:parent_interaction_apply_mem_ground_space}
    For each site $i$, $h_i\,V^{(N)}(A) = 0$.
\end{lemma}

\begin{proof}
    \leanok
    Combines MPV window membership with the fact that the parent
    interaction annihilates ground-space elements.
\end{proof}

\begin{lemma}[Parent Hamiltonian annihilates the MPV]\label{lem:parent_hamiltonian_annihilates}
    \lean{MPSTensor.parentHamiltonian_annihilates}
    \leanok
    \uses{def:parent_hamiltonian, def:mpv, lem:local_term_annihilates_mpv}
    For every $N$, the MPV state satisfies
    \[
        H_N(A,L)\,V^{(N)}(A) = 0.
    \]
\end{lemma}

\begin{proof}
    \leanok
    Sum of zeros, since each local term annihilates the MPV.
\end{proof}

\begin{lemma}[Frustration-freeness of the MPV]\label{lem:parent_hamiltonian_ff}
    \lean{MPSTensor.parentHamiltonian_frustrationFree}
    \leanok
    \uses{def:is_frustration_free, def:mpv, lem:local_term_annihilates_mpv}
    The MPS vector is a frustration-free ground state: for every site $i$,
    \[
        h_i V^{(N)}(A)=0.
    \]
\end{lemma}

\begin{proof}
    \leanok
    Immediate from Lemma~\ref{lem:local_term_annihilates_mpv}.
\end{proof}

\section{Intersection property and unique ground state}\label{sec:intersection_unique}

This section develops the intersection property of local ground spaces for
injective MPS and the resulting uniqueness of the ground state on the periodic
chain.

\subsection{Restriction maps}

Given a state $\psi$ on $L+1$ sites, we define two restriction maps.

\begin{remark}[Word and restriction identities]
    \lean{List.ofFn_snoc}
    \lean{MPSTensor.evalWord_ofFn_snoc}
    \lean{MPSTensor.evalWord_ofFn_cons}
    \lean{MPSTensor.restrictLastₗ}
    \lean{MPSTensor.restrictFirstₗ}
    \lean{MPSTensor.restrictLast_apply}
    \lean{MPSTensor.restrictFirst_apply}
    \leanok
    Appending a last letter to a word corresponds to right multiplication by the
    corresponding tensor letter, while prepending a first letter corresponds to left
    multiplication. The restriction maps are the associated linear maps on the
    Hilbert spaces in the computational-basis realization, together with their
    pointwise formulas.
\end{remark}

\begin{definition}[Left restriction]\label{def:restrict_last}
    \lean{MPSTensor.restrictLast}
    \leanok
    Fix the last physical index $j$:
    \[
        \sigma \longmapsto \psi(\sigma_1,\ldots,\sigma_L,j).
    \]
\end{definition}

\begin{definition}[Right restriction]\label{def:restrict_first}
    \lean{MPSTensor.restrictFirst}
    \leanok
    Fix the first physical index $i$:
    \[
        \sigma \longmapsto \psi(i,\sigma_1,\ldots,\sigma_L).
    \]
\end{definition}

\begin{definition}[Left ground membership]\label{def:in_left_ground}
    \lean{MPSTensor.InLeftGround}
    \leanok
    \uses{def:restrict_last, def:ground_space_parent}
    A state $\psi$ on $L+1$ sites satisfies the left ground condition if
    for every $j$, the state obtained by fixing the last index to $j$ lies in
    $G_L(A)$.
\end{definition}

\begin{definition}[Right ground membership]\label{def:in_right_ground}
    \lean{MPSTensor.InRightGround}
    \leanok
    \uses{def:restrict_first, def:ground_space_parent}
    A state $\psi$ on $L+1$ sites satisfies the right ground condition if
    for every $i$, the state obtained by fixing the first index to $i$ lies in
    $G_L(A)$.
\end{definition}

\subsection{Forward direction: ground space restricts}

\begin{lemma}[Left restriction]\label{lem:ground_space_in_left_ground}
    \lean{MPSTensor.groundSpace_inLeftGround}
    \leanok
    \uses{def:ground_space_parent, def:in_left_ground}
    If $\psi \in G_{L+1}(A)$, i.e.\ $\psi(\sigma) = \tr(A^\sigma X)$ for some~$X$,
    then for each $j$, the state obtained by fixing the last index to $j$
    lies in $G_L(A)$ with witness $A^j X$.
\end{lemma}

\begin{proof}
    \leanok
    The restriction simply absorbs the last tensor into the boundary matrix:
    \begin{center}
        \TNGroundSpaceMap{A}{\sigma_1}{\sigma_L}{L}{A^j X}
    \end{center}
    Direct computation:
    $\psi(\sigma_1,\ldots,\sigma_L,j) = \tr(A^{\sigma_1}\cdots A^{\sigma_L} \cdot A^j \cdot X)
        = \Gamma_L(A^j X)(\sigma)$.
\end{proof}

\begin{lemma}[Right restriction]\label{lem:ground_space_in_right_ground}
    \lean{MPSTensor.groundSpace_inRightGround}
    \leanok
    \uses{def:ground_space_parent, def:in_right_ground}
    If $\psi \in G_{L+1}(A)$ with $\psi(\sigma) = \tr(A^\sigma X)$, then
    for each $i$, the state obtained by fixing the first index to $i$
    lies in $G_L(A)$ with witness $X A^i$.
\end{lemma}

\begin{proof}
    \leanok
    On the right restriction one first rotates the trace and then reads the
    resulting boundary matrix on the shortened window:
    \begin{center}
        \TNGroundSpaceMap{A}{\sigma_1}{\sigma_L}{L}{X A^i}
    \end{center}
    By trace cyclicity:
    $\psi(i,\sigma_1,\ldots,\sigma_L) = \tr(A^i A^{\sigma} X)
        = \tr(A^{\sigma} \cdot X A^i) = \Gamma_L(X A^i)(\sigma)$.
\end{proof}

\subsection{Injectivity of the ground-space map}

\begin{theorem}[Ground-space map injectivity]\label{thm:ground_space_map_injective}
    \lean{MPSTensor.groundSpaceMap_injective}
    \leanok
    \uses{def:ground_space_map, def:injective}
    If $A$ is injective and $L \ge 1$, then $\Gamma_L$ is injective.
\end{theorem}

\begin{proof}
    \leanok
    If $\Gamma_L(X) = 0$ then $\tr(A^\sigma X) = 0$ for every length-$L$ word
    $\sigma$. Since $A$ is injective, the products $\{A^\sigma\}_\sigma$ span
    $\MN{D}$, so the trace pairing gives $X = 0$.
\end{proof}

\begin{theorem}[Ground-space map injectivity from word span]
    \label{thm:ground_space_map_injective_word_span}
    \lean{MPSTensor.groundSpaceMap_injective_of_wordSpan_eq_top}
    \leanok
    \uses{def:ground_space_map}
    If the length-$L$ products $\{A^\sigma : |\sigma| = L\}$ span $\MN{D}$, then
    $\Gamma_L$ is injective.
\end{theorem}

\begin{proof}
    \leanok
    If $\Gamma_L(X) = 0$, then $\tr(A^\sigma X) = 0$ for every word $\sigma$ of
    length $L$. The spanning hypothesis turns this into $\tr(MX) = 0$ for every
    $M \in \MN{D}$, and the trace pairing gives $X = 0$.
\end{proof}

\begin{lemma}[Ground-space map injectivity after blocking]
    \label{lem:ground_space_map_injective_blk}
    \lean{MPSTensor.groundSpaceMap_injective_of_isNBlkInjective}
    \leanok
    \uses{thm:ground_space_map_injective_word_span, def:l_blk_injective}
    If the length-$L_0$ products span the full matrix algebra,
    \[
        \spn\{A^\omega:|\omega|=L_0\}=\MN{D},
    \]
    then $\Gamma_{L_0}$ is injective.
\end{lemma}

\begin{proof}
    \leanok
    The hypothesis is precisely the fixed-length spanning condition needed by
    Theorem~\ref{thm:ground_space_map_injective_word_span}.
\end{proof}

\begin{theorem}[Ground-space dimension]\label{thm:ground_space_finrank_eq}
    \lean{MPSTensor.groundSpace_finrank_eq}
    \leanok
    \uses{thm:ground_space_map_injective, lem:ground_space_finrank_le}
    If $A$ is injective and $L \ge 1$, then $\dim G_L(A) = D^2$.
\end{theorem}

\begin{proof}
    \leanok
    The upper bound $\dim G_L(A) \le D^2$ is Lemma~\ref{lem:ground_space_finrank_le}.
    Injectivity of $\Gamma_L$ gives $\dim G_L(A) = \dim \MN{D} = D^2$.
\end{proof}

\subsection{The intersection property}

\begin{theorem}[One inclusion of the intersection property]\label{thm:ground_space_intersection}
    \lean{MPSTensor.groundSpace_intersection}
    \leanok
    \uses{thm:ground_space_map_injective, def:in_left_ground, def:in_right_ground,
        def:ground_space_parent, def:injective}
    If $A$ is injective and $L \ge 2$, then the nontrivial inclusion in the
    intersection property holds:
    \[
        G_L^{\mathrm{left}}\cap G_L^{\mathrm{right}}
        \subseteq
        G_{L+1}(A).
    \]
\end{theorem}

\begin{proof}
    \leanok
    The ``inverting and growing back'' argument compares the left and right
    restrictions on their common $(L-1)$-site overlap and then reconstructs the
    full $(L+1)$-site word from a single boundary matrix:
    \begin{center}
        \TNBoundaryRegrow{X'}{\sigma_1}{\sigma_{L+1}}{L+1}
    \end{center}
    \begin{enumerate}
        \item From the right ground condition, for each physical index~$i$ there
              exists a unique $Y_i \in \MN{D}$ (by injectivity of $\Gamma_L$) such
              that $\psi(i,\sigma) = \tr(A^\sigma Y_i)$.
        \item Similarly, from the left ground condition, for each~$j$ there exists
              a unique $Z_j$ with $\psi(\sigma,j) = \tr(A^\sigma Z_j)$.
        \item Matching on the $(L{-}1)$-site overlap and using nondegeneracy of
              the trace pairing:
              $A^j Y_i = Z_j A^i$ for all $i,j$.
        \item Since $A$ is injective, the matrices $\{A^j\}$ span $\MN{D}$
              (Definition~\ref{def:injective}), so $\Id = \sum_j c_j A^j$ for
              some coefficients $c_j \in \C$. Then
              $Y_i = \Id \cdot Y_i = \sum_j c_j A^j Y_i = \sum_j c_j Z_j A^i = X' A^i$,
              where $X' := \sum_j c_j Z_j$.
        \item By trace cyclicity,
              $\psi(\sigma) = \tr(A^\sigma X')$, so $\psi \in G_{L+1}(A)$.
    \end{enumerate}
\end{proof}

\begin{theorem}[Intersection property]\label{thm:ground_space_iff}
    \lean{MPSTensor.groundSpace_iff_left_right}
    \leanok
    \uses{lem:ground_space_in_left_ground, lem:ground_space_in_right_ground,
        thm:ground_space_intersection}
    For injective $A$ and $L \ge 2$:
    \[
        G_L^{\mathrm{left}}\cap G_L^{\mathrm{right}}
        =
        G_{L+1}(A),
    \]
    equivalently
    $\psi \in G_{L+1}(A) \iff
      \psi \in G_L^{\mathrm{left}} \cap G_L^{\mathrm{right}}$.
\end{theorem}

\begin{proof}
    \leanok
    Immediate from the two forward-direction lemmas and
    Theorem~\ref{thm:ground_space_intersection}.
\end{proof}

\subsection{Unique ground state}

\subsubsection{Contiguous and cyclic interval restrictions}

\begin{remark}[Contiguous interval restrictions]\label{rem:contiguous_window_operators}
    \lean{MPSTensor.contiguousCfg}
    \lean{MPSTensor.contiguousRestrictₗ}
    \lean{MPSTensor.contiguousRestrictₗ_apply}
    \lean{MPSTensor.contiguousCfg_zero_full}
    \lean{MPSTensor.contiguousRestrictₗ_restrictLast}
    \lean{MPSTensor.contiguousRestrictₗ_restrictFirst}
    \lean{MPSTensor.tailRestrictₗ_contiguousRestrictₗ}
    \leanok
    For a non-wrapping interval $[s,s+M-1] \subseteq \{0,\ldots,N-1\}$ and a
    full configuration $\rho$ supplying the complementary values, the
    restriction map is
    \[
        (\operatorname{Res}^{\rho}_{s,M}\psi)(\omega)
        =
        \psi\!\left(\rho^{[s,s+M-1]\leftarrow\omega}\right).
    \]
    The listed identities are the equations for $s=0$, $M=N$, for deleting the
    first or last site, and for first fixing $K$ sites and then restricting to
    $[s+K,s+K+L-1]$.
\end{remark}

\begin{remark}[Cyclic interval restrictions]\label{rem:cyclic_window_operators}
    \lean{MPSTensor.cyclicCfg}
    \lean{MPSTensor.cyclicRestrictₗ}
    \lean{MPSTensor.cyclicRestrictₗ_apply}
    \lean{MPSTensor.cyclicRestrictₗ_congr_outside}
    \lean{MPSTensor.cyclicRestrictₗ_cyclicCfg_outside}
    \lean{MPSTensor.cyclicCfg_cyclicForwardSite}
    \lean{MPSTensor.eq_cyclic_site_of_offset_eq}
    \lean{MPSTensor.cyclicRestrictₗ_restrictLast}
    \lean{MPSTensor.cyclicRestrictₗ_restrictFirst}
    \lean{MPSTensor.cyclicCfg_eq_contiguousCfg}
    \lean{MPSTensor.cyclicRestrictₗ_eq_contiguousRestrictₗ}
    \leanok
    On the periodic chain, for a full configuration $\rho$ supplying the
    complementary values, put
    \[
        \delta_i(k) := (k+N-i)\bmod N,\qquad
        \rho^{\mathrm{cyc}}_{i,L,\omega}(k)
        :=
        \begin{cases}
        \omega(\delta_i(k)), & \delta_i(k)<L,\\
        \rho(k), & \delta_i(k)\ge L.
        \end{cases}
    \]
    The restriction to the cyclic interval beginning at $i$ is
    \[
        (\operatorname{Res}^{\rho}_{i,L}\psi)(\omega)
        =
        \psi\!\left(\rho^{\mathrm{cyc}}_{i,L,\omega}\right),
    \]
    with all site labels read modulo $N$. If $L\le N$, this is the same
    configuration as $\rho^{[i,i+L-1]\leftarrow\omega}$. If two boundary
    conditions agree whenever $\delta_i(k)\ge L$, their restrictions are equal.
    In particular, filling the interval inside the outside configuration leaves the
    restriction unchanged:
    \[
        \operatorname{Res}^{\rho^{\mathrm{cyc}}_{i,L,\omega}}_{i,L}\psi
        =
        \operatorname{Res}^{\rho}_{i,L}\psi.
    \]
    For every $r<L$, the inserted word is recovered:
    \[
        \rho^{\mathrm{cyc}}_{i,L,\omega}(i+r)=\omega(r).
    \]
    Here site labels are taken modulo $N$. The remaining listed identities say that
    deleting the final site of an $(L+1)$-interval gives the $L$-interval
    beginning at $i$; if $L+1\le N$, deleting the first site gives the
    $L$-interval beginning at $i+1$; and the cyclic formula agrees with the
    contiguous formula when the interval does not wrap around the ring.
\end{remark}

\begin{lemma}[One-site restrictions of cyclic-window boundary conditions]
    \label{lem:cyclic_window_boundary_matrix_transport}
    \lean{MPSTensor.cyclicRestrictₗ_restrictFirst_groundSpaceMap}
    \lean{MPSTensor.cyclicRestrictₗ_restrictLast_groundSpaceMap}
    \lean{MPSTensor.adjacent_cyclicRestrictₗ_witness_overlap}
    \lean{MPSTensor.adjacent_cyclicRestrictₗ_witness_overlap_common_background}
    \lean{MPSTensor.boundary_witness_product_of_adjacent_overlaps}
    \lean{MPSTensor.adjacent_cyclicRestrictₗ_witness_product}
    \lean{MPSTensor.adjacent_cyclicRestrictₗ_witness_product_common_background}
    \lean{MPSTensor.adjacent_cyclicRestrictₗ_witness_product_common_background_named}
    \leanok
    \uses{rem:cyclic_window_operators, def:ground_space_parent, def:l_blk_injective}
    If a cyclic $(L+1)$-window is represented by a boundary matrix $Y$, then
    deleting one endpoint gives
    \[
        \Gamma_{L+1}(Y)\big|_{\text{last}=b}=\Gamma_L(A^bY),
        \qquad
        \Gamma_{L+1}(Y)\big|_{\text{first}=a}=\Gamma_L(YA^a).
    \]
    Hence two adjacent cyclic windows whose restrictions agree on the common
    length-$L$ overlap have boundary matrices satisfying
    \[
        Y_1 A^a = A^b Y_2.
    \]
    Consequently, for a finite chain of adjacent overlaps indexed by
    \(0\le r<n\),
    \[
        Y_rA^{a_r}=A^{b_r}Y_{r+1},
    \]
    one has
    \[
        Y_0A^{a_0}\cdots A^{a_{n-1}}
        =
        A^{b_0}\cdots A^{b_{n-1}}Y_n .
    \]
    In particular, if the two endpoint words are named by the equations
    \[
        a_r=\rho(i_0+r),\qquad
        b_r=\rho(i_0+r+L+1),
    \]
    with site labels read modulo $N$, then the same transport identity is
    written as
    \[
        Y_0A^a=A^bY_n .
    \]
\end{lemma}

\begin{proof}
    \leanok
    The first two identities are the cyclic first- and last-site deletion
    identities, followed by the corresponding ground-space restriction formulas.
    For adjacent windows, equality of the completed outside configurations makes
    the two restricted states on the common overlap equal. Injectivity of the
    length-$L$ ground-space map then identifies their boundary matrices.
    Iterating the displayed adjacent identity gives the word-product identity.
\end{proof}

\begin{lemma}[Iterated contiguous-window intersection]\label{lem:contiguous_mem_ground_space}
    \lean{MPSTensor.contiguous_mem_groundSpace}
    \leanok
    \uses{thm:ground_space_intersection}
    If an $N$-site state satisfies the $L$-site ground condition on every
    non-wrapping contiguous window (with $L \ge 2$), then it lies in $G_N(A)$.
\end{lemma}

\begin{proof}
    \leanok
    The induction repeatedly merges two neighbouring admissible windows
    $[i,i+L-1]$ and $[i+1,i+L]$ into the larger window $[i,i+L]$ by applying
    Theorem~\ref{thm:ground_space_intersection} to their common overlap
    $[i+1,i+L-1]$.
    Strong induction on window size, combining adjacent windows via
    Theorem~\ref{thm:ground_space_intersection} at each step.
\end{proof}

\begin{definition}[Suffix restriction]\label{def:tail_restrict_parent}
    \lean{MPSTensor.tailRestrictₗ}
    \leanok
    Fix a prefix $u : \Fin K \to \Fin d$. The suffix restriction of an
    $(K+L)$-site state $\psi$ is the $L$-site state given by
    \[
        (\operatorname{tailRestrict}_u \psi)(\sigma) = \psi(u,\sigma).
    \]
\end{definition}

\begin{definition}[Suffix ground membership]\label{def:in_tail_ground}
    \lean{MPSTensor.InTailGround}
    \leanok
    \uses{def:tail_restrict_parent, def:ground_space_parent}
    A state $\psi$ on $K+L$ sites lies in the suffix ground condition if every
    fixed prefix $u$ gives a suffix state in $G_L(A)$.
\end{definition}

\begin{lemma}[Ground-space suffix restriction]\label{lem:ground_space_in_tail_ground}
    \lean{MPSTensor.restrictLast_groundSpaceMap,
        MPSTensor.restrictFirst_groundSpaceMap,
        MPSTensor.tailRestrictₗ_groundSpaceMap,
        MPSTensor.groundSpace_inTailGround}
    \leanok
    \uses{def:in_tail_ground, def:ground_space_parent}
    If $\psi \in G_{K+L}(A)$, then every fixed-prefix suffix restriction of
    $\psi$ lies in $G_L(A)$. For a vector in $G_{L+1}(A)$, fixing either
    endpoint gives the expected left or right multiplication of its boundary
    matrix.
\end{lemma}

\begin{proof}
    \leanok
    Write $\psi(\rho) = \tr(A^\rho X)$ for a boundary matrix $X$. Fixing the
    final site gives the boundary matrix $A^jX$, while fixing the first site
    gives $XA^i$ by cyclicity of the trace. Fixing a prefix $u$ gives
    \[
        \psi(u,\sigma) = \tr(A^u A^\sigma X) = \tr(A^\sigma X A^u),
    \]
    so the suffix restriction again has the form of a ground-space vector.
\end{proof}

\begin{lemma}[Left-tail compatibility]\label{lem:left_tail_compatibility}
    \lean{MPSTensor.exists_left_tail_compatibility}
    \leanok
    \uses{def:in_left_ground, def:in_tail_ground, def:ground_space_parent}
    Assume $\psi$ on $K+L_0+1$ sites satisfies the left ground condition on the
    first $K+L_0$ sites and the suffix ground condition on every fixed prefix of
    length $K$. Then there exist matrices $(Z_j)_j$ and $(Y_u)_u$ such that
    \[
        \psi(\sigma,j) = \tr(A^\sigma Z_j), \qquad
        \psi(u,\tau) = \tr(A^\tau Y_u),
    \]
    and for all $j$ and $u$,
    \[
        Z_j A^u = A_j Y_u.
    \]
\end{lemma}

\begin{proof}
    \leanok
    Choose $Z_j$ from the long left-window condition and $Y_u$ from the suffix
    ground condition. Restricting the suffix state at its last site gives the
    boundary matrix $A_j Y_u$, while first fixing the last site and then taking
    the suffix gives $Z_j A^u$. The two restrictions are the same because
    appending the fixed prefix and then the last letter gives the same
    configuration as appending the last letter to the suffix and then placing
    the prefix in front. Injectivity of $\Gamma_{L_0}$, obtained from
    $L_0$-block injectivity, yields $Z_j A^u = A_j Y_u$.
\end{proof}

\begin{theorem}[Inverting step for word compatibility]%
    \label{thm:right_factor_word_compatibility}
    \lean{MPSTensor.exists_right_factor_of_evalWord_compatibility}
    \leanok
    \uses{def:l_blk_injective}
    Assume $A$ is $L_0$-block injective with $L_0 > 0$. If matrices $(Z_j)_j$
    satisfy
    \[
        Z_j A^\sigma = A_j Y_\sigma
    \]
    for every word $\sigma$ of length $K$, then there exists a matrix $X$ such
    that $Z_j = A_j X$ for all $j$.
\end{theorem}

\begin{proof}
    \leanok
    For $K = 1$, the hypothesis gives $Z_j A_i = A_j Y_i$ for every letter $i$.
    Every blocked word of length $L_0$ is nonempty, so this identity extends to
    the coefficients of the blocked tensor $A^{[L_0]}$. Since $A^{[L_0]}$ is
    injective, the identity matrix is a linear combination of those blocked
    coefficients, and substituting this decomposition yields $Z_j = A_j X$ for a
    common matrix $X$. For larger $K$, strip the first letter: for each $i$, the
    matrices $Z_j A_i$ satisfy the same compatibility hypothesis with words of
    length $K - 1$, so the induction hypothesis reduces the problem to the case
    $K = 1$.
\end{proof}

\begin{theorem}[Growing-back step at the injectivity length]%
    \label{thm:ground_space_extend_right_block_injective}
    \lean{MPSTensor.groundSpace_extend_right_of_isNBlkInjective}
    \leanok
    \uses{def:in_left_ground, def:in_tail_ground, def:ground_space_parent}
    Assume $A$ is $L_0$-block injective with $L_0 > 0$. If an element $\psi$
    on $K + L_0 + 1$ sites satisfies the left ground condition on
    the first $K + L_0$ sites and the suffix ground condition on every prefix of
    length $K$, then $\psi \in G_{K + L_0 + 1}(A)$.
\end{theorem}

\begin{proof}
    \leanok
    \uses{lem:left_tail_compatibility, thm:right_factor_word_compatibility}
    Lemma~\ref{lem:left_tail_compatibility} produces matrices $(Z_j)_j$ and
    $(Y_u)_u$ with $Z_j A^u = A_j Y_u$ for every prefix word $u$ of length $K$.
    Theorem~\ref{thm:right_factor_word_compatibility} gives a common right factor
    $X$ such that $Z_j = A_j X$ for all $j$. Therefore each last-site restriction
    of $\psi$ agrees with the corresponding last-site restriction of the ground-
    space vector determined by $X$. Since every configuration is obtained by
    adjoining its final letter to its initial $(K + L_0)$-tuple, the two states
    coincide, so $\psi$ lies in $G_{K + L_0 + 1}(A)$.
\end{proof}

\begin{theorem}[Open-chain iteration of the intersection property]%
    \label{thm:normal_open_chain_range_reduction}
    \lean{MPSTensor.contiguous_mem_groundSpace_of_isNBlkInjective}
    \leanok
    \uses{thm:ground_space_extend_right_block_injective, def:l_blk_injective,
        rem:contiguous_window_operators}
    Let $A$ be $L_0$-block injective with $D \ge 1$ and $L_0 > 0$.
    If an $N$-site state satisfies the ground-space constraint on every
    non-wrapping contiguous interval of length $L_0+1$, for every fixed choice
    of the complementary sites, with $N \ge L_0+1$, then it lies in $G_N(A)$.
\end{theorem}

\begin{proof}
    \leanok
    \uses{thm:ground_space_extend_right_block_injective,
        def:tail_restrict_parent, rem:contiguous_window_operators}
    Induct on the extra length beyond $L_0+1$. The induction hypothesis gives the
    long left-window ground condition after deleting the last site, while
    Definition~\ref{def:tail_restrict_parent} and the contiguous-interval
    restriction identities above identify every fixed-prefix suffix with one of
    the assumed length-$L_0+1$ intervals.
    Theorem~\ref{thm:ground_space_extend_right_block_injective} then
    grows back the rightmost site.
\end{proof}

\begin{lemma}[Open-chain iteration for a sequence of subspaces]%
    \label{lem:open_chain_iteration_subspace_sequence}
    \lean{MPSTensor.contiguous_mem_of_restriction_intersection_submodules}
    \leanok
    \uses{rem:contiguous_window_operators}
    Let \(S_m\subseteq(\mathbb C^d)^{\otimes m}\) be subspaces, and let
    \(0<L\le N\). Suppose that every non-wrapping length-\(L\) interval of
    \(\psi\), for every fixed choice of the complementary sites, lies in \(S_L\).
    Suppose also that, for every \(m\ge L\),
    \[
        \mathbb C^d\otimes S_m\cap S_m\otimes\mathbb C^d=S_{m+1}.
    \]
    Then \(\psi\in S_N\).
\end{lemma}

\begin{proof}
    \leanok
    \uses{rem:contiguous_window_operators}
    Induct on the length \(m\) of a non-wrapping interval.  The case \(m=L\)
    is the assumed local condition.  If all length-\(m\) intervals lie in
    \(S_m\), then the two one-site restrictions of a length-\(m+1\) interval are
    adjacent length-\(m\) intervals.  Hence the length-\(m+1\) interval lies in
    \(\mathbb C^d\otimes S_m\cap S_m\otimes\mathbb C^d\), and the displayed
    identity places it in \(S_{m+1}\).  Taking \(m=N\) and the interval starting
    at the first site gives \(\psi\in S_N\).
\end{proof}

\begin{lemma}[Periodic constraints and propagated subspaces]%
    \label{lem:periodic_constraints_propagated_subspace_sequence}
    \lean{MPSTensor.chainGroundSpace_le_of_local_le_restriction_intersection_submodules}
    \leanok
    \uses{lem:open_chain_iteration_subspace_sequence, rem:cyclic_window_operators}
    Let \(S_m\subseteq(\mathbb C^d)^{\otimes m}\) be subspaces, and let
    \(0<L\le N\).  Suppose that \(G_L(A)\subseteq S_L\).  Suppose also that, for
    every \(m\ge L\),
    \[
        \mathbb C^d\otimes S_m\cap S_m\otimes\mathbb C^d=S_{m+1}.
    \]
    Then the periodic local constraints imply
    \[
        \mathcal G_{N,L}(A)\subseteq S_N.
    \]
\end{lemma}

\begin{proof}
    \leanok
    \uses{lem:open_chain_iteration_subspace_sequence, rem:cyclic_window_operators}
    If \(\psi\in\mathcal G_{N,L}(A)\), then every length-\(L\) cyclic interval of
    \(\psi\) lies in \(G_L(A)\), and hence in \(S_L\).  A non-wrapping interval
    is the cyclic interval starting at the same site.  Therefore all
    non-wrapping length-\(L\) intervals lie in \(S_L\), and
    Lemma~\ref{lem:open_chain_iteration_subspace_sequence} gives
    \(\psi\in S_N\).
\end{proof}

\begin{lemma}[Cyclic-window range antitonicity]%
    \label{lem:cyclic_window_range_antitonicity}
    \lean{MPSTensor.chainGroundSpace_le_chainGroundSpace_succ,
        MPSTensor.chainGroundSpace_le_chainGroundSpace_of_le}
    \leanok
    \uses{rem:cyclic_window_operators}
    On a nonempty periodic chain, longer cyclic-window constraints imply all
    shorter cyclic-window constraints: if $L' \le L \le N$, then
    $\mathcal G_{N,L}(A) \subseteq \mathcal G_{N,L'}(A)$.
\end{lemma}

\begin{proof}
    \leanok
    \uses{lem:ground_space_in_left_ground}
    Fix a cyclic window of length $L+1$ and restrict its last site. The resulting
    length-$L$ vector is again in the ground space, so one-site peeling gives the
    inclusion for consecutive ranges. Iteration gives the general antitonicity.
\end{proof}

\begin{theorem}[Open-chain containment from cyclic local constraints]%
    \label{thm:cyclic_normal_open_chain_range_reduction}
    \lean{MPSTensor.chainGroundSpace_le_groundSpace_of_isNBlkInjective}
    \leanok
    \uses{lem:cyclic_window_range_antitonicity,
        thm:normal_open_chain_range_reduction}
    Let $A$ be $L_0$-block injective with $L_0>0$ and $D \ge 1$.
    If an $N$-site periodic-chain state satisfies all cyclic ground-space
    constraints at some range $L_0+1 \le L \le N$, then it lies in the
    open-chain ground space $G_N(A)$.
\end{theorem}

\begin{proof}
    \leanok
    \uses{lem:cyclic_window_range_antitonicity,
        thm:normal_open_chain_range_reduction}
    Lemma~\ref{lem:cyclic_window_range_antitonicity} reduces the cyclic
    hypotheses to range $L_0+1$. On every non-wrapping interval, the cyclic and
    contiguous restriction maps agree, so
    Theorem~\ref{thm:normal_open_chain_range_reduction} applies.
\end{proof}

\begin{lemma}[Boundary-crossing compatibility at the last site]%
    \label{lem:wrapped_window_one_sided_compatibility}
    \lean{MPSTensor.wrapping_window_compatibility_of_isNBlkInjective}
    \leanok
    \uses{def:l_blk_injective, rem:cyclic_window_operators}
    Assume $A$ is $L_0$-block injective with $L_0>0$ and $M \ge L_0$.
    If the reduced cyclic interval used in closing the boundary at the last site
    has boundary matrices $Y_\tau$, then for every physical letter $j$ one has
    \begin{equation}\label{eq:ph_wrapped_window_compat}
        C^+_\tau A^j X = Y_\tau A^j,
    \end{equation}
    where $C^+_\tau$ is the complementary word seen by that cyclic interval.
\end{lemma}

\begin{proof}
    \leanok
    \uses{lem:ground_space_map_injective_blk}
    The trace identity for this cyclic interval is tested against all words of
    length $L_0$. Injectivity of $\Gamma_{L_0}$, which follows from
    $L_0$-block injectivity, strips the test word and gives the displayed matrix
    identity~(\ref{eq:ph_wrapped_window_compat}).
\end{proof}

\begin{lemma}[Boundary-crossing compatibility at the opposite boundary position]%
    \label{lem:mirror_wrapped_window_one_sided_compatibility}
    \lean{MPSTensor.wrapping_window_mirror_compatibility_of_isNBlkInjective}
    \leanok
    \uses{def:l_blk_injective, rem:cyclic_window_operators}
    Assume $A$ is $L_0$-block injective with $L_0>0$ and $M \ge L_0$.
    If the opposite reduced cyclic interval used in closing the boundary has
    boundary matrices $Y_\tau$, then
    \[
        X A^j C^-_\tau = A^j Y_\tau
    \]
    for every physical letter $j$, where $C^-_\tau$ is the complementary word
    seen from the opposite cyclic position.
\end{lemma}

\begin{proof}
    \leanok
    \uses{lem:ground_space_map_injective_blk}
    Factor the cyclic word at the opposite cyclic position, rotate
    the trace, and use injectivity of $\Gamma_{L_0}$ to remove the
    length-$L_0$ test word.
\end{proof}

\begin{theorem}[The two boundary-crossing compatibilities]%
    \label{thm:reduced_wrapped_boundary_compatibilities}
    \lean{MPSTensor.chainGroundSpace_wrapped_boundary_compatibilities_of_isNBlkInjective}
    \leanok
    \uses{def:chain_ground_space, def:l_blk_injective,
        lem:cyclic_window_range_antitonicity,
        lem:wrapped_window_one_sided_compatibility,
        lem:mirror_wrapped_window_one_sided_compatibility}
    Let $A$ be $L_0$-block injective with $L_0>0$ and $D \ge 1$.
    Suppose $\psi \in \mathcal G_{N,L}(A)$, $N \ge 2$, $L_0 < L \le N$, and
    $\psi = \Gamma_N(X)$ for a boundary matrix $X$. Then there are two families
    of boundary matrices $Y^+_\tau$ and $Y^-_\tau$, indexed by fixed complement
    values $\tau$, such that for every physical letter $j$,
    \begin{equation}\label{eq:ph_reduced_wrapped_compat}
        C^+_\tau A^j X = Y^+_\tau A^j,
        \qquad
        X A^j C^-_\tau = A^j Y^-_\tau,
    \end{equation}
    where $C^+_\tau$ and $C^-_\tau$ are the complementary words exposed by the
    two opposite cyclic intervals used in closing the boundary, both of reduced
    length $L_0+1$.
\end{theorem}

\begin{proof}
    \leanok
    \uses{lem:cyclic_window_range_antitonicity,
        lem:wrapped_window_one_sided_compatibility,
        lem:mirror_wrapped_window_one_sided_compatibility}
    First use Lemma~\ref{lem:cyclic_window_range_antitonicity} to reduce the
    cyclic constraints from range $L$ to range $L_0+1$. Then apply the two
    cyclic-interval compatibility lemmas at the two boundary positions to obtain
    the displayed one-sided identities~(\ref{eq:ph_reduced_wrapped_compat}).
\end{proof}

\begin{lemma}[Boundary-crossing equations for $\Gamma_{M+1}(X)$]%
    \label{lem:closure_property_boundary_crossing_equations_ground_space_map}
    \lean{MPSTensor.closure_property_wrapped_mirror_compatibilities_of_groundSpaceMap}
    \leanok
    \uses{def:ground_space_map, rem:cyclic_window_operators,
        lem:wrapped_window_one_sided_compatibility,
        lem:mirror_wrapped_window_one_sided_compatibility}
    Let $A$ be $L_0$-block-injective with $L_0>0$, $D\ge1$, and $L_0\le M$.
    Suppose $\psi=\Gamma_{M+1}(X)$, and suppose matrices $Y_i(\tau)$ represent
    the length-$(L_0+1)$ cyclic restrictions
    \[
        \operatorname{Res}^{\tau}_{i,L_0+1}(\psi)
        =
        \Gamma_{L_0+1}(Y_i(\tau)).
    \]
    Then, for every physical letter $j$ and boundary condition $\tau$,
    \begin{equation}\label{eq:ph_closure_boundary_crossing_gamma_equations}
        A^{\tau_{L_0}}\cdots A^{\tau_{M-1}} A^j X
        =
        Y_M(\tau)A^j,
        \qquad
        X A^j A^{\tau_1}\cdots A^{\tau_{M-L_0}}
        =
        A^jY_{M+1-L_0}(\tau).
    \end{equation}
\end{lemma}

\begin{proof}
    \leanok
    Substitute $\psi=\Gamma_{M+1}(X)$ in the two cyclic restrictions at the
    last site and at the opposite boundary-crossing support. Lemmas
    \ref{lem:wrapped_window_one_sided_compatibility} and
    \ref{lem:mirror_wrapped_window_one_sided_compatibility} then give the two
    equations in~(\ref{eq:ph_closure_boundary_crossing_gamma_equations}).
\end{proof}

\begin{lemma}[One-sided products for a common complement word]
    \label{lem:closure_property_boundary_one_sided_products_ground_space_map}
    \lean{MPSTensor.closure_property_boundary_one_sided_products_of_groundSpaceMap}
    \leanok
    \uses{lem:closure_property_boundary_crossing_equations_ground_space_map,
        lem:wrapped_middle_backgrounds}
    Let $A$ be $L_0$-block-injective with $L_0>0$, $D\ge1$, and $L_0\le M$.
    Suppose $\psi=\Gamma_{M+1}(X)$, and suppose matrices $Y_i(\tau)$ represent
    the length-$(L_0+1)$ cyclic restrictions
    \[
        \operatorname{Res}^{\tau}_{i,L_0+1}(\psi)
        =
        \Gamma_{L_0+1}(Y_i(\tau)).
    \]
    For every boundary letter $\eta$, every physical letter $j$, and every
    complementary word $\mu$,
    \[
        Y_M(\tau^+_\eta(\mu))A^j
        =
        A^\mu A^j X,
        \qquad
        X A^j A^\mu
        =
        A^jY_{M+1-L_0}(\tau^-_\eta(\mu)).
    \]
\end{lemma}

\begin{proof}
    \leanok
    Substitute $\tau^+_\eta(\mu)$ and $\tau^-_\eta(\mu)$ into
    Lemma~\ref{lem:closure_property_boundary_crossing_equations_ground_space_map}.
    The exposed-complement identities of
    Lemma~\ref{lem:wrapped_middle_backgrounds} identify both complementary
    products with $A^\mu$.
\end{proof}

\begin{lemma}[Common complement word for the two boundary supports]%
    \label{lem:wrapped_middle_backgrounds}
    \lean{MPSTensor.wrappedMiddleBackground,
        MPSTensor.mirrorMiddleBackground,
        MPSTensor.wrappedMiddleBackground_complement,
        MPSTensor.mirrorMiddleBackground_complement}
    \leanok
    \uses{rem:cyclic_window_operators}
    On a chain of length $N=M+1$, with $L_0\le M$, fix a physical letter $\eta$
    used outside the complementary sites and a word $\mu$ of length
    $M+1-(L_0+1)$. Write $\tau^+_\eta(\mu)$ for the boundary condition at the
    support crossing the last site and $\tau^-_\eta(\mu)$ for the boundary
    condition at the opposite boundary-crossing support. Their exposed complements
    are both exactly $\mu$.
\end{lemma}

\begin{proof}
    \leanok
    Put the same letter $\eta$ on the remaining sites. Fill the physical sites
    $L_0,\ldots,M-1$ with $\mu$ for the support crossing the last site, and fill
    the sites $1,\ldots,M-L_0$ with $\mu$ for the opposite boundary-crossing
    support.
    The two complement-extraction identities are then immediate from the index
    arithmetic. The same outside-site calculation gives, for any $\rho$ with
    $\rho_{k+L_0}=\mu_k$,
    \[
        \operatorname{Res}^{\rho}_{M,L_0+1}\psi
        =
        \operatorname{Res}^{\tau^+_\eta(\mu)}_{M,L_0+1}\psi,
    \]
    and, for any $\rho$ with $\rho_{k+1}=\mu_k$,
    \[
        \operatorname{Res}^{\rho}_{M+1-L_0,L_0+1}\psi
        =
        \operatorname{Res}^{\tau^-_\eta(\mu)}_{M+1-L_0,L_0+1}\psi.
    \]
\end{proof}

\begin{lemma}[Boundary equations from a common complement word]
    \label{lem:boundary_assignment_first_letter_products}
    \lean{MPSTensor.cyclicRestrictₗ_first_products_eq_of_restriction_eq,
        MPSTensor.cyclicRestrictₗ_wrappedMiddleBackground_eq_of_complement_eq,
        MPSTensor.cyclicRestrictₗ_mirrorMiddleBackground_eq_of_complement_eq,
        MPSTensor.wrappedMiddleBackground_first_products_eq_of_complement_eq,
        MPSTensor.mirrorMiddleBackground_first_products_eq_of_complement_eq,
        MPSTensor.wrappedMiddleBackground_first_products_eq_of_complement_eq_right_word,
        MPSTensor.mirrorMiddleBackground_first_products_eq_of_complement_eq_right_word}
    \leanok
    \uses{rem:cyclic_window_operators, lem:wrapped_middle_backgrounds}
    Let $A$ be $L_0$-block-injective, with $L_0>0$, and let $L_0\le M$.
    Suppose two length-$(L_0+1)$ restrictions of a state $\psi$ at the same
    cyclic support are represented by
    \[
        \operatorname{Res}^{\rho}_{i,L_0+1}\psi=\Gamma_{L_0+1}(Y_\rho),
        \qquad
        \operatorname{Res}^{\tau}_{i,L_0+1}\psi=\Gamma_{L_0+1}(Y_\tau).
    \]
    If these two restrictions are equal, then for every physical letter $j$,
    \[
        Y_\rho A^j = Y_\tau A^j .
    \]
    In particular, the last-site boundary-crossing condition satisfies
    \[
        \rho_{k+L_0}=\mu_k
        \Longrightarrow
        Y_\rho A^j = Y_{\tau^+_\eta(\mu)}A^j,
    \]
    and the opposite boundary-crossing condition satisfies
    \[
        \rho_{k+1}=\mu_k
        \Longrightarrow
        Y_\rho A^j = Y_{\tau^-_\eta(\mu)}A^j .
    \]
    Consequently, for every word $\sigma$ of length $L_0$, the two displayed
    implications remain true after right multiplication by $A^\sigma$.
\end{lemma}

\begin{proof}
    \leanok
    Apply the first-letter restriction to both equal length-$(L_0+1)$
    restrictions. The two resulting length-$L_0$ vectors are represented by
    $\Gamma_{L_0}(Y_\rho A^j)$ and $\Gamma_{L_0}(Y_\tau A^j)$, and block
    injectivity makes $\Gamma_{L_0}$ injective. The two displayed special cases
    follow from Lemma~\ref{lem:wrapped_middle_backgrounds}. The
    length-$L_0$ word forms are obtained by multiplying these equations on the
    right by $A^\sigma$.
\end{proof}

\begin{lemma}[Reduction to compatible boundary conditions]
    \label{lem:boundary_closing_product_compatible_backgrounds}
    \lean{MPSTensor.boundary_closing_product_eq_of_compatible_backgrounds}
    \leanok
    \uses{lem:boundary_assignment_first_letter_products}
    Let $A$ be $L_0$-block-injective, with $L_0>0$, and let $L_0\le M$.
    Let $\psi$ be a state on the chain of length $M+1$, and let the matrices
    $Y_i(\tau)$ represent its length-$(L_0+1)$ cyclic restrictions:
    \[
        \operatorname{Res}^{\tau}_{i,L_0+1}(\psi)
        =
        \Gamma_{L_0+1}(Y_i(\tau)).
    \]
    Fix a boundary letter $\eta$ and a complementary word $\mu$. Let $\rho^+$
    and $\rho^-$ be two boundary conditions satisfying
    \[
        \rho^+_{k+L_0}=\mu_k,
        \qquad
        \rho^-_{k+1}=\mu_k .
    \]
    If the adjacent-window comparison gives, for every physical letter $j$ and
    every word $\sigma$ of length $L_0$,
    \[
        Y_M(\rho^+)A^jA^\sigma
        =
        Y_{M+1-L_0}(\rho^-)A^jA^\sigma ,
    \]
    then the same product equation holds for $\tau^+_\eta(\mu)$ and
    $\tau^-_\eta(\mu)$:
    \[
        Y_M(\tau^+_\eta(\mu))A^jA^\sigma
        =
        Y_{M+1-L_0}(\tau^-_\eta(\mu))A^jA^\sigma .
    \]
\end{lemma}

\begin{proof}
    \leanok
    By Lemma~\ref{lem:boundary_assignment_first_letter_products},
    \[
        Y_M(\rho^+)A^j = Y_M(\tau^+_\eta(\mu))A^j,
        \qquad
        Y_{M+1-L_0}(\rho^-)A^j
        =
        Y_{M+1-L_0}(\tau^-_\eta(\mu))A^j.
    \]
    Multiplying on the right by $A^\sigma$ and applying the hypothesis gives the
    displayed product equation.
\end{proof}

\begin{theorem}[Pointwise reduction to compatible boundary conditions]
    \label{thm:boundary_closing_product_pointwise_compatible_boundary_assignments}
    \lean{MPSTensor.boundary_closing_product_eq_of_pointwise_compatible_boundary_assignments}
    \leanok
    \uses{lem:boundary_assignment_first_letter_products}
    Let $A$ be $L_0$-block-injective, with $L_0>0$, and let $L_0\le M$.
    Let $\psi$ be a state on the chain of length $M+1$, and let the matrices
    $Y_i(\tau)$ represent its length-$(L_0+1)$ cyclic restrictions:
    \[
        \operatorname{Res}^{\tau}_{i,L_0+1}(\psi)
        =
        \Gamma_{L_0+1}(Y_i(\tau)).
    \]
    Fix a boundary letter $\eta$ and a complementary word $\mu$. Suppose that,
    for each physical letter $j$ and word $\sigma$ of length $L_0$, boundary
    conditions $\rho^+_{j,\sigma}$ and $\rho^-_{j,\sigma}$ satisfy, for every
    complementary position $k$,
    \[
        \rho^+_{j,\sigma}(k+L_0)=\mu_k,
        \qquad
        \rho^-_{j,\sigma}(k+1)=\mu_k .
    \]
    If, for every $j$ and $\sigma$,
    \[
        Y_M(\rho^+_{j,\sigma})A^jA^\sigma
        =
        Y_{M+1-L_0}(\rho^-_{j,\sigma})A^jA^\sigma ,
    \]
    then, for every $j$ and $\sigma$,
    \[
        Y_M(\tau^+_\eta(\mu))A^jA^\sigma
        =
        Y_{M+1-L_0}(\tau^-_\eta(\mu))A^jA^\sigma .
    \]
\end{theorem}

\begin{proof}
    \leanok
    Fix $j$ and $\sigma$. Lemma~\ref{lem:boundary_assignment_first_letter_products}
    gives
    \[
        Y_M(\rho^+_{j,\sigma})A^j
        =
        Y_M(\tau^+_\eta(\mu))A^j,
        \qquad
        Y_{M+1-L_0}(\rho^-_{j,\sigma})A^j
        =
        Y_{M+1-L_0}(\tau^-_\eta(\mu))A^j.
    \]
    Multiply the two identities on the right by $A^\sigma$ and use the assumed
    product equation for the chosen pair $j,\sigma$.
\end{proof}

\begin{lemma}[Endpoint words at the closing boundary]
    \label{lem:boundary_closing_endpoint_word_products}
    \lean{MPSTensor.boundary_closing_endpoint_word_products_common_background}
    \leanok
    \uses{lem:cyclic_window_boundary_matrix_transport}
    Let \(A\) be \(L_0\)-block-injective with \(L_0>0\), and let \(L_0\le M\).
    Let \(\psi\) be a state on the chain of length \(M+1\). Suppose that
    matrices \(Y_r(\rho)\), for \(0\le r\le L_0-1\), represent the
    length-\((L_0+1)\) cyclic restrictions beginning at the sites
    \(M+1-L_0+r\):
    \[
        \operatorname{Res}^{\rho}_{M+1-L_0+r,L_0+1}\psi
        =
        \Gamma_{L_0+1}(Y_r(\rho)).
    \]
    Then
    \[
        Y_0(\rho)
        A^{\rho_{M+1-L_0}}\cdots A^{\rho_{M-1}}
        =
        A^{\rho_1}\cdots A^{\rho_{L_0-1}}
        Y_{L_0-1}(\rho).
    \]
    For \(L_0=1\), both products are empty.
\end{lemma}

\begin{proof}
    \leanok
    Use Lemma~\ref{lem:cyclic_window_boundary_matrix_transport} for the
    \(L_0-1\) adjacent restrictions starting at \(i_0=M+1-L_0\). For
    \(0\le r<L_0-1\), the two endpoint words are determined by
    \[
        i_0+r=M+1-L_0+r,\qquad
        i_0+r+L_0+1\equiv r+1 \pmod {M+1},
    \]
    Substituting these two words in the iterated transport identity gives the
    displayed product equation.
\end{proof}

\begin{lemma}[Boundary-condition product at the closing boundary]
    \label{lem:closure_property_boundary_condition_product_window_witnesses}
    \lean{MPSTensor.closure_property_boundary_condition_product_of_window_witnesses,
        MPSTensor.closure_property_boundary_condition_product_of_window_witnesses_mul_right}
    \leanok
    \uses{rem:cyclic_window_operators, lem:cyclic_window_boundary_matrix_transport}
    Let $A$ be $L_0$-block-injective with $L_0>0$, and let $L_0\le M$. Let
    $\psi$ be a state on the chain of length $M+1$, and suppose matrices
    $Y_i(\rho)$ represent the length-$(L_0+1)$ cyclic restrictions
    \[
        \operatorname{Res}^{\rho}_{i,L_0+1}(\psi)
        =
        \Gamma_{L_0+1}(Y_i(\rho)).
    \]
    Then, for every boundary condition $\rho$,
    \[
        Y_{M+1-L_0}(\rho)
        A^{\rho_{M+1-L_0}}\cdots A^{\rho_{M-1}}
        =
        A^{\rho_1}\cdots A^{\rho_{L_0-1}}Y_M(\rho).
    \]
    For $L_0=1$, both products are empty.
    The same equation remains true after multiplying both sides on the right by
    any matrix \(R\).
\end{lemma}

\begin{proof}
    \leanok
    Apply Lemma~\ref{lem:cyclic_window_boundary_matrix_transport} to the
    $L_0-1$ restrictions beginning at $M+1-L_0+r$, with
    \[
        Y_r(\rho)=Y_{M+1-L_0+r}(\rho).
    \]
    With $i_0=M+1-L_0$, the site labels satisfy
    \[
        i_0+r=M+1-L_0+r,
        \qquad
        i_0+r+L_0+1\equiv r+1 \pmod {M+1}
    \]
    hence, for $r=0,\ldots,L_0-2$,
    \[
        Y_{M+1-L_0+r}(\rho) A^{\rho_{M+1-L_0+r}}
        =
        A^{\rho_{r+1}}Y_{M+2-L_0+r}(\rho).
    \]
    Multiplying these $L_0-1$ equations from left to right gives the displayed
    product equation. The right-multiplied form follows by multiplying this
    equality by \(R\).
\end{proof}

\begin{lemma}[Transport followed by the last-boundary equation]
    \label{lem:boundary_transport_wrapped_product_ground_space_map}
    \lean{MPSTensor.closure_property_boundary_condition_transport_wrapped_product_of_groundSpaceMap}
    \leanok
    \uses{lem:closure_property_boundary_condition_product_window_witnesses,
        lem:closure_property_boundary_crossing_equations_ground_space_map}
    Let $A$ be $L_0$-block-injective with $L_0>0$ and $D\ge1$, and let
    $L_0\le M$.  Let $\psi=\Gamma_{M+1}(X)$, and suppose that matrices
    $Y_i(\rho)$ represent all length-$(L_0+1)$ cyclic restrictions of $\psi$.
    Then, for every boundary condition $\rho$ and every physical letter $j$,
    \[
        Y_{M+1-L_0}(\rho)
        A^{\rho_{M+1-L_0}}\cdots A^{\rho_{M-1}}A^j
        =
        A^{\rho_1}\cdots A^{\rho_{L_0-1}}
        A^{\rho_{L_0}}\cdots A^{\rho_{M-1}}A^jX .
    \]
    For $\rho=\tau^-_\eta(\mu)$ this is the transport identity supplied by the
    adjacent-window argument.  It is not the remaining padded identity, where
    the factor after $Y_{M+1-L_0}(\tau^-_\eta(\mu))$ is \(A^jA^\sigma\).
\end{lemma}

\begin{proof}
    \leanok
    Lemma~\ref{lem:closure_property_boundary_condition_product_window_witnesses}
    gives
    \[
        Y_{M+1-L_0}(\rho)
        A^{\rho_{M+1-L_0}}\cdots A^{\rho_{M-1}}A^j
        =
        A^{\rho_1}\cdots A^{\rho_{L_0-1}}Y_M(\rho)A^j .
    \]
    The last-boundary one-sided equation for $\psi=\Gamma_{M+1}(X)$ gives
    \[
        Y_M(\rho)A^j
        =
        A^{\rho_{L_0}}\cdots A^{\rho_{M-1}}A^jX .
    \]
    Substitution gives the displayed formula.
\end{proof}

\begin{lemma}[Long boundary-condition product at the closing boundary]
    \label{lem:closure_property_boundary_condition_long_product_window_witnesses}
    \lean{MPSTensor.closure_property_boundary_condition_long_product_of_window_witnesses}
    \leanok
    \uses{rem:cyclic_window_operators, lem:cyclic_window_boundary_matrix_transport}
    Let $A$ be $L_0$-block-injective with $L_0>0$, and let $L_0\le M$. Let
    $\psi$ be a state on the chain of length $M+1$, and suppose matrices
    $Y_i(\rho)$ represent the length-$(L_0+1)$ cyclic restrictions
    \[
        \operatorname{Res}^{\rho}_{i,L_0+1}(\psi)
        =
        \Gamma_{L_0+1}(Y_i(\rho)).
    \]
    Then, for every boundary condition $\rho$,
    \[
        Y_M(\rho)A^{\rho_M}A^{\rho_0}\cdots A^{\rho_{M-L_0}}
        =
        A^{\rho_{L_0}}\cdots A^{\rho_M}A^{\rho_0}
        Y_{M+1-L_0}(\rho).
    \]
    Each side has $M+2-L_0$ one-site factors, with site labels read modulo
    $M+1$.
\end{lemma}

\begin{proof}
    \leanok
    Apply Lemma~\ref{lem:cyclic_window_boundary_matrix_transport} to the
    $M+2-L_0$ restrictions beginning at $M+r$ modulo $M+1$. With $i_0=M$, the
    site labels satisfy
    \[
        i_0+r\equiv M+r,
        \qquad
        i_0+r+L_0+1\equiv L_0+r
        \pmod {M+1}.
    \]
    The final window is
    \[
        i_0+M+2-L_0\equiv M+1-L_0 \pmod {M+1}.
    \]
    Substituting these site labels in the iterated product identity gives the
    displayed equation.
\end{proof}

\begin{lemma}[One matrix $Y_\mu$ from compared boundary matrices]%
    \label{lem:two_sided_middle_of_wrapped_witness_comparison}
    \lean{MPSTensor.two_sided_middle_compatibility_of_wrapped_witness_comparison}
    \leanok
    \uses{lem:wrapped_middle_backgrounds}
    Fix a physical letter $\eta$ used outside the complementary sites. Suppose
    the two one-sided boundary matrix families $Y^+$ and $Y^-$ have been
    extracted from the cyclic windows. If, for every word on the complementary sites
    $\mu$, these boundary matrices agree on the boundary conditions built from
    $\eta$,
    \[
        Y^+_{\tau^+_\eta(\mu)} = Y^-_{\tau^-_\eta(\mu)},
    \]
    then there is a single family $Y_\mu$ satisfying
    \[
        A^\mu A^b X = Y_\mu A^b,
        \qquad
        X A^a A^\mu = A^a Y_\mu
    \]
    for all letters $a,b$.
\end{lemma}

\begin{proof}
    \leanok
    Define $Y_\mu$ to be the first boundary matrix evaluated on the boundary
    condition
    $\tau^+_\eta(\mu)$. Lemma~\ref{lem:wrapped_middle_backgrounds} rewrites the
    first cyclic-window complement as $\mu$, while the assumed comparison at
    $\tau^-_\eta(\mu)$ rewrites the opposite boundary matrix to the same
    $Y_\mu$.
\end{proof}

\begin{lemma}[Commutation from two one-sided equations]%
    \label{lem:common_middle_same_witness_commutation}
    \lean{MPSTensor.commutes_words_of_two_sided_middle_compatibility}
    \leanok
    Suppose a single family of boundary matrices $Y_\mu$, indexed by words
    $\mu$ on the complementary sites, satisfies the two one-sided identities
    \begin{equation}\label{eq:ph_common_middle_one_sided}
        A^\mu A^b X = Y_\mu A^b,
        \qquad
        X A^a A^\mu = A^a Y_\mu
    \end{equation}
    for every pair of physical letters $a,b$. Then $X$ commutes with every word
    product $A^a A^\mu A^b$ obtained by adjoining one letter on each side of the
    word $A^\mu$.
\end{lemma}

\begin{proof}
    \leanok
    Multiply the second identity in~(\ref{eq:ph_common_middle_one_sided}) on the
    right by $A^b$ and the first identity on the left by $A^a$:
    \[
        X A^a A^\mu A^b = A^a Y_\mu A^b
        = A^a A^\mu A^b X.
    \]
\end{proof}

\begin{theorem}[Boundary matrix commutation from periodic closure]%
    \label{thm:boundary_matrix_commutes}
    \lean{MPSTensor.boundary_matrix_commutes}
    \leanok
    \uses{def:ground_space_map, def:ground_space_parent, thm:ground_space_map_injective}
    Let $A$ be injective with $D \ge 1$, and let $\psi = \Gamma_{N}(X)$ for some boundary matrix
    $X \in \MN{D}$. If every cyclic restriction of length $L$ of $\psi$ lies in $G_{L}(A)$,
    with $N \ge 2$ and $1 < L \le N$, then, for every $j = 0,\ldots,d-1$,
    \[
        X A^{j} = A^{j} X
    \]
\end{theorem}

\begin{proof}
    \leanok
    The periodic boundary condition produces matrix identities of the form
    $XMW=MY$ for arbitrary test matrices $M$ and complementary words $W$.
    Spanning first over the length-$(L-1)$ tails and then over the complementary words,
    injectivity forces $(XM-MX)W = 0$ for every $W$, hence $XM = MX$ for every matrix
    $M$. In particular, $X$ commutes with each letter $A^{j}$.
\end{proof}

\begin{definition}[MPV submodule]\label{def:mpv_submodule}
    \lean{MPSTensor.mpvSubmodule}
    \leanok
    \uses{def:mpv}
    The subspace spanned by the MPS vector
    $V^{(N)}(A)(\sigma) = \tr(A^{\sigma_0}\cdots A^{\sigma_{N-1}})$.
\end{definition}

\begin{lemma}[The MPV is the ground-space map at the identity]
    \label{lem:mpv_eq_ground_space_map_one}
    \lean{MPSTensor.mpv_eq_groundSpaceMap_one}
    \leanok
    \uses{def:ground_space_map, def:mpv}
    For every chain length $N$,
    \[
        V^{(N)}(A) = \Gamma_{N}(\Id).
    \]
\end{lemma}

\begin{lemma}[MPV lies in the ground space]\label{lem:mpv_mem_ground_space}
    \lean{MPSTensor.mpv_mem_groundSpace}
    \leanok
    \uses{def:ground_space_parent, def:mpv}
    $V^{(N)}(A) \in G_N(A)$ for every $N$, with boundary matrix $X = I$.
\end{lemma}

\begin{proof}
    \leanok
    $V^{(N)}(A)(\sigma) = \tr(A^\sigma) = \tr(A^\sigma \cdot I) = \Gamma_N(I)(\sigma)$.
\end{proof}

\begin{definition}[Unique ground state]\label{def:has_unique_ground_state}
    \lean{MPSTensor.HasUniqueGroundState}
    \leanok
    A submodule $S$ has a unique ground state if $\dim S = 1$.
\end{definition}

\begin{theorem}[One-dimensionality criterion]\label{thm:has_unique_ground_state_iff_proportional}
    \lean{MPSTensor.hasUniqueGroundState_iff_proportional}
    \leanok
    \uses{def:has_unique_ground_state}
    For finite-dimensional $S$, $\dim S = 1$ iff there exists a nonzero
    $\psi_0 \in S$ such that every $\psi \in S$ is of the form $c\,\psi_0$.
\end{theorem}

\begin{proof}
    \leanok
    Apply the standard finite-dimensional criterion that a nonzero space has
    dimension one exactly when all of its vectors are proportional to a single
    nonzero vector.
\end{proof}

\begin{definition}[Periodic chain ground space]\label{def:chain_ground_space}
    \lean{MPSTensor.chainGroundSpace}
    \leanok
    \uses{def:ground_space_parent}
    For $N > 0$ and $L \le N$, define the \emph{periodic chain ground space}
    \[
        \mathcal{G}_{N,L}(A) := \bigl\{\psi \in (\{0,\ldots,d{-}1\}^N \to \C)
        : \forall i \in \{0,\ldots,N{-}1\},\
        \psi|_{[i,i+L-1]} \in G_L(A)\bigr\},
    \]
    where the condition $\psi|_{[i,i+L-1]} \in G_L(A)$ means that for every
    choice of physical indices outside the window, the restriction of $\psi$
    to that window lies in $G_L(A)$.
    When $N = 0$ or $L > N$, set $\mathcal{G}_{N,L}(A) := \top$ by convention.
\end{definition}

\begin{lemma}[Cyclic-window witnesses inside the chain ground space]
    \label{lem:chain_ground_space_window_witnesses}
    \lean{MPSTensor.chainGroundSpace_window_witnesses}
    \leanok
    \uses{def:chain_ground_space, def:ground_space_parent}
    If $0<N$, $L\le N$, and $\psi\in\mathcal G_{N,L}(A)$, then for every
    cyclic window and every outside configuration there is a boundary matrix $Y$
    such that the corresponding restricted $L$-site state is $\Gamma_L(Y)$.
\end{lemma}

\begin{proof}
    \leanok
    This is the defining cyclic-window condition for membership in
    $\mathcal G_{N,L}(A)$, together with the definition of the local ground
    space as the range of $\Gamma_L$.
\end{proof}

\begin{lemma}[The MPV satisfies every cyclic window constraint]
    \label{lem:mpv_mem_chain_ground_space}
    \lean{MPSTensor.mpv_mem_chainGroundSpace}
    \leanok
    \uses{def:chain_ground_space, lem:mpv_window_mem_ground_space}
    If $0 < N$ and $L \le N$, then
    \[
        V^{(N)}(A) \in \mathcal G_{N,L}(A).
    \]
\end{lemma}

\begin{theorem}[MPV nonvanishing for block-injective tensors]%
    \label{thm:mpv_ne_zero_blk_injective}
    \lean{MPSTensor.mpv_ne_zero_of_isNBlkInjective}
    \leanok
    \uses{def:mpv, def:l_blk_injective}
    If $A$ is $L_0$-block-injective with $L_0 > 0$, $D \ge 1$, and
    $N \ge L_0 + 1$, then $V^{(N)}(A) \ne 0$ in the $N$-site space.
\end{theorem}

\begin{proof}
    \leanok
    Suppose for contradiction that $V^{(N)}(A) = 0$, i.e.\
    $\tr(A^\sigma) = 0$ for every word $\sigma$ of length~$N$.
    For any words $w$ and $u$ with $|w| = N - L_0$ and $|u| = L_0$,
    we have $\tr(A^w A^u) = 0$.
    Since $A$ is $L_0$-block-injective, $\{A^u : |u| = L_0\}$
    spans $\MN{D}$. By nondegeneracy of the trace pairing,
    $A^w = 0$ for every $|w| = N - L_0$.
    Repeating the argument with words of length $N - 2L_0$ (since
    $|w'| + L_0 = N - L_0$, the product $A^{w'} A^{u_1}$ has length
    $N - L_0$, so $A^{w'} A^{u_1} = 0$ by the previous step) and iterating,
    we eventually reach $\tr(A^u) = 0$ for $|u| = L_0$, which gives
    $I = 0$, a contradiction.
\end{proof}

\subsection{Block-word stripping}

\begin{theorem}[Common right factor from block-word compatibility]
    \label{thm:block_word_strip_right_factor}
    \lean{MPSTensor.exists_right_factor_of_block_word_compatibility}
    \leanok
    \uses{def:l_blk_injective}
    Assume $A$ is $L_0$-block-injective with $L_0 > 0$. Let $\{Z_u\}_{|u|=L_0}$
    be a family indexed by words of length $L_0$. If for some $K \ge 0$ and every
    suffix word $w$ of length $K$ there exists $Y_w \in \MN{D}$ such that,
    for every word $u$ of length~$L_0$,
    \[
        Z_u A^w = A^u Y_w
    \]
    then there exists $X \in \MN{D}$ with $Z_u = A^u X$ for every word $u$ of
    length $L_0$.
\end{theorem}

\begin{proof}
    \leanok
    Induct on $K$. For $K = 0$ the empty suffix gives $Z_u = A^u Y_{\emptyset}$.
    For the successor step, fix the first physical index and apply the induction
    hypothesis to the shortened suffix. This reduces the statement to letter
    compatibility on the right. Extending that letter compatibility to all
    length-$L_0$ words and decomposing $\Id$ in the spanning family
    $\{A^u : |u| = L_0\}$ yields a common factor $X$.
\end{proof}

\begin{theorem}[Block-word commutation from long-word commutation]
    \label{thm:block_word_strip_commutes}
    \lean{MPSTensor.commutes_block_words_of_commutes_long_words_of_isNBlkInjective}
    \leanok
    \uses{thm:block_word_strip_right_factor, def:l_blk_injective}
    If $A$ is $L_0$-block-injective with $L_0 > 0$, $m \ge L_0$, and
    $X \in \MN{D}$ commutes with every length-$m$ product $A^\omega$, then $X$
    already commutes with every length-$L_0$ product $A^u$.
\end{theorem}

\begin{proof}
    \leanok
    Write each length-$m$ word as a concatenation $uw$ with $|u| = L_0$ and
    $|w| = m - L_0$. The hypothesis gives
    \[
        X A^u A^w = A^u A^w X = A^u (A^w X).
    \]
    Thus the family $Z_u := X A^u$ satisfies the hypothesis of
    Theorem~\ref{thm:block_word_strip_right_factor}. One obtains
    $X A^u = A^u R$ for all $|u| = L_0$. Decomposing $\Id$ in the span of the
    length-$L_0$ words shows $R = X$, hence $X A^u = A^u X$.
\end{proof}

\begin{theorem}[Centrality from long-word commutation]
    \label{thm:block_word_strip_central}
    \lean{MPSTensor.commutes_all_of_commutes_long_words_of_isNBlkInjective}
    \leanok
    \uses{thm:block_word_strip_commutes, def:l_blk_injective}
    If $A$ is $L_0$-block-injective with $L_0 > 0$, $m \ge L_0$, and
    $X \in \MN{D}$ commutes with every length-$m$ word product, then $X$
    commutes with every matrix in $\MN{D}$.
\end{theorem}

\begin{proof}
    \leanok
    By Theorem~\ref{thm:block_word_strip_commutes}, $X$ commutes with every
    length-$L_0$ word product. These products span $\MN{D}$ by block
    injectivity, so linearity gives $XM = MX$ for every $M \in \MN{D}$.
\end{proof}

\begin{lemma}[Amplifying fixed-length word commutation]
    \label{lem:word_commutation_multiple_lengths}
    \lean{MPSTensor.commutes_words_mul_of_commutes_words}
    \leanok
    If a boundary matrix $X$ commutes with every word product of length $m$, then
    it commutes with every word product whose length is a multiple $q\,m$.
\end{lemma}

\begin{proof}
    \leanok
    Proceed by induction on $q$. For $q=0$, the claim is trivial. For the
    inductive step, let a word of length $(q+1)m$ be given and write it as a
    length-$m$ prefix followed by a length-$q\,m$ suffix. The prefix commutes with
    $X$ by hypothesis, and the suffix commutes by the induction hypothesis.
\end{proof}

\begin{lemma}[Padded complement annihilation]
    \label{lem:padded_complement_annihilation}
    \lean{MPSTensor.eq_zero_of_mul_evalWord_eq_zero_of_isNBlkInjective_of_le_mul,
        MPSTensor.eq_zero_of_evalWord_mul_eq_zero_of_isNBlkInjective_of_le_mul}
    \leanok
    \uses{def:l_blk_injective}
    If $A$ is $L_0$-block-injective, $q \ge 1$, and \(k \le qL_0\), then the
    two one-sided annihilation statements are:
    \[
        Z A^w = 0 \quad (|w|=k) \qquad\Longrightarrow\qquad Z=0.
    \]
    The left-handed form is
    \[
        A^w Z = 0 \quad (|w|=k) \qquad\Longrightarrow\qquad Z=0.
    \]
\end{lemma}

\begin{proof}
    \leanok
    The hypothesis $\spn\{A^u: |u|=L_0\}=\MN{D}$ implies the same spanning
    statement at every positive multiple $qL_0$. Each length-$qL_0$ word factors
    as a length-$k$ prefix followed by padding, so $Z$ annihilates all generators
    of the full span and hence annihilates the identity.
    For the left-handed form, factor each length-$qL_0$ word as a prefix followed
    by a length-$k$ suffix and use the same spanning argument.
\end{proof}

\begin{lemma}[One-sided boundary witnesses are unique]
    \label{lem:one_sided_boundary_witness_unique}
    \lean{MPSTensor.right_witness_unique_of_isNBlkInjective}
    \lean{MPSTensor.left_witness_unique_of_isNBlkInjective}
    \leanok
    \uses{def:l_blk_injective, lem:padded_complement_annihilation,
        lem:wordspan_blk_inj}
    Let $A$ be $L_0$-block-injective with $L_0>0$.  If two boundary matrices
    have the same product with every one-site tensor on one fixed side, then the
    two boundary matrices are equal.
\end{lemma}

\begin{proof}
    \leanok
    If $Y_1 A^j = Y_2 A^j$ for every physical index $j$, then for every
    nonempty word $w$, induction on the word gives
    \[
        Y_1 A^w = Y_2 A^w
    \]
    Since $A$ is $L_0$-block-injective,
    \[
        \spn\{A^w : |w|=L_0\}=\MN{D}.
    \]
    Thus the left multiplication maps $M\mapsto Y_1 M$ and
    $M\mapsto Y_2 M$ are equal on $\MN{D}$, and evaluating on the identity
    gives $Y_1=Y_2$.  The other one-sided statement is the same argument, with
    right multiplication in place of left multiplication.
\end{proof}

\begin{lemma}[Generator commutation from long-word commutation]
    \label{lem:boundary_matrix_commutes_long_words}
    \lean{MPSTensor.boundary_matrix_commutes_of_isNBlkInjective_of_long_word_commutes}
    \leanok
    \uses{thm:block_word_strip_central}
    Under the hypotheses of Theorem~\ref{thm:block_word_strip_central}, one has
    $X A^i = A^i X$ for every physical index $i$.
\end{lemma}

\begin{proof}
    \leanok
    Apply Theorem~\ref{thm:block_word_strip_central} with $M = A^i$.
\end{proof}

\begin{theorem}[MPS-line containment from long-word commutation]
    \label{thm:ground_space_map_mpv_of_long_word_commutes}
    \lean{MPSTensor.groundSpaceMap_mem_mpvSubmodule_of_isNBlkInjective_of_long_word_commutes}
    \leanok
    \uses{def:mpv_submodule, thm:block_word_strip_central}
    Let $A$ be injective after blocking $L_0>0$ sites, and let $m \ge L_0$.
    If a boundary matrix $X$ commutes with every length-$m$ word product, then
    \[
        \Gamma_N(X)\in\spn\{V^{(N)}(A)\}.
    \]
\end{theorem}

\begin{proof}
    \leanok
    Theorem~\ref{thm:block_word_strip_central} makes $X$ commute with the whole
    matrix algebra. The center of $\MN{D}$ consists of scalar matrices, so
    $X=c\Id$, and therefore $\Gamma_N(X)=cV^{(N)}(A)$.
\end{proof}

\begin{theorem}[MPS-line containment from positive-length commutation]
    \label{thm:ground_space_map_mpv_of_positive_word_commutes}
    \lean{MPSTensor.groundSpaceMap_mem_mpvSubmodule_of_isNBlkInjective_of_positive_word_commutes}
    \leanok
    \uses{lem:word_commutation_multiple_lengths,
        thm:ground_space_map_mpv_of_long_word_commutes}
    Let $A$ be injective after blocking $L_0>0$ sites. If $X$ commutes with
    every word product of some positive length $m$, then
    \[
        \Gamma_N(X)\in\spn\{V^{(N)}(A)\}.
    \]
\end{theorem}

\begin{proof}
    \leanok
    Lemma~\ref{lem:word_commutation_multiple_lengths} amplifies the length-$m$
    commutation hypothesis to length $L_0\,m \ge L_0$. Then
    Theorem~\ref{thm:ground_space_map_mpv_of_long_word_commutes} applies.
\end{proof}

\begin{theorem}[MPS-line containment from two one-sided equations]
    \label{thm:ground_space_map_mpv_of_common_middle}
    \lean{MPSTensor.groundSpaceMap_mem_mpvSubmodule_of_isNBlkInjective_of_two_sided_middle_compatibility}
    \leanok
    \uses{lem:common_middle_same_witness_commutation,
        thm:ground_space_map_mpv_of_positive_word_commutes}
    Let $A$ be injective after blocking $L_0>0$ sites. Suppose there are matrices
    $Y_\mu$ such that
    \[
        A^\mu A^b X = Y_\mu A^b,
        \qquad
        X A^a A^\mu = A^a Y_\mu
    \]
    for every pair of physical letters $a,b$. Then
    \[
        \Gamma_N(X)\in\spn\{V^{(N)}(A)\}.
    \]
\end{theorem}

\begin{proof}
    \leanok
    The identities in~(\ref{eq:ph_common_middle_one_sided}) give commutation
    with all words of length $|\mu|+2$, a positive length. The positive-length
    containment theorem then applies.
\end{proof}

\begin{theorem}[Boundary equality implies containment in the MPS line]
    \label{thm:ground_space_map_mpv_of_wrapped_witness_comparison}
    \lean{MPSTensor.groundSpaceMap_mem_mpvSubmodule_of_isNBlkInjective_of_wrapped_witness_comparison}
    \leanok
    \uses{lem:two_sided_middle_of_wrapped_witness_comparison,
        thm:ground_space_map_mpv_of_common_middle}
    Let $A$ be injective after blocking $L_0>0$ sites, and fix a physical letter
    $\eta$ used outside the complementary sites. Suppose that, for every
    boundary condition $\tau$ and every physical letter $j$,
    \[
        C^+_\tau A^jX=Y^+_\tau A^j,
        \qquad
        XA^jC^-_\tau=A^jY^-_\tau,
    \]
    where $C^+_\tau$ and $C^-_\tau$ are the complementary words exposed by the
    two boundary-crossing cyclic windows. Suppose also that the two boundary
    matrices satisfy
    \[
        Y^+_{\tau^+_\eta(\mu)} = Y^-_{\tau^-_\eta(\mu)}
    \]
    for every word $\mu$ on the complementary sites, then
    \[
        \Gamma_N(X)\in\spn\{V^{(N)}(A)\}.
    \]
\end{theorem}

\begin{proof}
    \leanok
    Lemma~\ref{lem:two_sided_middle_of_wrapped_witness_comparison} gives a
    family \(Y_\mu\) satisfying
    \[
        A^\mu A^bX=Y_\mu A^b,
        \qquad
        XA^aA^\mu=A^aY_\mu.
    \]
    Theorem~\ref{thm:ground_space_map_mpv_of_common_middle} applies.
\end{proof}

\begin{theorem}[Closure-property step for periodic chains]
    \label{thm:conditional_normal_range_reduction_witness_comparison}
    \lean{MPSTensor.chainGroundSpace_le_mpvSubmodule_of_isNBlkInjective_of_wrapped_witness_comparison}
    \leanok
    \uses{thm:cyclic_normal_open_chain_range_reduction,
        thm:reduced_wrapped_boundary_compatibilities,
        thm:ground_space_map_mpv_of_wrapped_witness_comparison}
    Let $A$ be injective after blocking $L_0>0$ sites and let $D \ge 1$. Assume
    $N \ge 2$ and $L_0 < L \le N$. Suppose that, for every physical letter
    $\eta$ used in the two boundary conditions and every boundary matrix
    obtained from an $N$-site chain ground state, the identities
    \[
        A^\mu A^j X = Y^+_{\tau^+_\eta(\mu)}A^j,
        \qquad
        X A^j A^\mu = A^jY^-_{\tau^-_\eta(\mu)}
    \]
    hold for every physical letter $j$ and every complementary word $\mu$, and
    are supplemented, for every $\mu$, by
    \[
        Y^+_{\tau^+_\eta(\mu)}
        =
        Y^-_{\tau^-_\eta(\mu)}.
    \]
    Then
    \[
        \mathcal G_{N,L}(A)\subseteq \spn\{V^{(N)}(A)\}.
    \]
\end{theorem}

\begin{proof}
    \leanok
    \uses{thm:cyclic_normal_open_chain_range_reduction,
        thm:reduced_wrapped_boundary_compatibilities,
        thm:ground_space_map_mpv_of_wrapped_witness_comparison}
    The cyclic-to-open-chain reduction writes any chain ground state as
    $\Gamma_N(X)$. The two boundary-crossing cyclic windows give
    \[
        A^\mu A^jX=Y^+_{\tau^+_\eta(\mu)}A^j,
        \qquad
        XA^jA^\mu=A^jY^-_{\tau^-_\eta(\mu)}.
    \]
    Together with
    \[
        Y^+_{\tau^+_\eta(\mu)}
        =
        Y^-_{\tau^-_\eta(\mu)},
    \]
    Theorem~\ref{thm:ground_space_map_mpv_of_wrapped_witness_comparison}
    finishes.
\end{proof}

\begin{lemma}[Products with one-site tensors determine the witness]
    \label{lem:wrapped_mirror_right_products_determine}
    \lean{MPSTensor.wrapped_mirror_witness_agree_of_right_products}
    \leanok
    \uses{lem:one_sided_boundary_witness_unique}
    Let $A$ be $L_0$-block-injective with $L_0>0$.  Fix two matrix
    families $Y^+$ and $Y^-$ and the two reindexed boundary conditions
    $\tau^+_\eta(\mu)$ and $\tau^-_\eta(\mu)$.  If, for every physical letter
    $j$,
    \[
        Y^+_{\tau^+_\eta(\mu)} A^j
        =
        Y^-_{\tau^-_\eta(\mu)} A^j,
    \]
    then
    \[
        Y^+_{\tau^+_\eta(\mu)}
        =
        Y^-_{\tau^-_\eta(\mu)}.
    \]
\end{lemma}

\begin{proof}
    \leanok
    Apply Lemma~\ref{lem:one_sided_boundary_witness_unique} to the two
    matrices $Y^+_{\tau^+_\eta(\mu)}$ and $Y^-_{\tau^-_\eta(\mu)}$.
\end{proof}

\begin{lemma}[First-letter restrictions determine a vector]
    \label{lem:first_letter_restrictions_determine}
    \lean{MPSTensor.eq_of_forall_restrictFirst_eq}
    \leanok
    Let $R_j$ be restriction to first letter $j$:
    \[
        (R_j\phi)(\sigma_1,\ldots,\sigma_L)
        =
        \phi(j,\sigma_1,\ldots,\sigma_L).
    \]
    For $\phi,\psi\in(\mathbb C^d)^{\otimes(L+1)}$,
    \[
        \bigl(\forall j,\ R_j\phi=R_j\psi\bigr)
        \Longrightarrow
        \phi=\psi .
    \]
\end{lemma}

\begin{proof}
    \leanok
    At $(\sigma_0,\sigma_1,\ldots,\sigma_L)$ the required equality is
    \[
        \phi(\sigma_0,\sigma_1,\ldots,\sigma_L)
        =
        (R_{\sigma_0}\phi)(\sigma_1,\ldots,\sigma_L)
        =
        (R_{\sigma_0}\psi)(\sigma_1,\ldots,\sigma_L)
        =
        \psi(\sigma_0,\sigma_1,\ldots,\sigma_L).
    \]
\end{proof}

\begin{lemma}[Auxiliary restriction equality for the closure property]
    \label{lem:boundary_restriction_from_fixed_letters}
    \lean{MPSTensor.closure_property_boundary_restriction_eq_of_fixed_boundary_letters}
    \leanok
    \uses{lem:first_letter_restrictions_determine, rem:cyclic_window_operators}
    Let $L_0>0$ and $L_0\le M$, and set
    \[
        B^+_{\eta,\mu}(\psi)
        =
        \operatorname{Res}^{\tau^+_{\eta}(\mu)}_{M,L_0+1}(\psi),
        \qquad
        B^-_{\eta,\mu}(\psi)
        =
        \operatorname{Res}^{\tau^-_{\eta}(\mu)}_{M+1-L_0,L_0+1}(\psi).
    \]
    With $R_j$ as in Lemma~\ref{lem:first_letter_restrictions_determine},
    \[
        \bigl(\forall j,\ R_jB^+_{\eta,\mu}(\psi)
        =
        R_jB^-_{\eta,\mu}(\psi)\bigr)
        \Longrightarrow
        B^+_{\eta,\mu}(\psi)=B^-_{\eta,\mu}(\psi).
    \]
\end{lemma}

\begin{proof}
    \leanok
    Apply Lemma~\ref{lem:first_letter_restrictions_determine} to the two
    length-$(L_0+1)$ restrictions, using the displayed family of first-letter
    equalities.
\end{proof}

\begin{lemma}[Closed-boundary restriction from right products]
    \label{lem:closed_boundary_restriction_from_right_products}
    \lean{MPSTensor.closure_property_boundary_restriction_eq_of_first_products}
    \leanok
    \uses{rem:cyclic_window_operators, lem:first_letter_restrictions_determine}
    Let $A$ be an MPS tensor, let $L_0>0$, and let $L_0\le M$. Suppose the two
    length-$(L_0+1)$ restrictions at the closed boundary are represented by
    boundary matrices $Y^+_{\eta,\mu}$ and $Y^-_{\eta,\mu}$:
    \[
        B^+_{\eta,\mu}(\psi)=\Gamma_{L_0+1}(Y^+_{\eta,\mu}),
        \qquad
        B^-_{\eta,\mu}(\psi)=\Gamma_{L_0+1}(Y^-_{\eta,\mu}).
    \]
    If, for every physical letter $j$,
    \[
        Y^+_{\eta,\mu}A^j
        =
        Y^-_{\eta,\mu}A^j,
    \]
    then
    \[
        B^+_{\eta,\mu}(\psi)=B^-_{\eta,\mu}(\psi).
    \]
\end{lemma}

\begin{proof}
    \leanok
    Fix $j$ and take the first-letter restriction of the two displayed
    length-$(L_0+1)$ restrictions. These are
    \[
        \Gamma_{L_0}(Y^+_{\eta,\mu}A^j),
        \qquad
        \Gamma_{L_0}(Y^-_{\eta,\mu}A^j).
    \]
    The one-site product equality gives, for every $j$,
    \[
        \Gamma_{L_0}(Y^+_{\eta,\mu}A^j)
        =
        \Gamma_{L_0}(Y^-_{\eta,\mu}A^j).
    \]
    Lemma~\ref{lem:first_letter_restrictions_determine} identifies the two
    original restrictions.
\end{proof}

\begin{lemma}[Boundary restrictions from first-letter products]
    \label{lem:closed_boundary_restrictions_from_first_letter_products}
    \lean{MPSTensor.closure_property_boundary_restrictions_eq_of_first_products}
    \leanok
    \uses{lem:closed_boundary_restriction_from_right_products}
    Let \(L_0>0\) and \(L_0\le M\). Let \(Y_i(\tau)\) represent the
    length-\((L_0+1)\) cyclic restrictions of \(\psi\). Fix a complementary
    word \(\mu\). If, for every boundary letter \(\eta\) and physical letter
    \(j\),
    \[
        Y_M(\tau^+_\eta(\mu))A^j
        =
        Y_{M+1-L_0}(\tau^-_\eta(\mu))A^j,
    \]
    then for every \(\eta\),
    \[
        \operatorname{Res}^{\tau^+_\eta(\mu)}_{M,L_0+1}(\psi)
        =
        \operatorname{Res}^{\tau^-_\eta(\mu)}_{M+1-L_0,L_0+1}(\psi).
    \]
\end{lemma}

\begin{proof}
    \leanok
    Apply Lemma~\ref{lem:closed_boundary_restriction_from_right_products} for
    each boundary letter \(\eta\), using the corresponding first-letter product
    equations.
\end{proof}

\begin{lemma}[Restriction form of the closure property]
    \label{lem:closure_property_boundary_restriction_chain_ground_space}
    \lean{MPSTensor.closure_property_boundary_restriction_eq_of_chainGroundSpace}
    \leanok
    \uses{def:chain_ground_space, rem:cyclic_window_operators,
        lem:chain_ground_space_window_witnesses,
        lem:closed_boundary_restriction_from_right_products,
        lem:boundary_closing_right_products_chain_ground_space}
    Let $A$ be $L_0$-block-injective with $L_0>0$ and $D\ge1$, let
    $L_0\le M$, and let
    \[
        \psi\in\mathcal G_{M+1,L_0+1}(A),
        \qquad
        \psi=\Gamma_{M+1}(X).
    \]
    For every boundary letter $\eta$ and complementary word $\mu$, the local
    coordinate restriction equality representing the boundary-closing step of
    \cite[lines~2078--2079]{Cirac2021Matrix} is
    \[
        \operatorname{Res}^{\tau^+_{\eta}(\mu)}_{M,L_0+1}(\psi)
        =
        \operatorname{Res}^{\tau^-_{\eta}(\mu)}_{M+1-L_0,L_0+1}(\psi).
    \]
\end{lemma}

\begin{proof}
    \notready
    The window-witness lemma gives matrices
    \(Y_M(\tau^+_\eta(\mu))\) and
    \(Y_{M+1-L_0}(\tau^-_\eta(\mu))\) such that
    \[
        \operatorname{Res}^{\tau^+_\eta(\mu)}_{M,L_0+1}(\psi)
        =
        \Gamma_{L_0+1}\!\left(Y_M(\tau^+_\eta(\mu))\right),
        \qquad
        \operatorname{Res}^{\tau^-_\eta(\mu)}_{M+1-L_0,L_0+1}(\psi)
        =
        \Gamma_{L_0+1}\!\left(Y_{M+1-L_0}(\tau^-_\eta(\mu))\right).
    \]
    Lemma~\ref{lem:boundary_closing_right_products_chain_ground_space}
    gives, for every $j$,
    \[
        Y_M(\tau^+_\eta(\mu))A^j
        =
        Y_{M+1-L_0}(\tau^-_\eta(\mu))A^j .
    \]
    Lemma~\ref{lem:closed_boundary_restriction_from_right_products} identifies
    the two length-$(L_0+1)$ restrictions.
\end{proof}

\begin{lemma}[Auxiliary product from closing-boundary restrictions]
    \label{lem:closure_property_auxiliary_product_from_closing_restrictions}
    \lean{MPSTensor.closure_property_auxiliary_boundary_product_eq_of_closing_restrictions}
    \leanok
    \uses{def:l_blk_injective, def:ground_space_map,
        rem:cyclic_window_operators, lem:ground_space_map_injective_blk,
        lem:cyclic_window_boundary_matrix_transport,
        lem:wrapped_middle_backgrounds}
    Let $A$ be $L_0$-block-injective with $L_0>0$, and let $L_0\le M$.
    Let $\psi$ be a state on the chain of length $M+1$, and let
    $Y_i(\tau)$ be matrices satisfying
    \[
        \operatorname{Res}^{\tau}_{i,L_0+1}(\psi)
        =
        \Gamma_{L_0+1}(Y_i(\tau)).
    \]
    Suppose that, for every outside letter $\eta$,
    \[
        \operatorname{Res}^{\tau^+_\eta(\mu)}_{M,L_0+1}(\psi)
        =
        \operatorname{Res}^{\tau^-_\eta(\mu)}_{M+1-L_0,L_0+1}(\psi).
    \]
    Then there are boundary conditions $\rho^+_{j,\sigma}$ and
    $\rho^-_{j,\sigma}$ such that, for every complementary position $k$,
    \[
        \rho^+_{j,\sigma}(k+L_0)=\mu_k,
        \qquad
        \rho^-_{j,\sigma}(k+1)=\mu_k,
    \]
    and, for every physical letter $j$ and every word $\sigma$ of length
    $L_0$,
    \[
        Y_M(\rho^+_{j,\sigma})A^jA^\sigma
        =
        Y_{M+1-L_0}(\rho^-_{j,\sigma})A^jA^\sigma.
    \]
\end{lemma}

\begin{proof}
    \leanok
    For the product indexed by $j$ and $\sigma$, specialize the preceding
    equality at the outside letter $\eta=j$. Take
    $\rho^+_{j,\sigma}=\tau^+_j(\mu)$ and
    $\rho^-_{j,\sigma}=\tau^-_j(\mu)$.  The complement equations are the two
    identities of Lemma~\ref{lem:wrapped_middle_backgrounds}.  Fix $j$.  Applying
    the first-letter restriction to the displayed equality gives
    \[
        \Gamma_{L_0}\!\left(Y_M(\tau^+_j(\mu))A^j\right)
        =
        \Gamma_{L_0}\!\left(Y_{M+1-L_0}(\tau^-_j(\mu))A^j\right),
    \]
    where the two sides are identified by
    Lemma~\ref{lem:cyclic_window_boundary_matrix_transport}.  By
    Lemma~\ref{lem:ground_space_map_injective_blk},
    \[
        Y_M(\tau^+_j(\mu))A^j
        =
        Y_{M+1-L_0}(\tau^-_j(\mu))A^j.
    \]
    Right multiplication by $A^\sigma$ gives the asserted product equation.
\end{proof}

\begin{lemma}[Auxiliary product from right products]
    \label{lem:closure_property_auxiliary_product_from_right_products}
    \lean{MPSTensor.closure_property_auxiliary_boundary_product_eq_of_right_products}
    \leanok
    \uses{lem:closed_boundary_restriction_from_right_products,
        lem:closure_property_auxiliary_product_from_closing_restrictions}
    Let $A$ be $L_0$-block-injective with $L_0>0$, and let $L_0\le M$.
    Let $\psi$ be a state on the chain of length $M+1$, and let
    $Y_i(\tau)$ be matrices satisfying
    \[
        \operatorname{Res}^{\tau}_{i,L_0+1}(\psi)
        =
        \Gamma_{L_0+1}(Y_i(\tau)).
    \]
    Fix a complementary word $\mu$ of length $M+1-(L_0+1)$.
    Suppose that, for every boundary letter $\eta$ and every physical letter
    $j$,
    \[
        Y_M(\tau^+_{\eta}(\mu))A^j
        =
        Y_{M+1-L_0}(\tau^-_{\eta}(\mu))A^j .
    \]
    Then there are boundary conditions $\rho^+_{j,\sigma}$ and
    $\rho^-_{j,\sigma}$ such that, for every complementary position $k$,
    \[
        \rho^+_{j,\sigma}(k+L_0)=\mu_k,
        \qquad
        \rho^-_{j,\sigma}(k+1)=\mu_k,
    \]
    and, for every physical letter $j$ and every word $\sigma$ of length
    $L_0$,
    \[
        Y_M(\rho^+_{j,\sigma})A^jA^\sigma
        =
        Y_{M+1-L_0}(\rho^-_{j,\sigma})A^jA^\sigma.
    \]
\end{lemma}

\begin{proof}
    \leanok
    For each boundary letter $\eta$,
    Lemma~\ref{lem:closed_boundary_restriction_from_right_products} gives
    \[
        \operatorname{Res}^{\tau^+_\eta(\mu)}_{M,L_0+1}(\psi)
        =
        \operatorname{Res}^{\tau^-_\eta(\mu)}_{M+1-L_0,L_0+1}(\psi).
    \]
    Applying these equalities,
    Lemma~\ref{lem:closure_property_auxiliary_product_from_closing_restrictions}
    gives the displayed boundary conditions and product equation.
\end{proof}

\begin{lemma}[Opposite-boundary comparison after multiplication by words]
    \label{lem:closure_property_opposite_boundary_comparison_from_right_words}
    \lean{MPSTensor.closure_property_mirror_right_product_eq_of_right_word_products}
    \leanok
    \uses{def:l_blk_injective, lem:padded_complement_annihilation}
    Let $A$ be $L_0$-block-injective with $L_0>0$, and let $L_0\le M$.
    Fix a complementary word $\mu$, a matrix $X$, and a family of matrices
    $Y_i(\tau)$. Suppose that, for every boundary letter $\eta$, every
    physical letter $j$, and every word $\sigma$ of length $L_0$,
    \[
        Y_{M+1-L_0}(\tau^-_\eta(\mu))A^jA^\sigma
        =
        A^\mu A^jXA^\sigma .
    \]
    Then, for every $\eta$ and $j$,
    \[
        Y_{M+1-L_0}(\tau^-_\eta(\mu))A^j
        =
        A^\mu A^jX .
    \]
\end{lemma}

\begin{proof}
    \leanok
    Fix $\eta$ and $j$, and set
    \[
        Z_{\eta,j}
        =
        Y_{M+1-L_0}(\tau^-_\eta(\mu))A^j-A^\mu A^jX .
    \]
    The hypothesis gives, for every $\sigma$ with $|\sigma|=L_0$,
    \[
        Z_{\eta,j}A^\sigma=0.
    \]
    Since $A$ is $L_0$-block-injective, the products $A^\sigma$ with
    $|\sigma|=L_0$ span the full matrix algebra.  Applying
    Lemma~\ref{lem:padded_complement_annihilation} gives $Z_{\eta,j}=0$.
\end{proof}

\begin{lemma}[Opposite-boundary comparison after left multiplication by words]
    \label{lem:closure_property_opposite_boundary_comparison_from_left_words}
    \lean{MPSTensor.closure_property_mirror_padded_products_of_left_word_products}
    \leanok
    \uses{def:l_blk_injective, lem:padded_complement_annihilation}
    Let $A$ be $L_0$-block-injective with $L_0>0$, and let $L_0\le M$.
    Fix a complementary word $\mu$, a matrix $X$, and a family of matrices
    $Y_i(\tau)$. Suppose that, for every boundary letter $\eta$, every physical
    letter $j$, and every pair of words $\sigma,\alpha$ of length $L_0$,
    \[
        A^\alpha\,
        \left(Y_{M+1-L_0}(\tau^-_\eta(\mu))A^jA^\sigma\right)
        =
        A^\alpha\,\left(A^\mu A^jXA^\sigma\right).
    \]
    Then, for every $\eta$, $j$, and $\sigma$,
    \[
        Y_{M+1-L_0}(\tau^-_\eta(\mu))A^jA^\sigma
        =
        A^\mu A^jXA^\sigma .
    \]
\end{lemma}

\begin{proof}
    \leanok
    Fix $\eta$, $j$, and $\sigma$, and set
    \[
        Z_{\eta,j,\sigma}
        =
        Y_{M+1-L_0}(\tau^-_\eta(\mu))A^jA^\sigma
        -
        A^\mu A^jXA^\sigma .
    \]
    The hypothesis says that \(A^\alpha Z_{\eta,j,\sigma}=0\) for every word
    \(\alpha\) of length \(L_0\).  The left-handed form of
    Lemma~\ref{lem:padded_complement_annihilation} gives
    \(Z_{\eta,j,\sigma}=0\).
\end{proof}

\begin{lemma}[Auxiliary product from the opposite-boundary comparison]
    \label{lem:closure_property_auxiliary_product_from_opposite_boundary_words}
    \lean{MPSTensor.closure_property_auxiliary_boundary_product_eq_of_mirror_padded_products}
    \leanok
    \uses{lem:closure_property_opposite_boundary_comparison_from_right_words,
        lem:closure_property_auxiliary_product_from_right_products}
    Let $A$ be $L_0$-block-injective with $L_0>0$, and let $L_0\le M$.
    Let $\psi$ be a state on the chain of length $M+1$, and let
    $Y_i(\tau)$ be matrices satisfying
    \[
        \operatorname{Res}^{\tau}_{i,L_0+1}(\psi)
        =
        \Gamma_{L_0+1}(Y_i(\tau)).
    \]
    Fix a complementary word $\mu$ and a matrix $X$. Suppose that, for every
    boundary letter $\eta$ and every physical letter $j$,
    \[
        Y_M(\tau^+_{\eta}(\mu))A^j
        =
        A^\mu A^jX .
    \]
    Suppose also that, for every $\eta$, $j$, and every word $\sigma$ of length
    $L_0$,
    \[
        Y_{M+1-L_0}(\tau^-_{\eta}(\mu))A^jA^\sigma
        =
        A^\mu A^jXA^\sigma .
    \]
    Then there are boundary conditions $\rho^+_{j,\sigma}$ and
    $\rho^-_{j,\sigma}$ such that, for every complementary position $k$,
    \[
        \rho^+_{j,\sigma}(k+L_0)=\mu_k,
        \qquad
        \rho^-_{j,\sigma}(k+1)=\mu_k,
    \]
    and, for every physical letter $j$ and every word $\sigma$ of length
    $L_0$,
    \[
        Y_M(\rho^+_{j,\sigma})A^jA^\sigma
        =
        Y_{M+1-L_0}(\rho^-_{j,\sigma})A^jA^\sigma.
    \]
\end{lemma}

\begin{proof}
    \leanok
    Lemma~\ref{lem:closure_property_opposite_boundary_comparison_from_right_words}
    gives, for every $\eta$ and $j$,
    \[
        Y_{M+1-L_0}(\tau^-_{\eta}(\mu))A^j
        =
        A^\mu A^jX .
    \]
    Combining this equation with
    \[
        Y_M(\tau^+_{\eta}(\mu))A^j=A^\mu A^jX
    \]
    gives
    \[
        Y_M(\tau^+_{\eta}(\mu))A^j
        =
        Y_{M+1-L_0}(\tau^-_{\eta}(\mu))A^j .
    \]
    Lemma~\ref{lem:closure_property_auxiliary_product_from_right_products}
    then gives the displayed boundary conditions and auxiliary product.
\end{proof}

\begin{lemma}[Auxiliary product from the left-multiplied opposite-boundary comparison]
    \label{lem:closure_property_auxiliary_product_from_opposite_boundary_left_words}
    \lean{MPSTensor.closure_property_auxiliary_boundary_product_eq_of_mirror_left_word_products}
    \leanok
    \uses{lem:closure_property_opposite_boundary_comparison_from_left_words,
        lem:closure_property_auxiliary_product_from_opposite_boundary_words}
    Let $A$ be $L_0$-block-injective with $L_0>0$, and let $L_0\le M$.
    Let $\psi$ be a state on the chain of length $M+1$, and let
    $Y_i(\tau)$ be matrices satisfying
    \[
        \operatorname{Res}^{\tau}_{i,L_0+1}(\psi)
        =
        \Gamma_{L_0+1}(Y_i(\tau)).
    \]
    Fix a complementary word $\mu$ and a matrix $X$. Suppose that, for every
    boundary letter $\eta$ and every physical letter $j$,
    \[
        Y_M(\tau^+_\eta(\mu))A^j=A^\mu A^jX .
    \]
    Suppose also that, for every $\eta$, $j$, and every pair of words
    $\sigma,\alpha$ of length $L_0$,
    \[
        A^\alpha\,
        \left(Y_{M+1-L_0}(\tau^-_\eta(\mu))A^jA^\sigma\right)
        =
        A^\alpha\,\left(A^\mu A^jXA^\sigma\right).
    \]
    Then there are boundary conditions $\rho^+_{j,\sigma}$ and
    $\rho^-_{j,\sigma}$ satisfying the complementary-word equations
    \[
        \rho^+_{j,\sigma}(k+L_0)=\mu_k,
        \qquad
        \rho^-_{j,\sigma}(k+1)=\mu_k
    \]
    and the product equation
    \[
        Y_M(\rho^+_{j,\sigma})A^jA^\sigma
        =
        Y_{M+1-L_0}(\rho^-_{j,\sigma})A^jA^\sigma .
    \]
\end{lemma}

\begin{proof}
    \leanok
    Lemma~\ref{lem:closure_property_opposite_boundary_comparison_from_left_words}
    gives
    \[
        Y_{M+1-L_0}(\tau^-_\eta(\mu))A^jA^\sigma
        =
        A^\mu A^jXA^\sigma .
    \]
    Lemma~\ref{lem:closure_property_auxiliary_product_from_opposite_boundary_words}
    then gives the displayed boundary conditions and product equation.
\end{proof}

\begin{lemma}[Auxiliary product from the left-multiplied closure-property comparison]
    \label{lem:closure_property_auxiliary_boundary_product_ground_space_left_words}
    \lean{MPSTensor.closure_property_auxiliary_boundary_product_eq_of_groundSpaceMap_left_words}
    \leanok
    \uses{lem:closure_property_boundary_one_sided_products_ground_space_map,
        lem:closure_property_auxiliary_product_from_opposite_boundary_left_words}
    Let $A$ be $L_0$-block-injective with $L_0>0$, $L_0\le M$, and
    $D\ge1$. Let
    \[
        \psi=\Gamma_{M+1}(X)
    \]
    be an open-chain representation, and let $Y_i(\tau)$ represent the
    length-$(L_0+1)$ cyclic restrictions of $\psi$. Fix a complementary word
    $\mu$. Suppose that, for every boundary letter $\eta$, physical letter
    $j$, and length-$L_0$ words $\sigma,\alpha$,
    \[
        A^\alpha
        \left(Y_{M+1-L_0}(\tau^-_\eta(\mu))A^jA^\sigma\right)
        =
        A^\alpha
        \left(A^\mu A^jXA^\sigma\right).
    \]
    Then there are boundary conditions $\rho^+_{j,\sigma}$ and
    $\rho^-_{j,\sigma}$ satisfying the complementary-word equations
    \[
        \rho^+_{j,\sigma}(k+L_0)=\mu_k,
        \qquad
        \rho^-_{j,\sigma}(k+1)=\mu_k
    \]
    and the product equation
    \[
        Y_M(\rho^+_{j,\sigma})A^jA^\sigma
        =
        Y_{M+1-L_0}(\rho^-_{j,\sigma})A^jA^\sigma .
    \]
\end{lemma}

\begin{proof}
    \leanok
    Lemma~\ref{lem:closure_property_boundary_one_sided_products_ground_space_map}
    gives the last-boundary equation
    \[
        Y_M(\tau^+_\eta(\mu))A^j=A^\mu A^jX .
    \]
    The assumed left-multiplied coordinate equation and
    Lemma~\ref{lem:closure_property_auxiliary_product_from_opposite_boundary_left_words}
    then give the displayed boundary conditions and product equation.
    % Source note: arXiv:2011.12127, Section IV.C, lines 2078--2079,
    % says that the corresponding boundary-closing argument is analogous.
    % The assumed equation is the coordinate form isolated in
    % docs/paper-gaps/cpgsv21_normal_range_reduction.tex.
\end{proof}

\begin{lemma}[Boundary restrictions from the left-multiplied comparison]
    \label{lem:closure_property_boundary_restrictions_ground_space_map_left_words}
    \lean{MPSTensor.closure_property_boundary_restrictions_eq_of_groundSpaceMap_left_words}
    \leanok
    \uses{lem:closure_property_boundary_one_sided_products_ground_space_map,
        lem:closure_property_opposite_boundary_comparison_from_left_words,
        lem:closure_property_opposite_boundary_comparison_from_right_words,
        lem:closed_boundary_restrictions_from_first_letter_products}
    Let $A$ be $L_0$-block-injective with $L_0>0$, $L_0\le M$, and
    $D\ge1$. Let
    \[
        \psi=\Gamma_{M+1}(X)
    \]
    be an open-chain representation, and let $Y_i(\tau)$ represent the
    length-$(L_0+1)$ cyclic restrictions of $\psi$. Fix a complementary word
    $\mu$. Suppose that, for every boundary letter $\eta$, physical letter $j$,
    and length-$L_0$ words $\alpha,\sigma$,
    \[
        A^\alpha
        \left(Y_{M+1-L_0}(\tau^-_\eta(\mu))A^jA^\sigma\right)
        =
        A^\alpha
        \left(A^\mu A^jXA^\sigma\right).
    \]
    Then, for every boundary letter $\eta$,
    \[
        \operatorname{Res}^{\tau^+_\eta(\mu)}_{M,L_0+1}(\psi)
        =
        \operatorname{Res}^{\tau^-_\eta(\mu)}_{M+1-L_0,L_0+1}(\psi).
    \]
\end{lemma}

\begin{proof}
    \leanok
    The assumed equations and
    Lemma~\ref{lem:closure_property_opposite_boundary_comparison_from_left_words}
    give, for every $\eta$, $j$, and $\sigma$,
    \[
        Y_{M+1-L_0}(\tau^-_\eta(\mu))A^jA^\sigma
        =
        A^\mu A^jXA^\sigma .
    \]
    Lemma~\ref{lem:closure_property_opposite_boundary_comparison_from_right_words}
    removes the right word:
    \[
        Y_{M+1-L_0}(\tau^-_\eta(\mu))A^j
        =
        A^\mu A^jX .
    \]
    The last-boundary equation from
    Lemma~\ref{lem:closure_property_boundary_one_sided_products_ground_space_map}
    gives
    \[
        Y_M(\tau^+_\eta(\mu))A^j=A^\mu A^jX .
    \]
    Hence the first-letter products agree, and
    Lemma~\ref{lem:closed_boundary_restrictions_from_first_letter_products}
    identifies the two restrictions.
\end{proof}

\begin{lemma}[Equality of the boundary-closing cyclic restrictions]
    \label{lem:closure_property_boundary_restrictions_ground_space_map}
    \lean{MPSTensor.closure_property_boundary_restrictions_eq_of_groundSpaceMap}
    \leanok
    \uses{def:ground_space_map, rem:cyclic_window_operators,
        lem:closure_property_boundary_one_sided_products_ground_space_map,
        lem:closure_property_opposite_boundary_comparison_from_left_words,
        lem:closure_property_opposite_boundary_comparison_from_right_words,
        lem:closed_boundary_restrictions_from_first_letter_products,
        lem:closure_property_boundary_restrictions_ground_space_map_left_words}
    Let $A$ be $L_0$-block-injective with $L_0>0$, $L_0\le M$, and
    $D\ge1$. Let
    \[
        \psi=\Gamma_{M+1}(X)
    \]
    be an open-chain representation. Assume that for all $i$ and $\tau$,
    \[
        \operatorname{Res}^{\tau}_{i,L_0+1}(\psi)\in G_{L_0+1}(A).
    \]
    Fix a complementary word $\mu$. For every boundary letter $\eta$, choose the
    two local boundary conditions so that, for each complementary index $k$,
    \[
        \tau^+_\eta(\mu)_{k+L_0}=\mu_k,
        \qquad
        \tau^-_\eta(\mu)_{k+1}=\mu_k,
    \]
    and put the letter $\eta$ on the remaining sites. These two choices have the
    same complementary word. The notation $\tau^\pm_\eta(\mu)$ is a local
    coordinate device for the boundary-closing comparison; it is not notation
    used in the source. In these coordinates the boundary-closing step is the
    equality
    \[
        \operatorname{Res}^{\tau^+_\eta(\mu)}_{M,L_0+1}(\psi)
        =
        \operatorname{Res}^{\tau^-_\eta(\mu)}_{M+1-L_0,L_0+1}(\psi).
    \]
\end{lemma}

\begin{proof}
    \notready
    Choose matrices $Y_i(\tau)$ representing the length-$(L_0+1)$
    restrictions. It remains to establish the following coordinate comparison:
    for every boundary letter $\eta$, physical letter $j$, and length-$L_0$
    words $\alpha,\sigma$,
    \[
        A^\alpha
        \left(Y_{M+1-L_0}(\tau^-_\eta(\mu))A^jA^\sigma\right)
        =
        A^\alpha
        \left(A^\mu A^jXA^\sigma\right).
    \]
    Lemma~\ref{lem:closure_property_boundary_restrictions_ground_space_map_left_words}
    records that these equations imply the equality of the two restrictions.
    The displayed left-multiplied family is the coordinate form of the
    boundary-closing sentence in \cite[lines~2078--2079]{Cirac2021Matrix} that
    is still open in the formalization; see
    \path{docs/paper-gaps/cpgsv21_normal_range_reduction.tex}.
\end{proof}

\begin{lemma}[First-letter products from equal boundary restrictions]
    \label{lem:closure_property_boundary_first_products_restrictions}
    \lean{MPSTensor.closure_property_boundary_first_products_of_restrictions}
    \leanok
    \uses{def:ground_space_map, rem:cyclic_window_operators}
    Let $A$ be $L_0$-block-injective with $L_0>0$ and $L_0\le M$. Let
    $\psi$ be a vector on $M+1$ sites, and suppose the two length-$(L_0+1)$
    boundary restrictions are represented by
    \[
        \Gamma_{L_0+1}(Y_M(\tau^+_\eta(\mu))),
        \qquad
        \Gamma_{L_0+1}(Y_{M+1-L_0}(\tau^-_\eta(\mu))).
    \]
    If these two restrictions are equal, then for every physical letter $j$,
    \[
        Y_M(\tau^+_\eta(\mu))A^j
        =
        Y_{M+1-L_0}(\tau^-_\eta(\mu))A^j .
    \]
    Thus the first-letter restriction of the preceding equality is precisely the
    displayed product equation.
\end{lemma}

\begin{proof}
    \leanok
    Restricting both sides to the first physical letter $j$ gives
    \[
        \Gamma_{L_0}\!\left(Y_M(\tau^+_\eta(\mu))A^j\right)
        =
        \Gamma_{L_0}\!\left(Y_{M+1-L_0}(\tau^-_\eta(\mu))A^j\right).
    \]
    Since $A$ is $L_0$-block-injective, $\Gamma_{L_0}$ is injective, and the
    displayed product equality follows.
\end{proof}

\begin{lemma}[Left-multiplied comparison from equal boundary restrictions]
    \label{lem:closure_property_wrapped_mirror_left_word_products_boundary_restrictions}
    \lean{MPSTensor.closure_property_wrapped_mirror_left_word_products_of_boundary_restrictions}
    \leanok
    \uses{def:ground_space_map, rem:cyclic_window_operators,
        lem:closure_property_boundary_first_products_restrictions}
    Let $A$ be $L_0$-block-injective with $L_0>0$ and $L_0\le M$. Let
    $\psi$ be a vector on $M+1$ sites, and let $Y_i(\tau)$ represent the
    length-$(L_0+1)$ cyclic restrictions of $\psi$. Assume that for every
    boundary letter $\eta$,
    \[
        \operatorname{Res}^{\tau^+_\eta(\mu)}_{M,L_0+1}(\psi)
        =
        \operatorname{Res}^{\tau^-_\eta(\mu)}_{M+1-L_0,L_0+1}(\psi).
    \]
    Then for every boundary letter $\eta$, physical letter $j$, and
    length-$L_0$ words $\alpha,\sigma$,
    \[
        A^\alpha
        \left(Y_{M+1-L_0}(\tau^-_\eta(\mu))A^jA^\sigma\right)
        =
        A^\alpha
        \left(Y_M(\tau^+_\eta(\mu))A^jA^\sigma\right).
    \]
\end{lemma}

\begin{proof}
    \leanok
    The preceding lemma gives
    \[
        Y_M(\tau^+_\eta(\mu))A^j
        =
        Y_{M+1-L_0}(\tau^-_\eta(\mu))A^j .
    \]
    Multiplying on the left by $A^\alpha$ and on the right by $A^\sigma$
    gives the desired comparison.
\end{proof}

\begin{lemma}[Left-multiplied comparison of boundary matrices]
    \label{lem:closure_property_wrapped_mirror_left_word_products_ground_space}
    \lean{MPSTensor.closure_property_wrapped_mirror_left_word_products_of_groundSpaceMap}
    \leanok
    \uses{def:ground_space_map, rem:cyclic_window_operators,
        lem:closure_property_boundary_restrictions_ground_space_map,
        lem:closure_property_wrapped_mirror_left_word_products_boundary_restrictions}
    Let $A$ be $L_0$-block-injective with $L_0>0$, $L_0\le M$, and
    $D\ge1$. Let
    \[
        \psi=\Gamma_{M+1}(X)
    \]
    be an open-chain representation, and let $Y_i(\tau)$ represent the
    length-$(L_0+1)$ cyclic restrictions of $\psi$. Fix a complementary word
    $\mu$. Then for every boundary letter $\eta$, physical letter $j$, and
    length-$L_0$ words $\alpha,\sigma$,
    \[
        A^\alpha
        \left(Y_{M+1-L_0}(\tau^-_\eta(\mu))A^jA^\sigma\right)
        =
        A^\alpha
        \left(Y_M(\tau^+_\eta(\mu))A^jA^\sigma\right).
    \]
\end{lemma}

\begin{proof}
    \notready
    The representations by the matrices $Y_i(\tau)$ supply the local
    ground-space hypothesis for the preceding lemma.
    Once the preceding boundary-restriction equality is proved, it gives
    \[
        \operatorname{Res}^{\tau^+_\eta(\mu)}_{M,L_0+1}(\psi)
        =
        \operatorname{Res}^{\tau^-_\eta(\mu)}_{M+1-L_0,L_0+1}(\psi).
    \]
    Applying the left-multiplied comparison for equal boundary restrictions
    gives
    \[
        A^\alpha
        \left(Y_{M+1-L_0}(\tau^-_\eta(\mu))A^jA^\sigma\right)
        =
        A^\alpha
        \left(Y_M(\tau^+_\eta(\mu))A^jA^\sigma\right).
    \]
\end{proof}

\begin{lemma}[Left-multiplied boundary-closing comparison]
    \label{lem:closure_property_mirror_left_word_products_ground_space}
    \lean{MPSTensor.closure_property_mirror_left_word_products_of_groundSpaceMap}
    \leanok
    \uses{def:ground_space_map, rem:cyclic_window_operators,
        lem:closure_property_wrapped_mirror_left_word_products_ground_space,
        lem:closure_property_boundary_one_sided_products_ground_space_map}
    Let $A$ be $L_0$-block-injective with $L_0>0$, $L_0\le M$, and
    $D\ge1$. Let
    \[
        \psi=\Gamma_{M+1}(X)
    \]
    be an open-chain representation, and let $Y_i(\tau)$ represent the
    length-$(L_0+1)$ cyclic restrictions of $\psi$. Fix a complementary word
    $\mu$. Then for every boundary letter $\eta$, physical letter $j$, and
    length-$L_0$ words $\alpha,\sigma$,
    \[
        A^\alpha
        \left(Y_{M+1-L_0}(\tau^-_\eta(\mu))A^jA^\sigma\right)
        =
        A^\alpha
        \left(A^\mu A^jXA^\sigma\right).
    \]
\end{lemma}

\begin{proof}
    \notready
    Once the comparison of the two cyclic restrictions is available, it gives
    \[
        A^\alpha
        \left(Y_{M+1-L_0}(\tau^-_\eta(\mu))A^jA^\sigma\right)
        =
        A^\alpha
        \left(Y_M(\tau^+_\eta(\mu))A^jA^\sigma\right).
    \]
    The last-boundary one-sided equation gives
    \[
        Y_M(\tau^+_\eta(\mu))A^j=A^\mu A^jX.
    \]
    Substitution gives the displayed identity.
    % The cited paragraph does not display the first comparison. It is
    % the coordinate reconstruction isolated here for the closing-boundary step
    % of \cite[lines~2078--2079]{Cirac2021Matrix}. Remaining gap recorded in
    % docs/paper-gaps/cpgsv21_normal_range_reduction.tex.
\end{proof}

\begin{lemma}[Letterwise restriction comparison at the closing boundary]
    \label{lem:closure_property_fixed_boundary_letter_chain_ground_space}
    \lean{MPSTensor.closure_property_fixed_boundary_letter_eq_of_chainGroundSpace}
    \leanok
    \uses{def:chain_ground_space, rem:cyclic_window_operators,
        lem:chain_ground_space_window_witnesses,
        lem:boundary_closing_right_products_chain_ground_space}
    Let $A$ be $L_0$-block-injective with $L_0>0$ and $D\ge1$, let
    $L_0\le M$, and let
    \[
        \psi\in\mathcal G_{M+1,L_0+1}(A),
        \qquad
        \psi=\Gamma_{M+1}(X).
    \]
    Let $R_j$ denote restriction to first physical letter $j$.  For every
    outside letter $\eta$, complementary word $\mu$, and physical letter $j$,
    one has
    \[
        R_j\operatorname{Res}^{\tau^+_\eta(\mu)}_{M,L_0+1}(\psi)
        =
        R_j\operatorname{Res}^{\tau^-_\eta(\mu)}_{M+1-L_0,L_0+1}(\psi).
    \]
    This is the letterwise coordinate form of the local closure-property comparison
    corresponding to \cite[lines~2078--2079]{Cirac2021Matrix}.
\end{lemma}

\begin{proof}
    \notready
    Let $Y_i(\tau)$ be the window witnesses of
    Lemma~\ref{lem:chain_ground_space_window_witnesses}.  By
    Lemma~\ref{lem:boundary_closing_right_products_chain_ground_space},
    \[
        Y_M(\tau^+_\eta(\mu))A^j
        =
        Y_{M+1-L_0}(\tau^-_\eta(\mu))A^j .
    \]
    Fixing the first physical letter in the two length-$(L_0+1)$ windows gives
    \[
    \begin{aligned}
        R_j\operatorname{Res}^{\tau^+_\eta(\mu)}_{M,L_0+1}(\psi)
        &=
        \Gamma_{L_0}\!\left(Y_M(\tau^+_\eta(\mu))A^j\right),\\
        R_j\operatorname{Res}^{\tau^-_\eta(\mu)}_{M+1-L_0,L_0+1}(\psi)
        &=
        \Gamma_{L_0}\!\left(Y_{M+1-L_0}(\tau^-_\eta(\mu))A^j\right).
    \end{aligned}
    \]
    The two restrictions are therefore equal.
\end{proof}

\begin{lemma}[Auxiliary boundary-condition product equation]
    \label{lem:closure_property_auxiliary_boundary_product_chain_ground_space}
    \lean{MPSTensor.closure_property_auxiliary_boundary_product_eq_of_groundSpaceMap}
    \leanok
    \uses{def:l_blk_injective, def:ground_space_map, rem:cyclic_window_operators,
        lem:closure_property_boundary_one_sided_products_ground_space_map,
        lem:closure_property_auxiliary_product_from_opposite_boundary_words}
    Let $A$ be $L_0$-block-injective with $L_0>0$ and $D\ge1$. Let $\psi$ be a
    state on the chain of length $M+1$ with open-chain representation
    $\psi=\Gamma_{M+1}(X)$.  Let $Y_i(\tau)$ be matrices satisfying
    \[
        \operatorname{Res}^{\tau}_{i,L_0+1}(\psi)
        =
        \Gamma_{L_0+1}(Y_i(\tau)).
    \]
    If $L_0\le M$, then for every complementary word $\mu$ there are boundary
    conditions $\rho^+_{j,\sigma}$ and $\rho^-_{j,\sigma}$, depending on the
    physical letter $j$ and the word $\sigma$ of length $L_0$, such that, for
    every complementary position $k$,
    \[
        \rho^+_{j,\sigma}(k+L_0)=\mu_k,
        \qquad
        \rho^-_{j,\sigma}(k+1)=\mu_k
    \]
    and, for every $j$ and $\sigma$,
    \[
        Y_M(\rho^+_{j,\sigma})A^jA^\sigma
        =
        Y_{M+1-L_0}(\rho^-_{j,\sigma})A^jA^\sigma
    \]
\end{lemma}

\begin{proof}
    \notready
    Lemma~\ref{lem:closure_property_boundary_one_sided_products_ground_space_map}
    gives
    \[
        Y_M(\tau^+_{\eta}(\mu))A^j
        =
        A^\mu A^j X .
    \]
    By
    Lemma~\ref{lem:closure_property_auxiliary_product_from_opposite_boundary_words},
    it remains to prove, for every boundary letter $\eta$, every physical
    letter $j$, and every word $\sigma$ of length $L_0$,
    \[
        Y_{M+1-L_0}(\tau^-_{\eta}(\mu))A^jA^\sigma
        =
        A^\mu A^jXA^\sigma .
    \]
    % The cited paragraph does not display this formula. It is the coordinate
    % reconstruction isolated here for the closing-boundary step of
    % \cite[lines~2078--2079]{Cirac2021Matrix}. Remaining gap recorded in
    % docs/paper-gaps/cpgsv21_normal_range_reduction.tex.
\end{proof}

\begin{lemma}[Product equality after closing the boundary]
    \label{lem:boundary_closing_word_equation_chain_ground_space}
    \lean{MPSTensor.closure_property_boundary_closing_product_eq_of_chainGroundSpace}
    \leanok
    \uses{def:chain_ground_space, def:ground_space_map, rem:cyclic_window_operators}
    Let $A$ be $L_0$-block-injective with $L_0>0$ and $D\ge1$. Let
    $\psi\in\mathcal G_{M+1,L_0+1}(A)$ have an open-chain representation
    $\psi=\Gamma_{M+1}(X)$.  Let $Y_i(\tau)$ be matrices satisfying
    \[
        \operatorname{Res}^{\tau}_{i,L_0+1}(\psi)
        =
        \Gamma_{L_0+1}(Y_i(\tau)).
    \]
    If $L_0\le M$, then for every
    boundary letter $\eta$, complementary word $\mu$, physical letter $j$, and
    word $\sigma$ of length $L_0$,
    \[
        Y_M(\tau^{+}_{\eta}(\mu))A^jA^\sigma
        =
        Y_{M+1-L_0}(\tau^{-}_{\eta}(\mu))A^jA^\sigma .
    \]
\end{lemma}

\begin{proof}
    \notready
    \uses{lem:closure_property_auxiliary_boundary_product_chain_ground_space,
        thm:boundary_closing_product_pointwise_compatible_boundary_assignments}
    Lemma~\ref{lem:closure_property_auxiliary_boundary_product_chain_ground_space}
    gives boundary conditions $\rho^+_{j,\sigma}$ and
    $\rho^-_{j,\sigma}$ satisfying
    \[
        \rho^+_{j,\sigma}(k+L_0)=\mu_k,
        \qquad
        \rho^-_{j,\sigma}(k+1)=\mu_k
    \]
    and
    \[
        Y_M(\rho^+_{j,\sigma})A^jA^\sigma
        =
        Y_{M+1-L_0}(\rho^-_{j,\sigma})A^jA^\sigma .
    \]
    Theorem~\ref{thm:boundary_closing_product_pointwise_compatible_boundary_assignments}
    replaces these auxiliary conditions by the displayed boundary conditions:
    \[
        Y_M(\tau^{+}_{\eta}(\mu))A^jA^\sigma
        =
        Y_{M+1-L_0}(\tau^{-}_{\eta}(\mu))A^jA^\sigma,
    \]
    which is the product equality needed after closing the boundary.
\end{proof}

\begin{lemma}[Right-product equality after closing the boundary]
    \label{lem:boundary_closing_right_products_chain_ground_space}
    \lean{MPSTensor.closure_property_boundary_right_products_eq_of_chainGroundSpace}
    \leanok
    \uses{lem:boundary_closing_word_equation_chain_ground_space,
        lem:padded_complement_annihilation}
    Let $A$ be $L_0$-block-injective with $L_0>0$ and $D\ge1$. Let
    $\psi\in\mathcal G_{M+1,L_0+1}(A)$ have an open-chain representation
    $\psi=\Gamma_{M+1}(X)$.  Let $Y_i(\tau)$ be matrices satisfying
    \[
        \operatorname{Res}^{\tau}_{i,L_0+1}(\psi)
        =
        \Gamma_{L_0+1}(Y_i(\tau)).
    \]
    If $L_0\le M$, then for every
    letter $\eta$ used in the two boundary conditions, every word $\mu$ on
    the complementary sites, and every physical letter $j$,
    \[
        Y_M(\tau^{+}_{\eta}(\mu))A^j
        =
        Y_{M+1-L_0}(\tau^{-}_{\eta}(\mu))A^j.
    \]
    This is the coordinate one-site product equality used to compare the two
    boundary matrices after closing the boundary.
\end{lemma}

\begin{proof}
    \notready
    Fix $j$ and set
    \[
        Z_j=
        Y_M(\tau^{+}_{\eta}(\mu))A^j
        -
        Y_{M+1-L_0}(\tau^{-}_{\eta}(\mu))A^j .
    \]
    Lemma~\ref{lem:boundary_closing_word_equation_chain_ground_space}
    gives, after subtracting the two sides, for every word $\sigma$ of length $L_0$,
    \[
        Z_jA^\sigma=0.
    \]
    Since $A$ is $L_0$-block-injective,
    \[
        \operatorname{span}\{A^\sigma:|\sigma|=L_0\}=M_D(\mathbb C).
    \]
    Applying Lemma~\ref{lem:padded_complement_annihilation} gives $Z_j=0$, hence
    \[
        Y_M(\tau^{+}_{\eta}(\mu))A^j
        =
        Y_{M+1-L_0}(\tau^{-}_{\eta}(\mu))A^j.
    \]
\end{proof}

\begin{theorem}[Boundary conditions agree in the closure property]
    \label{thm:wrapped_mirror_witness_agree_chain_ground_space}
    \lean{MPSTensor.wrapped_mirror_witness_agree_of_chainGroundSpace}
    \leanok
    \uses{thm:reduced_wrapped_boundary_compatibilities,
        lem:wrapped_mirror_right_products_determine,
        lem:boundary_closing_right_products_chain_ground_space}
    Let $A$ be $L_0$-block-injective with $L_0>0$ and $D\ge1$, let $N\ge2$
    and $L_0<L\le N$, and let
    $\psi\in\mathcal G_{N,L}(A)$ have an open-chain representation
    $\psi=\Gamma_N(X)$.  Suppose $Y^{+}$ and $Y^{-}$ are two families of
    matrices obtained from the two boundary-crossing local constraints used
    in the closure property, and that they satisfy,
    for every physical letter $j$ and every boundary condition $\tau$,
    \[
        A^{\tau_{L_0}\cdots\tau_{N-2}} A^j X
        =
        Y^{+}_{\tau} A^j,
        \qquad
        X A^j A^{\tau_1\cdots\tau_{N-L_0-1}}
        =
        A^j Y^{-}_{\tau}.
    \]
    For every physical letter
    $\eta\in\{0,\ldots,d-1\}$ and every word $\mu$ on the complementary sites,
    define two boundary conditions by
    \[
        \tau^{+}_{\eta}(\mu)_i =
        \begin{cases}
            \mu_{i-L_0}, & L_0\le i<N-1,\\
            \eta,        & \text{otherwise},
        \end{cases}
        \qquad
        \tau^{-}_{\eta}(\mu)_i =
        \begin{cases}
            \mu_{i-1}, & 1\le i<N-L_0,\\
            \eta,      & \text{otherwise}.
        \end{cases}
    \]
    Then the two families agree on these two boundary conditions:
    \[
        Y^{+}_{\tau^{+}_{\eta}(\mu)}
        =
        Y^{-}_{\tau^{-}_{\eta}(\mu)}.
    \]
    This is the boundary-matrix comparison required by the closure property.
\end{theorem}

\begin{proof}
    \notready
    The reduced cyclic-window compatibilities at the boundary give, for every
    physical letter $j$,
    \[
        A^{\mu} A^j X
        =
        Y^{+}_{\tau^{+}_{\eta}(\mu)} A^j,
        \qquad
        X A^j A^{\mu}
        =
        A^j Y^{-}_{\tau^{-}_{\eta}(\mu)}.
    \]
    By Lemma~\ref{lem:wrapped_mirror_right_products_determine}, it is enough to
    prove, for every $j$,
    \[
        Y^{+}_{\tau^{+}_{\eta}(\mu)} A^j
        =
        Y^{-}_{\tau^{-}_{\eta}(\mu)} A^j.
    \]
    Write $N=M+1$. The two boundary-crossing supports start at
    $M$ and $M+1-L_0$.  Let $Y_i(\tau)$
    denote the cyclic-window witness supplied by
    Lemma~\ref{lem:chain_ground_space_window_witnesses} at the window beginning at
    $i$.  The first step identifies the two abstract boundary-condition families
    with the two boundary-crossing witnesses:
    \[
        Y^{+}_{\tau^{+}_{\eta}(\mu)}
        =
        Y_M(\tau^{+}_{\eta}(\mu)),
        \qquad
        Y^{-}_{\tau^{-}_{\eta}(\mu)}
        =
        Y_{M+1-L_0}(\tau^{-}_{\eta}(\mu)).
    \]
    Together with the preceding identifications,
    Lemma~\ref{lem:boundary_closing_right_products_chain_ground_space}
    gives, for every physical letter $j$,
    \[
        Y^{+}_{\tau^{+}_{\eta}(\mu)} A^j
        =
        Y^{-}_{\tau^{-}_{\eta}(\mu)} A^j.
    \]
    Lemma~\ref{lem:wrapped_mirror_right_products_determine} then gives the
    desired equality of boundary matrices.
\end{proof}

\begin{theorem}[Chain ground space is spanned by the MPS vector in the injective case]%
    \label{thm:chain_ground_space_eq_mpv_blk}
    \lean{MPSTensor.chainGroundSpace_eq_mpvSubmodule}
    \leanok
    \uses{def:chain_ground_space, def:mpv_submodule, def:injective}
    If $A$ is injective, $D \ge 1$, $N \ge 2$, and $1 < L \le N$, then
    \[
        \mathcal G_{N,L}(A) = \spn\bigl\{V^{(N)}(A)\bigr\}.
    \]
\end{theorem}

\begin{proof}
    \leanok
    \uses{lem:contiguous_mem_ground_space, thm:boundary_matrix_commutes,
        thm:ground_space_map_injective, lem:mpv_mem_chain_ground_space}
    Let $\psi \in \mathcal G_{N,L}(A)$. The non-wrapping window conditions and
    Lemma~\ref{lem:contiguous_mem_ground_space} imply $\psi \in G_{N}(A)$, so
    $\psi = \Gamma_{N}(X)$ for some boundary matrix $X$. The periodic boundary condition now
    applies Theorem~\ref{thm:boundary_matrix_commutes} and forces $X$ to commute with
    every letter $A^{j}$. Since $A$ is injective, the letters generate the full matrix
    algebra, hence $X$ is scalar. Therefore $\psi$ is proportional to $V^{(N)}(A)$, while
    the reverse inclusion is Lemma~\ref{lem:mpv_mem_chain_ground_space}.
\end{proof}

\begin{theorem}[Closure-property step toward uniqueness of the ground state]
    \label{thm:normal_range_reduction_containment}
    \lean{MPSTensor.chainGroundSpace_le_mpvSubmodule_of_normal_range_reduction}
    \leanok
    \uses{def:chain_ground_space, def:mpv_submodule, def:normal,
        def:l_blk_injective, thm:cyclic_normal_open_chain_range_reduction,
        thm:reduced_wrapped_boundary_compatibilities,
        thm:wrapped_mirror_witness_agree_chain_ground_space,
        thm:conditional_normal_range_reduction_witness_comparison}
    If $A$ is normal, $D \ge 1$, and injective after blocking $L_0>0$ sites, with
    $N \ge 2$ and $L_0 < L \le N$, then
    \[
        \mathcal G_{N,L}(A) \subseteq \spn\bigl\{V^{(N)}(A)\bigr\}.
    \]
\end{theorem}

\begin{proof}
    \notready
    \uses{thm:reduced_wrapped_boundary_compatibilities,
        thm:conditional_normal_range_reduction_witness_comparison,
        lem:padded_complement_annihilation}
    For a chain ground state, the open-chain containment gives
    $\psi=\Gamma_N(X)$. The two boundary-crossing local constraints give
    \[
        A^\mu A^j X = Y^+_{\tau^+_\eta(\mu)}A^j,
        \qquad
        X A^j A^\mu = A^jY^-_{\tau^-_\eta(\mu)}.
    \]
    The remaining closed-boundary comparison is the equality
    \[
        Y^+_{\tau^+_\eta(\mu)} = Y^-_{\tau^-_\eta(\mu)}.
    \]
    These equations imply
    \[
        \psi\in\spn\{V^{(N)}(A)\}.
    \]
\end{proof}

\begin{theorem}[Chain ground space is spanned by the MPS vector in the normal case]%
    \label{thm:chain_ground_space_eq_mpv_normal}
    \lean{MPSTensor.chainGroundSpace_eq_mpvSubmodule_normal}
    \leanok
    \uses{thm:normal_range_reduction_containment, lem:mpv_mem_chain_ground_space,
        def:normal, def:l_blk_injective}
    If $A$ is normal, $D \ge 1$, and injective after blocking $L_0>0$ sites, with
    $N \ge 2$ and $L_0 < L \le N$, then
    \[
        \mathcal G_{N,L}(A) = \spn\bigl\{V^{(N)}(A)\bigr\}.
    \]
\end{theorem}

\begin{theorem}[Unique ground state on the periodic chain]%
    \label{thm:ground_space_unique_periodic}
    \lean{MPSTensor.groundSpace_unique_periodic}
    \leanok
    \uses{thm:chain_ground_space_eq_mpv_blk, def:has_unique_ground_state,
        lem:mpv_eq_ground_space_map_one, thm:ground_space_map_injective}
    For an injective tensor $A$ with $D \ge 1$, if $N \ge 2$ and
    $1 < L \le N$, then the periodic chain ground space is one-dimensional.
\end{theorem}

\begin{proof}
    \leanok
    \uses{thm:chain_ground_space_eq_mpv_blk, lem:mpv_eq_ground_space_map_one,
        thm:ground_space_map_injective}
    By Theorem~\ref{thm:chain_ground_space_eq_mpv_blk}, the ground space is the line
    $\spn\bigl\{V^{(N)}(A)\bigr\}$. It remains to show that the MPV is nonzero. If
    $V^{(N)}(A)=0$, then Lemma~\ref{lem:mpv_eq_ground_space_map_one} gives
    \[
        \Gamma_{N}(\Id) = \Gamma_{N}(0).
    \]
    Since $A$ is injective and $N > 0$, Theorem~\ref{thm:ground_space_map_injective}
    forces $\Id = 0$, a contradiction. Therefore the spanning line is one-dimensional.
\end{proof}

\begin{theorem}[Unique ground state for block-injective tensors]%
    \label{thm:unique_gs_injective}
    \lean{MPSTensor.parentHamiltonian_unique_gs_injective}
    \leanok
    \uses{thm:chain_ground_space_eq_mpv_normal}
    If $A$ is $L_0$-block-injective with $L_0 > 0$ and $D \ge 1$, the parent Hamiltonian
    with interaction range $2L_0$ has a unique ground state on every periodic chain
    of $N \ge 2L_0$ sites.
\end{theorem}

\begin{theorem}[Unique ground state for normal tensors at range $L_0+1$]%
    \label{thm:unique_gs_normal}
    \lean{MPSTensor.parentHamiltonian_unique_gs_normal}
    \leanok
    \uses{thm:unique_gs_injective}
    If $A$ is normal (hence $L_0$-block-injective for some $L_0 > 0$),
    $D \ge 1$, then the parent Hamiltonian with interaction range $L_0 + 1$
    has a unique ground state on every periodic chain of $N \ge L_0+1$ sites.
\end{theorem}

\section{Commuting parent Hamiltonians}\label{sec:commuting_parent_ham}

A parent Hamiltonian has \emph{commuting} local terms when all translated
interaction projectors commute with each other.

\begin{definition}[Commuting parent Hamiltonian]\label{def:is_commuting_parent_ham}
    \lean{MPSTensor.IsCommutingParentHam}
    \leanok
    \uses{def:local_term_parent}
    The parent Hamiltonian with block length~$L$ on $N$ sites has
    \emph{commuting} local terms if
    $h_i h_j = h_j h_i$ for all $i,j \in \{0,\ldots,N{-}1\}$.
    See Ref.~\cite[Definition~3.9]{Cirac2016MPDO_arXiv}.
\end{definition}

\begin{definition}[Nearest-neighbor commuting parent Hamiltonian]\label{def:is_nncph}
    \lean{MPSTensor.IsNNCPH}
    \leanok
    \uses{def:is_commuting_parent_ham}
    A \emph{nearest-neighbor commuting parent Hamiltonian} (NNCPH) is a
    commuting parent Hamiltonian with block length $L = 2$.
    See Ref.~\cite[Definition~3.9]{Cirac2016MPDO_arXiv}.
\end{definition}

\begin{remark}[Elementary consequences of the commuting definition]
    \lean{MPSTensor.IsNNCPH.isCommutingParentHam}
    \lean{MPSTensor.IsCommutingParentHam.symm}
    \lean{MPSTensor.IsCommutingParentHam.ham_comm_localTerm}
    \leanok
    The nearest-neighbor condition is simply the length-$2$ specialization of the general
    commuting condition; the latter is symmetric in the two sites, and it implies that the
    full parent Hamiltonian commutes with each local interaction term.
\end{remark}

\begin{theorem}[RFP implies NNCPH]\label{thm:rfp_implies_nncph}
    \lean{MPSTensor.rfp_implies_nncph}
    \leanok
    \uses{def:is_rfp, def:is_nncph, def:normal}
    If $A$ is a normal renormalization fixed point, then for every chain length
    $N \ge 2$ the parent Hamiltonian with $L = 2$ has commuting local terms.
    This implication is cited here from the product-of-entangled-pairs
    description in Appendix~B of
    Theorem~3.10 in~\cite{Cirac2016MPDO_arXiv}.
\end{theorem}

\begin{remark}[Product-of-pairs equations for NNCPH]%
    \label{rem:product_pair_projectors}
    \lean{MPSTensor.productPairWindow}
    \lean{MPSTensor.productPairWindow_one}
    \lean{MPSTensor.productPairState}
    \lean{MPSTensor.productPairState_one}
    \lean{MPSTensor.HasProductPairMPV}
    \lean{MPSTensor.HasProductPairLocalProjectors}
    \lean{MPSTensor.HasProductPairLocalProjectors.isNNCPH}
    \lean{MPSTensor.ProductPairBridge}
    \lean{MPSTensor.ProductPairBridge.localTerm_idempotent}
    \lean{MPSTensor.ProductPairBridge.commuting_twoSite_localTerms}
    \lean{MPSTensor.ProductPairBridge.isNNCPH}
    \lean{MPSTensor.AppendixBStructuralData}
    \lean{MPSTensor.AppendixBStructuralData.exists_ofRFP}
    \lean{MPSTensor.AppendixBStructuralData.ofRFP}
    \lean{MPSTensor.AppendixBStructuralData.twoSiteAmplitude}
    \lean{MPSTensor.AppendixBStructuralData.coreTensor}
    \lean{MPSTensor.AppendixBStructuralData.coreTensor_apply}
    \lean{MPSTensor.AppendixBStructuralData.gaugeEquiv_coreTensor}
    \lean{MPSTensor.AppendixBStructuralData.mpv_eq_coreTensor}
    \lean{MPSTensor.AppendixBStructuralData.twoSiteAmplitude_eq_mpv}
    \lean{MPSTensor.AppendixBStructuralData.twoSiteAmplitude_eq_coreTensor_mpv}
    \lean{MPSTensor.AppendixBStructuralData.mpv_eq_productPairState_one}
    \lean{MPSTensor.AppendixBProductPairExtraction}
    \lean{MPSTensor.AppendixBProductPairExtraction.ofCoreTensorFactorization}
    \lean{MPSTensor.AppendixBProductPairExtraction.toProductPairBridge}
    \lean{MPSTensor.AppendixBProductPairExtraction.commuting_twoSite_localTerms}
    \lean{MPSTensor.commuting_twoSite_localTerms_of_rfp_of_appendixBExtraction}
    \lean{MPSTensor.rfp_implies_nncph_of_appendixBExtraction}
    \leanok
    The Appendix~B structural form supplies
    \[
        A^i=X\Lambda U^iX^{-1}.
    \]
    The change of basis does not alter periodic MPS coefficients, so the
    corresponding two-site amplitude is
    \[
        \phi_2(a,b)
        =
        V^{(2)}(A)(a,b)
        =
        V^{(2)}(\Lambda U)(a,b).
    \]
    The missing product-of-pairs input is the even-chain factorization
    \[
        V^{(2n)}(A)(\sigma)
        =
        \prod_{p=0}^{n-1}
        \phi_2\bigl(\sigma_{2p},\sigma_{2p+1}\bigr).
    \]
    The nearest-neighbor parent projectors must also be identified with
    idempotents $p_i$ satisfying
    \[
        h_i(A,2)=p_i,\qquad p_i^2=p_i,\qquad p_ip_j=p_jp_i.
    \]
    These equations give
    \[
        h_i(A,2)h_j(A,2)=h_j(A,2)h_i(A,2),
    \]
    which is the NNCPH conclusion for every finite chain. Thus the remaining
    mathematical input after the Appendix~B structural form is exactly the
    factorization above together with the two-site projector identification.
\end{remark}

\begin{theorem}[NNCPH implies RFP]\label{thm:nncph_implies_rfp}
    \lean{MPSTensor.nncph_implies_rfp}
    \leanok
    \uses{def:is_rfp, def:is_nncph, def:normal}
    If $A$ is normal, in left-canonical form, and its parent Hamiltonian with
    $L = 2$ has commuting local terms for every chain length $N \ge 2$, then $A$
    is a renormalization fixed point. This implication is cited here from Beigi's
    ground-space
    characterization of one-dimensional commuting nearest-neighbor Hamiltonians; the
    left-canonical hypothesis fixes the normalization of the transfer map.
\end{theorem}

\section{Decorrelation and commuting parent Hamiltonians}%
\label{sec:decorrelation}

This section develops the abstract operator-algebraic formulation of
decorrelation and its connection to commuting parent Hamiltonians, following
\cite[Appendix~D, Section~D.2]{Cirac2016MPDO_arXiv}.

\begin{definition}[Decorrelation]\label{def:is_decorrelated}
    \lean{Decorrelation.IsDecorrelated}
    \leanok
    Given an idempotent endomorphism $P_K$ and sets of observables
    $\mathcal{O}_A$, $\mathcal{O}_B$, regions $A$ and $B$ are
    \emph{decorrelated} w.r.t.\ $K$ when
    $P_K \circ O_A \circ (1 - P_K) \circ O_B \circ P_K = 0$
    for all $O_A \in \mathcal{O}_A$, $O_B \in \mathcal{O}_B$.
\end{definition}

\begin{definition}[Commuting parent Hamiltonian structure]\label{def:has_commuting_parent_ham}
    \lean{Decorrelation.HasCommutingParentHam}
    \leanok
    A subspace $K$ (given by its idempotent endomorphism $P_K$) has a
    \emph{commuting parent Hamiltonian structure} if there exist idempotent
    endomorphisms $P_{AX}$ and $P_{XB}$ that commute and satisfy
    $P_{AX} \circ P_{XB} = P_K$.
\end{definition}

\begin{theorem}[Idempotence gives a trivial commuting parent Hamiltonian]%
    \label{thm:has_commuting_parent_ham_of_idem}
    \lean{Decorrelation.HasCommutingParentHam.ofIdem}
    \leanok
    \uses{def:has_commuting_parent_ham}
    If
    \[
        P_K \circ P_K = P_K,
    \]
    then $P_K$ has a commuting parent Hamiltonian structure.
\end{theorem}

\begin{proof}
    \leanok
    \uses{def:has_commuting_parent_ham}
    Take $P_{AX} = P_{XB} = P_K$. Their commutativity is immediate, and the
    product identity is exactly the assumed idempotence of $P_K$.
\end{proof}

\begin{lemma}[Cancellation for commuting idempotents]
    \label{lem:comp_complement_comm_zero}
    \lean{LinearMap.comp_complement_comm_zero}
    \leanok
    If $P$ and $Q$ are commuting idempotents, then
    \[
        Q \circ (1 - P \circ Q) \circ P = 0.
    \]
\end{lemma}

\begin{proof}
    \leanok
    Expand the middle factor, use commutativity to rewrite $QPQP$ as $QPPQ$, and
    then apply idempotence of $P$ and $Q$.
\end{proof}

\begin{theorem}[Idempotence of the product projector]%
    \label{thm:pk_idem}
    \lean{Decorrelation.HasCommutingParentHam.pK_idem}
    \leanok
    \uses{def:has_commuting_parent_ham}
    If $P_K = P_{AX} \circ P_{XB}$ is a commuting parent Hamiltonian
    decomposition, then
    \[
        P_K \circ P_K = P_K.
    \]
\end{theorem}

\begin{proof}
    \leanok
    \uses{def:has_commuting_parent_ham}
    Since $P_{AX}$ and $P_{XB}$ are idempotent and commute, their product
    $P_{AX} \circ P_{XB}$ is idempotent. Substituting
    $P_K = P_{AX} \circ P_{XB}$ gives the claim.
\end{proof}

\begin{theorem}[Absorption by the left projector]%
    \label{thm:pAX_comp_pK}
    \lean{Decorrelation.HasCommutingParentHam.pAX_comp_pK}
    \leanok
    \uses{def:has_commuting_parent_ham}
    If $P_K = P_{AX} \circ P_{XB}$ is a commuting parent Hamiltonian
    decomposition, then
    \[
        P_{AX} \circ P_K = P_K.
    \]
\end{theorem}

\begin{proof}
    \leanok
    \uses{def:has_commuting_parent_ham}
    Using $P_K = P_{AX} \circ P_{XB}$ and the idempotence of $P_{AX}$,
    \[
        P_{AX} \circ P_K
        = P_{AX} \circ P_{AX} \circ P_{XB}
        = P_{AX} \circ P_{XB}
        = P_K.
    \]
\end{proof}

\begin{theorem}[Absorption by the right projector]%
    \label{thm:pK_comp_pXB}
    \lean{Decorrelation.HasCommutingParentHam.pK_comp_pXB}
    \leanok
    \uses{def:has_commuting_parent_ham}
    If $P_K = P_{AX} \circ P_{XB}$ is a commuting parent Hamiltonian
    decomposition, then
    \[
        P_K \circ P_{XB} = P_K.
    \]
\end{theorem}

\begin{proof}
    \leanok
    \uses{def:has_commuting_parent_ham}
    Using $P_K = P_{AX} \circ P_{XB}$ and the idempotence of $P_{XB}$,
    \[
        P_K \circ P_{XB}
        = P_{AX} \circ P_{XB} \circ P_{XB}
        = P_{AX} \circ P_{XB}
        = P_K.
    \]
\end{proof}

\begin{theorem}[Cross-absorption by the right projector]%
    \label{thm:pXB_comp_pK}
    \lean{Decorrelation.HasCommutingParentHam.pXB_comp_pK}
    \leanok
    \uses{def:has_commuting_parent_ham}
    If $P_K = P_{AX} \circ P_{XB}$ is a commuting parent Hamiltonian
    decomposition, then
    \[
        P_{XB} \circ P_K = P_K.
    \]
\end{theorem}

\begin{proof}
    \leanok
    \uses{def:has_commuting_parent_ham}
    By commutativity,
    $P_{XB} \circ P_K = P_{XB} \circ P_{AX} \circ P_{XB}
        = P_{AX} \circ P_{XB} \circ P_{XB}$, and idempotence of $P_{XB}$
    yields $P_{XB} \circ P_K = P_K$.
\end{proof}

\begin{theorem}[Cross-absorption by the left projector]%
    \label{thm:pK_comp_pAX}
    \lean{Decorrelation.HasCommutingParentHam.pK_comp_pAX}
    \leanok
    \uses{def:has_commuting_parent_ham}
    If $P_K = P_{AX} \circ P_{XB}$ is a commuting parent Hamiltonian
    decomposition, then
    \[
        P_K \circ P_{AX} = P_K.
    \]
\end{theorem}

\begin{proof}
    \leanok
    \uses{def:has_commuting_parent_ham}
    By commutativity,
    $P_K \circ P_{AX} = P_{AX} \circ P_{XB} \circ P_{AX}
        = P_{AX} \circ P_{AX} \circ P_{XB}$, and idempotence of $P_{AX}$
    yields $P_K \circ P_{AX} = P_K$.
\end{proof}

\begin{theorem}[Reverse product identity]%
    \label{thm:reverse_product}
    \lean{Decorrelation.HasCommutingParentHam.reverse_product}
    \leanok
    \uses{def:has_commuting_parent_ham}
    If $P_K = P_{AX} \circ P_{XB}$ is a commuting parent Hamiltonian
    decomposition, then
    \[
        P_{XB} \circ P_{AX} = P_K.
    \]
\end{theorem}

\begin{proof}
    \leanok
    \uses{def:has_commuting_parent_ham}
    This is immediate from commutativity of $P_{AX}$ and $P_{XB}$ together
    with the identity $P_{AX} \circ P_{XB} = P_K$.
\end{proof}

\begin{theorem}[Commuting complementary projectors]%
    \label{thm:complement_comm}
    \lean{Decorrelation.HasCommutingParentHam.complement_comm}
    \leanok
    \uses{def:has_commuting_parent_ham}
    If $P_K = P_{AX} \circ P_{XB}$ is a commuting parent Hamiltonian
    decomposition and
    \[
        Q_{AX} = 1 - P_{AX}, \qquad Q_{XB} = 1 - P_{XB},
    \]
    then
    \[
        Q_{AX} \circ Q_{XB} = Q_{XB} \circ Q_{AX}.
    \]
\end{theorem}

\begin{proof}
    \leanok
    \uses{def:has_commuting_parent_ham}
    Expanding both products and using $P_{AX} \circ P_{XB}
        = P_{XB} \circ P_{AX}$ shows that the complementary projectors commute.
\end{proof}

\begin{theorem}[Commuting parent Hamiltonian implies decorrelation]%
    \label{thm:commuting_ham_is_decorrelated}
    \lean{Decorrelation.commutingHam_isDecorrelated}
    \leanok
    \uses{def:is_decorrelated, def:has_commuting_parent_ham,
        lem:comp_complement_comm_zero}
    If a subspace $K$ has a commuting parent Hamiltonian decomposition
    $P_K = P_{AX} \circ P_{XB}$ with $[P_{AX}, P_{XB}] = 0$, and
    $A$-observables commute with $P_{XB}$ while $B$-observables commute
    with $P_{AX}$, then regions $A$ and $B$ are decorrelated w.r.t.\ $K$.
    This is the backward direction of
    \cite[Appendix~D, Section~D.2, Proposition~D.3]{Cirac2016MPDO_arXiv}.
\end{theorem}

\begin{proof}
    \leanok
    The key identity is $P_{XB} \circ (1 - P_{AX} \circ P_{XB}) \circ P_{AX} = 0$.
    Using the commutation hypotheses to slide $O_A$ past $P_{XB}$ and $O_B$
    past $P_{AX}$, the five-fold composition factors through this zero term.
\end{proof}

\begin{theorem}[Commuting parent Hamiltonian implies decorrelation]%
    \label{thm:has_commuting_parent_ham_is_decorrelated}
    \lean{Decorrelation.HasCommutingParentHam.isDecorrelated}
    \leanok
    \uses{thm:commuting_ham_is_decorrelated, def:has_commuting_parent_ham, def:is_decorrelated}
    Corollary of Theorem~\ref{thm:commuting_ham_is_decorrelated}: if
    $P_K$ has a commuting parent Hamiltonian decomposition and observables
    respect locality, then regions $A$ and $B$ are decorrelated w.r.t.\ $K$.
\end{theorem}

\begin{lemma}[Frustration-free identity]\label{lem:frustration_free_ham_eq}
    \lean{LinearMap.frustration_free_ham_eq}
    \leanok
    For any endomorphisms $P$ and $Q$,
    \[
        (1-P) + (1-Q) - (1-P)\circ(1-Q) = 1 - P\circ Q.
    \]
\end{lemma}

\begin{proof}
    \leanok
    Expanding the left-hand side and cancelling gives the result.
\end{proof}

\begin{theorem}[Commuting parent Hamiltonian identity]%
    \label{thm:hamiltonian_eq}
    \lean{Decorrelation.HasCommutingParentHam.hamiltonian_eq}
    \leanok
    \uses{def:has_commuting_parent_ham, lem:frustration_free_ham_eq}
    If $P_K = P_{AX} \circ P_{XB}$ is a commuting parent Hamiltonian
    decomposition, then
    \[
        Q_{AX} + Q_{XB} - Q_{AX} \circ Q_{XB} = 1 - P_K,
    \]
    where $Q_{AX} = 1 - P_{AX}$ and $Q_{XB} = 1 - P_{XB}$ are the
    complementary projectors.
    This is the frustration-free form of the parent Hamiltonian from
    \cite[Appendix~D, Section~D.2]{Cirac2016MPDO_arXiv}.
\end{theorem}

\begin{proof}
    \leanok
    \uses{lem:frustration_free_ham_eq}
    Applying Lemma~\ref{lem:frustration_free_ham_eq} to $P_{AX}$ and
    $P_{XB}$, and substituting $P_{AX} \circ P_{XB} = P_K$.
\end{proof}

\begin{theorem}[Ground-space membership]\label{thm:mem_ground_iff}
    \lean{Decorrelation.HasCommutingParentHam.mem_ground_iff}
    \leanok
    \uses{def:has_commuting_parent_ham}
    If $P_K = P_{AX} \circ P_{XB}$ is a commuting parent Hamiltonian
    decomposition, then
    \[
        P_K v = v
        \quad\Longleftrightarrow\quad
        P_{AX} v = v \text{ and }  P_{XB} v = v.
    \]
    This is the algebraic form of equation~(D.2) from
    \cite{Cirac2016MPDO_arXiv}:
    $K_{AXB} = (K_{AX} \otimes H_B) \cap (H_A \otimes K_{XB})$.
\end{theorem}

\begin{proof}
    \leanok
    \uses{def:has_commuting_parent_ham}
    ($\Rightarrow$) If $P_K v = v$, then $P_{AX}(P_K v) = (P_{AX} \circ P_K) v
        = P_K v = v$ by left absorption, and similarly for $P_{XB}$.

    ($\Leftarrow$) If $P_{AX} v = v$ and $P_{XB} v = v$, then
    $P_K v = (P_{AX} \circ P_{XB}) v = P_{AX}(P_{XB} v) = P_{AX} v = v$.
\end{proof}

\section{Spectral gap}\label{sec:spectral_gap}

A frustration-free Hamiltonian is \emph{gapped} if the first excited
eigenvalue is bounded away from zero uniformly in the chain length $N$.
For MPS parent Hamiltonians, the spectral gap follows from a martingale
criterion due to Kastoryano and
Lucia~\cite{Kastoryano2018Martingale}.

\begin{definition}[Transported parent ground space]\label{def:parent_ground_space_es}
    \lean{MPSTensor.parentHamiltonianGroundSpaceES}
    \leanok
    \uses{def:parent_hamiltonian, def:nsite_space_parent}
    Denote by $\mathcal G_{N,L}^{\mathrm{ES}}(A)$ the parent ground space after
    identifying the coefficient space with its Euclidean $\ell^{2}$-model.
\end{definition}

\begin{definition}[Transported parent Hamiltonian]\label{def:parent_hamiltonian_es}
    \lean{MPSTensor.parentHamiltonianES}
    \leanok
    \uses{def:parent_hamiltonian, def:nsite_space_parent}
    Denote by $H_{N,L}^{\mathrm{ES}}(A)$ the conjugate of $H_{N}(A,L)$ by the
    canonical identification between the coefficient space and its Euclidean avatar.
\end{definition}

\begin{theorem}[Euclidean transport of the parent ground space]
    \label{thm:parent_ground_space_es_kernel}
    \lean{MPSTensor.parentHamiltonianGroundSpaceES_eq_ker_parentHamiltonianES}
    \leanok
    Under the canonical $\ell^2$ transport of the $N$-site state space, the
    transported parent ground space is exactly the kernel of the transported
    parent Hamiltonian.
\end{theorem}

\begin{proof}
    \leanok
    This is an immediate unraveling of the definitions: both constructions are
    obtained from the same parent Hamiltonian by conjugating with the canonical
    identification between the site space and its Euclidean-space avatar, so the
    transported ground space is precisely the transported kernel.
\end{proof}

\begin{lemma}[Membership in the transported parent ground space]
    \label{lem:mem_parent_ground_space_es}
    \lean{MPSTensor.mem_parentHamiltonianGroundSpaceES_iff}
    \leanok
    \uses{thm:parent_ground_space_es_kernel}
    A vector belongs to the transported parent ground space exactly when the transported
    parent Hamiltonian annihilates it.
\end{lemma}

\begin{remark}[Euclidean local projector ingredients]
    \label{rem:euclidean_local_projector_ingredients}
    \lean{MPSTensor.parentInteractionES}
    \lean{MPSTensor.cyclicRestrictES}
    \lean{MPSTensor.localTermES}
    \lean{MPSTensor.localTermESSummand}
    \leanok
    \uses{def:ground_space_es, rem:cyclic_window_operators}
    Let $P_L^{\mathrm{ES}}(A)$ denote the orthogonal projector onto
    $(G_L^{\mathrm{ES}}(A))^\perp$ on the Euclidean $L$-site space. For a cyclic
    window starting at $i$ and an outside configuration $\tau$, the transported
    restriction map $R_{i,\tau}$ evaluates an $N$-site Euclidean vector on that
    window. The transported local term $h_{i,\mathrm{ES}}$ is the corresponding
    $N$-site local projector, and $R_{i,\tau}^{\dagger} P_L^{\mathrm{ES}}(A) R_{i,\tau}$ is
    the basic positive summand in its cyclic-averaging formula.
\end{remark}

\begin{lemma}[Positivity of the Euclidean local summands]
    \label{lem:euclidean_local_projector_positive}
    \lean{MPSTensor.parentInteractionES_isPositive}
    \lean{MPSTensor.localTermESSummand_isPositive}
    \lean{MPSTensor.localTermESSummand_sum_isPositive}
    \leanok
    \uses{rem:euclidean_local_projector_ingredients}
    The Euclidean local projector $P_L^{\mathrm{ES}}(A)$ is positive. Consequently,
    every summand $R_{i,\tau}^{\dagger} P_L^{\mathrm{ES}}(A) R_{i,\tau}$ is positive, and so
    is the finite sum over $\tau$.
\end{lemma}

\begin{proof}
    \leanok
    Orthogonal projectors are positive. Conjugation by $R_{i,\tau}$ preserves
    positivity, and finite sums of positive operators remain positive.
\end{proof}

\begin{lemma}[Cyclic averaging for the transported ES local terms]
    \label{lem:local_term_es_average}
    \lean{MPSTensor.parentHamiltonianES_eq_sum_localTermES}
    \lean{MPSTensor.localTermES_eq_average_localTermESSummand}
    \leanok
    \uses{rem:euclidean_local_projector_ingredients, lem:euclidean_local_projector_positive}
    The transported parent Hamiltonian is the sum of the transported local terms
    $H_{N,\mathrm{ES}} = \sum_i h_{i,\mathrm{ES}}$, and each transported local term admits
    the exact cyclic-averaging formula
    \begin{equation}\label{eq:ph_cyclic_averaging}
        h_{i,\mathrm{ES}}
        = d^{-L} \sum_{\tau} R_{i,\tau}^{\dagger} P_L^{\mathrm{ES}}(A) R_{i,\tau}.
    \end{equation}
\end{lemma}

\begin{proof}
    \leanok
    Unfold the transported local term and compare it configuration-by-configuration
    with the cyclic restrictions. The $d^{-L}$ factor compensates for the $d^L$
    choices of the window coordinates in the ambient configuration parameter $\tau$.
\end{proof}

\begin{lemma}[Positivity of the transported ES Hamiltonian]
    \label{lem:transported_es_positive}
    \lean{MPSTensor.localTermES_isPositive}
    \lean{MPSTensor.parentHamiltonianES_isPositive}
    \leanok
    \uses{lem:local_term_es_average}
    Each transported local term $h_{i,\mathrm{ES}}$ is positive, hence the full
    transported parent Hamiltonian $H_{N,\mathrm{ES}}$ is positive.
\end{lemma}

\begin{proof}
    \leanok
    Apply the averaging formula~(\ref{eq:ph_cyclic_averaging}): $h_{i,\mathrm{ES}}$ is a non-negative scalar
    multiple of a finite sum of positive conjugates of $P_L^{\mathrm{ES}}(A)$, hence is
    itself positive. Summing over $i$ yields positivity of $H_{N,\mathrm{ES}}$.
\end{proof}

\begin{lemma}[Symmetric-projection structure of transported ES local terms]
    \label{lem:transported_es_local_projections}
    \lean{MPSTensor.parentInteractionES_isSymmetricProjection}
    \lean{MPSTensor.localTermES_isSymmetricProjection}
    \leanok
    \uses{rem:euclidean_local_projector_ingredients, lem:transported_es_positive}
    The Euclidean local projector $P_L^{\mathrm{ES}}(A)$ is a symmetric projection, and
    every transported local term $h_{i,\mathrm{ES}}$ on the $N$-site chain is a
    symmetric projection.
\end{lemma}

\begin{proof}
    \leanok
    \uses{lem:transported_es_positive}
    The local operator $P_L^{\mathrm{ES}}(A)$ is an orthogonal projection onto
    $(G_L^{\mathrm{ES}}(A))^\perp$. For $L\le N$, restricting
    $h_{i,\mathrm{ES}}v$ to a cyclic window gives
    $P_L^{\mathrm{ES}}(A)$ applied to the corresponding restriction of $v$; applying
    $h_{i,\mathrm{ES}}$ a second time therefore applies
    $(P_L^{\mathrm{ES}}(A))^2=P_L^{\mathrm{ES}}(A)$ on that window. The positivity
    result above supplies symmetry, and the case $L>N$ is the zero projection by
    definition.
\end{proof}

\begin{lemma}[Kernel of the Euclidean parent interaction]\label{lem:parent_interaction_es_kernel}
    \lean{MPSTensor.parentInteractionES_apply_eq_zero_iff}
    \leanok
    \uses{def:ground_space_es, rem:euclidean_local_projector_ingredients}
    The Euclidean local interaction $P_L^{\mathrm{ES}}(A)$ annihilates exactly the
    Euclidean local ground space:
    \[
        P_L^{\mathrm{ES}}(A)v=0 \quad\Longleftrightarrow\quad
        v\in G_L^{\mathrm{ES}}(A).
    \]
\end{lemma}

\begin{proof}
    \leanok
    Since $P_L^{\mathrm{ES}}(A)$ is the orthogonal projector onto
    $(G_L^{\mathrm{ES}}(A))^\perp$, its kernel is the original subspace
    $G_L^{\mathrm{ES}}(A)$.
\end{proof}

\begin{lemma}[Local-term kernel via cyclic restrictions]\label{lem:local_term_es_kernel_restrictions}
    \lean{MPSTensor.localTermES_eq_zero_iff_forall_cyclicRestrictES_mem_groundSpaceES}
    \lean{MPSTensor.cyclicRestrictES_mem_groundSpaceES_of_localTermES_eq_zero}
    \leanok
    \uses{lem:parent_interaction_es_kernel, rem:euclidean_local_projector_ingredients}
    For $L\le N$, a transported local term $h_{i,\mathrm{ES}}$ annihilates a vector
    $v$ iff every boundary-filled restriction of $v$ to the cyclic $L$-window
    starting at $i$ lies in $G_L^{\mathrm{ES}}(A)$.
\end{lemma}

\begin{proof}
    \leanok
    Restricting $h_{i,\mathrm{ES}}v$ to the same cyclic window gives
    $P_L^{\mathrm{ES}}(A)$ applied to the corresponding restriction of $v$.
    Lemma~\ref{lem:parent_interaction_es_kernel} identifies the vanishing of this
    local projector with membership in $G_L^{\mathrm{ES}}(A)$.
\end{proof}

\begin{theorem}[Adjacent local kernels realize the intersection property]
    \label{thm:adjacent_local_kernels_ground_space_es}
    \lean{MPSTensor.mem_groundSpaceES_succ_of_adjacent_localTermES_eq_zero}
    \lean{MPSTensor.adjacent_localTermES_eq_zero_of_mem_groundSpaceES_succ}
    \lean{MPSTensor.adjacent_localTermES_eq_zero_iff_mem_groundSpaceES_succ}
    \leanok
    \uses{thm:ground_space_intersection, lem:local_term_es_kernel_restrictions,
        lem:mem_ground_space_es}
    If $A$ is injective and $L\ge2$, then on an $(L+1)$-site chain the kernels of
    the two adjacent transported $L$-site local terms are exactly the transported
    $(L+1)$-site MPS ground space:
    \[
        h_{0,\mathrm{ES}}v=0\ \text{ and }\ h_{1,\mathrm{ES}}v=0
        \quad\Longleftrightarrow\quad
        v\in G_{L+1}^{\mathrm{ES}}(A).
    \]
\end{theorem}

\begin{proof}
    \leanok
    \uses{thm:ground_space_intersection, lem:local_term_es_kernel_restrictions,
        lem:mem_ground_space_es}
    Lemma~\ref{lem:local_term_es_kernel_restrictions} translates the two local
    kernel assumptions into the left and right $L$-site ground conditions. The
    open-chain intersection property, Theorem~\ref{thm:ground_space_intersection},
    then grows back the shared $(L-1)$-site overlap to the $(L+1)$-site ground space.
    The converse uses the forward restriction lemmas for ground-space vectors and
    the same kernel--restriction characterization.
\end{proof}

\begin{lemma}[Non-overlap positivity for transported ES local terms]
    \label{lem:transported_es_nonoverlap_positive}
    \lean{MPSTensor.CyclicWindowsDisjoint}
    \lean{MPSTensor.CyclicWindowsDisjoint.symm}
    \lean{MPSTensor.CyclicWindowsDisjoint.of_not_cyclicWindowsOverlap}
    \lean{MPSTensor.localTermES_commute_of_cyclic_windows_disjoint}
    \lean{MPSTensor.localTermES_re_inner_nonneg_of_cyclic_windows_disjoint}
    \leanok
    \uses{lem:transported_es_local_projections, rem:projection_geometry_rowsum_identities}
    Let $L\le N$ and let $h_{i,\mathrm{ES}},h_{j,\mathrm{ES}}$ be transported local
    terms whose cyclic $L$-windows are site-disjoint: no site of the periodic chain
    has cyclic offset $<L$ from both $i$ and $j$. Then the terms commute pointwise:
    for every vector $v$,
    \[
        h_{i,\mathrm{ES}}(h_{j,\mathrm{ES}}v)
        = h_{j,\mathrm{ES}}(h_{i,\mathrm{ES}}v),
    \]
    and their ordered cross term is non-negative:
    \[
        0\le \operatorname{Re}\langle h_{i,\mathrm{ES}}v,
        h_{j,\mathrm{ES}}v\rangle.
    \]
\end{lemma}

\begin{proof}
    \leanok
    \uses{lem:transported_es_local_projections, rem:projection_geometry_rowsum_identities}
    In coordinates, $h_{i,\mathrm{ES}}$ applies $P_L^{\mathrm{ES}}(A)$ only to the
    $i$-window variables and leaves the complementary variables fixed. Site
    disjointness gives two replacement identities: replacing the $i$-window does
    not change the extracted $j$-window, and the $i$- and $j$-window replacements
    commute. Therefore the two embedded local operators commute. The preceding
    projection-geometry identity for commuting symmetric projections gives the
    non-negative ordered cross term.
\end{proof}

\begin{lemma}[Quadratic form to norm bound]\label{lem:quadratic_form_to_norm_bound}
    \lean{FrustrationFree.spectralGap_of_martingale}
    \leanok
    Let $H$ be a positive semidefinite self-adjoint operator on a
    finite-dimensional complex Hilbert space, and suppose that, for all~$v$,
    \[
        \gamma\,\langle v, H v\rangle \le \langle H v, H v\rangle
    \]
    i.e.\ $H^2 \ge \gamma H$ as a quadratic form, for some $\gamma > 0$.
    Then for every $v$ in the orthogonal complement of $\ker H$,
    \[
        \gamma\,\|v\| \le \|H v\|.
    \]
    In particular, every positive eigenvalue of $H$ is at least $\gamma$.
    The positive-semidefiniteness hypothesis is essential: without it,
    $H = -\mathrm{Id}$ with $\gamma = 2$ satisfies the quadratic-form
    inequality vacuously but violates the norm bound.
\end{lemma}

\begin{proof}
    \leanok
    By the finite-dimensional spectral theorem, $H$ diagonalises in an
    orthonormal eigenbasis with real eigenvalues $\lambda_k$. For a unit
    eigenvector $\psi_k$ with eigenvalue $\lambda_k$ the hypothesis
    gives $\gamma\,\lambda_k \le \lambda_k^2$, whence $\lambda_k \le 0$
    or $\lambda_k \ge \gamma$. Since $H$ is positive semidefinite
    (its quadratic form is $\ge 0$), the negative case is excluded and
    every nonzero eigenvalue satisfies $\lambda_k \ge \gamma$. Writing
    $v = \sum_k c_k \psi_k$ and projecting onto $(\ker H)^{\perp}$
    keeps only terms with $\lambda_k \ge \gamma$, so
    $\|H v\|^2 = \sum_k |c_k|^2 \lambda_k^2
        \ge \gamma^2 \sum_k |c_k|^2 = \gamma^2 \|v\|^2$.
\end{proof}

\begin{lemma}[Parent-Hamiltonian quadratic-form reduction]
    \label{lem:parent_hamiltonian_es_quadratic_reduction}
    \lean{MPSTensor.parentHamiltonianES_norm_bound_of_quadratic_form}
    \lean{MPSTensor.parentHamiltonianES_gap_bound_of_quadratic_form}
    \leanok
    \uses{lem:transported_es_positive, thm:parent_ground_space_es_kernel,
        lem:quadratic_form_to_norm_bound}
    Let $H_{N,L}^{\mathrm{ES}}(A)$ be the transported parent Hamiltonian. If a
    constant $\gamma > 0$ satisfies the uniform quadratic-form estimate
    \begin{equation}\label{eq:ph_uniform_quadratic_form}
        \gamma\,\operatorname{Re}\langle H_{N,L}^{\mathrm{ES}}(A)v, v\rangle
        \le
        \operatorname{Re}\langle H_{N,L}^{\mathrm{ES}}(A)v,
        H_{N,L}^{\mathrm{ES}}(A)v\rangle
    \end{equation}
    for every $N \ge 2L$ and every vector $v$, then
    \[
        \gamma\,\|v\| \le \|H_{N,L}^{\mathrm{ES}}(A)v\|
    \]
    for every $v \in (\mathcal G_{N,L}^{\mathrm{ES}}(A))^\perp$. In particular,
    with $L>1$ and $\gamma = 1/(4L)$, the explicit gap-bound conclusion follows
    once this same quadratic-form estimate~(\ref{eq:ph_uniform_quadratic_form}) is supplied.
\end{lemma}

\begin{proof}
    \leanok
    The transported Hamiltonian is positive by
    Lemma~\ref{lem:transported_es_positive}, and its transported ground space is
    its kernel by Theorem~\ref{thm:parent_ground_space_es_kernel}. Therefore the
    abstract spectral-theorem conversion of
    Lemma~\ref{lem:quadratic_form_to_norm_bound} applies directly. The
    explicit constant is positive because $L>1$ implies $4L>0$.
\end{proof}

\begin{remark}[Projection-geometry identities for row-sum arguments]
    \label{rem:projection_geometry_rowsum_identities}
    \lean{LinearMap.IsSymmetricProjection.re_inner_apply_apply_self}
    \lean{LinearMap.IsSymmetricProjection.re_inner_nonneg}
    \lean{LinearMap.IsSymmetricProjection.re_inner_apply_apply_nonneg_of_commute}
    \lean{LinearMap.IsSymmetricProjection.re_inner_apply_apply_ge_neg_of_norm_apply_le}
    \lean{ProjectionGeometry.re_inner_sum_apply_left}
    \lean{ProjectionGeometry.re_inner_sum_apply_apply}
    \lean{ProjectionGeometry.sum_sum_eq_diag_add_offdiag}
    \lean{ProjectionGeometry.crossTerm_sum_bound_of_ordered_rowSum}
    \leanok
    The projection-geometry reduction uses the following elementary Hilbert-space
    identities: the diagonal identity
    $\operatorname{Re}\langle Pv,Pv\rangle=\operatorname{Re}\langle Pv,v\rangle$
    for symmetric projections; nonnegativity of this quadratic form; nonnegativity
    of $\operatorname{Re}\langle Pv,Qv\rangle$ for commuting symmetric projections
    $P,Q$; the Cauchy--Schwarz conversion from
    $\|P Q v\|\le c\|Pv\|$ to
    $\operatorname{Re}\langle Pv,Qv\rangle\ge -c\operatorname{Re}\langle Pv,v\rangle$;
    finite-sum expansions of
    $\operatorname{Re}\langle (\sum_i P_i)v,v\rangle$ and
    $\operatorname{Re}\langle (\sum_i P_i)v,(\sum_i P_i)v\rangle$; the
    diagonal/off-diagonal split of the ordered double sum; and the aggregate
    off-diagonal estimate obtained from row sums $\sum_{j\ne i}c_{ij}\le1$.
\end{remark}

\begin{lemma}[Ordered projection row-sum quadratic-form reduction]
    \label{lem:ordered_projection_rowsum_reduction}
    \lean{ProjectionGeometry.quadraticForm_sum_projections_of_ordered_rowSum}
    \leanok
    \uses{rem:projection_geometry_rowsum_identities}
    Let $P_i$ be a finite family of symmetric projections and let
    $H=\sum_i P_i$. If the ordered off-diagonal forms obey
    \[
        \operatorname{Re}\langle P_i v, P_j v\rangle
        \ge -(1-\gamma)c_{ij}\operatorname{Re}\langle P_i v,v\rangle
        \qquad (i\ne j),
    \]
    with row sums $\sum_{j\ne i} c_{ij}\le 1$ and $\gamma\le 1$, then
    $H^2\ge \gamma H$ as a quadratic form.
\end{lemma}

\begin{proof}
    \leanok
    Expand $\operatorname{Re}\langle H v,Hv\rangle$ as an ordered double sum.
    Symmetric idempotence identifies each diagonal term with
    $\operatorname{Re}\langle P_i v,v\rangle$, while the row-summable ordered
    cross-term bounds control the off-diagonal contribution by
    $-(1-\gamma)\sum_i\operatorname{Re}\langle P_i v,v\rangle$.
\end{proof}

\begin{lemma}[Finite-overlap projection row-sum reduction]
    \label{lem:finite_overlap_projection_rowsum_reduction}
    \lean{ProjectionGeometry.quadraticForm_sum_projections_of_finite_overlap}
    \leanok
    \uses{lem:ordered_projection_rowsum_reduction}
    Let $P_i$ be a finite family of symmetric projections and let $\gamma\le1$.
    Suppose an interaction predicate $i\sim j$ has at most $m$ off-diagonal
    neighbours in each row for some positive integer $m$. If noninteracting
    ordered cross terms are non-negative and, on each interacting pair,
    \[
        \operatorname{Re}\langle P_i v,P_jv\rangle
        \ge -(1-\gamma)\frac{1}{m}\operatorname{Re}\langle P_i v,v\rangle,
    \]
    then $H=\sum_i P_i$ satisfies $H^2\ge \gamma H$ as a quadratic form.
\end{lemma}

\begin{proof}
    \leanok
    Choose $c_{ij}=1/m$ on interacting pairs and $c_{ij}=0$ otherwise. The row
    cardinality bound gives $\sum_{j\ne i}c_{ij}\le1$; noninteracting
    nonnegativity supplies the zero-coefficient cross-term estimates. Apply
    Lemma~\ref{lem:ordered_projection_rowsum_reduction}.
\end{proof}

\begin{lemma}[Finite-overlap projection norm-compression reduction]
    \label{lem:finite_overlap_projection_norm_reduction}
    \lean{ProjectionGeometry.quadraticForm_sum_projections_of_finite_overlap_norm_bound}
    \lean{ProjectionGeometry.quadraticForm_sum_projections_of_finite_overlap_norm_bound_of_le}
    \leanok
    \uses{lem:finite_overlap_projection_rowsum_reduction,
        rem:projection_geometry_rowsum_identities}
    Let $P_i$ be a finite family of symmetric projections and let $\gamma\le1$.
    Suppose an interaction predicate has at most $m>0$ off-diagonal neighbours
    in each row, noninteracting ordered cross terms are non-negative, and
    interacting pairs satisfy the norm-compression estimate
    \begin{equation}\label{eq:ph_finite_overlap_norm_compression}
        \|P_iP_jv\|\le (1-\gamma)\frac{1}{m}\|P_iv\|.
    \end{equation}
    Then $H=\sum_iP_i$ satisfies $H^2\ge\gamma H$ as a quadratic form. The same
    conclusion holds from a separate compression coefficient $\eta$ whenever
    $\eta\le (1-\gamma)/m$.
\end{lemma}

\begin{proof}
    \leanok
    \uses{lem:finite_overlap_projection_rowsum_reduction,
        rem:projection_geometry_rowsum_identities}
    Cauchy--Schwarz and symmetric idempotence turn the norm-compression estimate
    into
    $\operatorname{Re}\langle P_i v,P_jv\rangle
        \ge -(1-\gamma)m^{-1}\operatorname{Re}\langle P_iv,v\rangle$.
    If a separate coefficient $\eta$ is used, the assumption
    $\eta\le(1-\gamma)/m$ first weakens that estimate to
    (\ref{eq:ph_finite_overlap_norm_compression}). The finite-overlap row-sum
    reduction then applies.
\end{proof}

\begin{lemma}[Parent-Hamiltonian ordered row-sum quadratic-form reduction]
    \label{lem:parent_hamiltonian_ordered_rowsum_quadratic}
    \lean{MPSTensor.parentHamiltonianES_quadratic_form_of_ordered_local_term_bounds}
    \leanok
    \uses{lem:ordered_projection_rowsum_reduction, def:parent_hamiltonian_es,
        rem:euclidean_local_projector_ingredients, lem:transported_es_local_projections}
    For a fixed chain length $N$, suppose the transported local terms
    $h_{i,\mathrm{ES}}$ satisfy the ordered row-sum estimate
    \begin{equation}\label{eq:ph_ordered_rowsum_estimate}
        \operatorname{Re}\langle h_{i,\mathrm{ES}} v,h_{j,\mathrm{ES}}v\rangle
        \ge -(1-\gamma)c_{ij}\operatorname{Re}\langle h_{i,\mathrm{ES}}v,v\rangle
        \qquad (i\ne j),
    \end{equation}
    with $\sum_{j\ne i}c_{ij}\le1$ and $\gamma\le1$. Then the transported
    parent Hamiltonian satisfies
    \[
        \gamma\operatorname{Re}\langle H_{N,L}^{\mathrm{ES}}v,v\rangle
        \le
        \operatorname{Re}\langle H_{N,L}^{\mathrm{ES}}v,
        H_{N,L}^{\mathrm{ES}}v\rangle
    \]
    for every vector $v$.
\end{lemma}

\begin{proof}
    \leanok
    \uses{lem:transported_es_local_projections, lem:ordered_projection_rowsum_reduction}
    Lemma~\ref{lem:transported_es_local_projections} supplies the symmetric-projection
    hypothesis for the family $P_i=h_{i,\mathrm{ES}}$. Apply
    Lemma~\ref{lem:ordered_projection_rowsum_reduction} to this family and use
    $H_{N,L}^{\mathrm{ES}}=\sum_i h_{i,\mathrm{ES}}$.
\end{proof}

\begin{lemma}[Parent-Hamiltonian gap from ordered local row bounds]
    \label{lem:parent_hamiltonian_ordered_rowsum_gap}
    \lean{MPSTensor.parentHamiltonianES_gap_bound_of_ordered_local_term_bounds}
    \leanok
    \uses{lem:parent_hamiltonian_ordered_rowsum_quadratic,
        lem:parent_hamiltonian_es_quadratic_reduction, lem:transported_es_local_projections}
    Assume uniformly for every $N\ge2L$ that the transported local terms satisfy
    the ordered row-sum estimate~(\ref{eq:ph_ordered_rowsum_estimate}) with $\gamma=1/(4L)$. If $L>1$, then
    $1/(4L)>0$ and, for every
    $v\in(\mathcal G_{N,L}^{\mathrm{ES}}(A))^\perp$,
    \[
        \frac{1}{4L}\|v\|\le
        \|H_{N,L}^{\mathrm{ES}}(A)v\|.
    \]
\end{lemma}

\begin{proof}
    \leanok
    \uses{lem:parent_hamiltonian_ordered_rowsum_quadratic,
        lem:parent_hamiltonian_es_quadratic_reduction}
    The fixed-chain row-sum reduction supplies the quadratic-form estimate with
    $\gamma=1/(4L)$ for every admissible $N$. The positivity and kernel
    identification of the transported Hamiltonian then allow
    Lemma~\ref{lem:parent_hamiltonian_es_quadratic_reduction} to convert this
    quadratic-form estimate into the norm lower bound on the orthogonal
    complement of the ground space.
\end{proof}

\begin{lemma}[Parent-Hamiltonian finite-overlap Friedrichs quadratic reduction]
    \label{lem:parent_hamiltonian_finite_overlap_friedrichs_quadratic}
    \lean{MPSTensor.parentHamiltonianES_quadratic_form_of_finite_overlap_friedrichs}
    \leanok
    \uses{lem:finite_overlap_projection_rowsum_reduction, def:parent_hamiltonian_es}
    For a fixed chain length $N$, let $m$ be a positive integer and let
    $\gamma\le1$. Suppose an overlap predicate on local windows has at most $m$
    overlapping off-diagonal windows in each row. If non-overlapping local terms
    have non-negative ordered cross terms, overlapping local terms satisfy the
    finite-overlap Friedrichs estimate
    \[
        \operatorname{Re}\langle h_{i,\mathrm{ES}}v,h_{j,\mathrm{ES}}v\rangle
        \ge -(1-\gamma)\frac{1}{m}
        \operatorname{Re}\langle h_{i,\mathrm{ES}}v,v\rangle,
    \]
    and the transported local terms are symmetric projections, then the
    transported Hamiltonian satisfies $H_{N,L}^{\mathrm{ES}}{}^2\ge
        \gamma H_{N,L}^{\mathrm{ES}}$ as a quadratic form.
\end{lemma}

\begin{proof}
    \leanok
    Apply the finite-overlap projection reduction with
    $P_i=h_{i,\mathrm{ES}}$.
\end{proof}

\begin{lemma}[Finite-overlap norm-compression reduction]
    \label{lem:parent_hamiltonian_finite_overlap_norm_quadratic}
    \lean{MPSTensor.parentHamiltonianES_quadratic_form_of_finite_overlap_norm_bound}
    \lean{MPSTensor.parentHamiltonianES_quadratic_form_of_finite_overlap_norm_bound_of_le}
    \leanok
    \uses{lem:finite_overlap_projection_norm_reduction, def:parent_hamiltonian_es,
        lem:transported_es_local_projections}
    For a fixed chain length $N$, let $m>0$ and $\gamma\le1$. Suppose an overlap
    predicate on local windows has at most $m$ overlapping off-diagonal windows in
    each row, non-overlapping local terms have non-negative ordered cross terms,
    and overlapping local terms satisfy
    \[
        \|h_{i,\mathrm{ES}}(h_{j,\mathrm{ES}}v)\|
        \le (1-\gamma)\frac{1}{m}\|h_{i,\mathrm{ES}}v\|.
    \]
    Then $H_{N,L}^{\mathrm{ES}}{}^2\ge \gamma H_{N,L}^{\mathrm{ES}}$ as a
    quadratic form. The same conclusion holds from a separate coefficient $\eta$
    satisfying $\eta\le(1-\gamma)/m$.
\end{lemma}

\begin{proof}
    \leanok
    \uses{lem:finite_overlap_projection_norm_reduction,
        lem:transported_es_local_projections}
    Apply the finite-overlap norm-compression projection reduction to the family
    $P_i=h_{i,\mathrm{ES}}$. If a separate coefficient $\eta$ is used, first weaken
    the compression estimate using $\eta\le(1-\gamma)/m$.
\end{proof}

\begin{definition}[Cyclic-window support]
    \label{def:cyclic_window_support}
    \lean{MPSTensor.cyclicWindowSupport}
    \leanok
    \uses{rem:cyclic_window_operators}
    For a length-$L$ cyclic window on $\Fin{N}$ starting at $i$, its support is
    \[
        S_i=\{i,i+1,\ldots,i+L-1\}\pmod N,
    \]
    with repeated visits counted once.
\end{definition}

\begin{lemma}[Cyclic-window support offsets]
    \label{lem:cyclic_window_support_offsets}
    \lean{MPSTensor.mem_cyclicWindowSupport_iff}
    \leanok
    \uses{def:cyclic_window_support}
    If $L\le N$, then a site $k$ belongs to the cyclic support $S_i$ exactly when
    its cyclic offset from $i$ is below $L$:
    \[
        k\in S_i
        \quad\Longleftrightarrow\quad
        (k-i)\bmod N < L.
    \]
\end{lemma}

\begin{proof}
    \leanok
    Write $k=i+r\pmod N$ for a representative offset $r$. Membership in
    $S_i$ gives an offset $r<L$, and conversely an offset below $L$ gives the
    corresponding point of the support.
\end{proof}

\begin{definition}[Cyclic-window overlap]
    \label{def:cyclic_windows_overlap}
    \lean{MPSTensor.cyclicWindowsOverlap}
    \leanok
    \uses{def:cyclic_window_support}
    Two length-$L$ cyclic windows overlap, written $i\sim_L j$, when their
    supports meet: $S_i\cap S_j\ne\varnothing$.
\end{definition}

\begin{lemma}[Concrete non-overlap positivity for cyclic supports]
    \label{lem:cyclic_window_nonoverlap_positive}
    \lean{MPSTensor.CyclicWindowsDisjoint.of_not_cyclicWindowsOverlap}
    \lean{MPSTensor.localTermES_re_inner_nonneg_of_not_cyclicWindowsOverlap}
    \leanok
    \uses{lem:cyclic_window_support_offsets, def:cyclic_windows_overlap,
        lem:transported_es_nonoverlap_positive}
    If $L\le N$ and the cyclic supports $S_i$ and $S_j$ do not meet, then the
    transported local ES terms have non-negative ordered cross term:
    \[
        0\le\operatorname{Re}\langle h_{i,\mathrm{ES}}v,
        h_{j,\mathrm{ES}}v\rangle.
    \]
\end{lemma}

\begin{proof}
    \leanok
    \uses{lem:cyclic_window_support_offsets, lem:transported_es_nonoverlap_positive}
    The support-offset characterization converts failure of $S_i\cap S_j\ne\varnothing$
    into site-disjointness of the two cyclic windows. The disjoint-window
    positivity lemma then applies.
\end{proof}

\begin{lemma}[Cyclic-window overlap row cardinality]
    \label{lem:cyclic_window_overlap_cardinality}
    \lean{MPSTensor.cyclicWindowsOverlap_card_le}
    \leanok
    \uses{def:cyclic_window_support, def:cyclic_windows_overlap}
    If $2L\le N$ and $L>1$, then for every $i\in\Fin{N}$,
    \[
        \#\{j\ne i: i\sim_L j\}\le 2(L-1).
    \]
\end{lemma}

\begin{proof}
    \leanok
    \uses{def:cyclic_window_support, def:cyclic_windows_overlap}
    If $S_i$ and $S_j$ meet, choose offsets $a,b<L$ with
    $i+a\equiv j+b\pmod N$. Since $2L\le N$, the offset from $i$ to $j$ is
    either the clockwise distance $a-b\in\{1,\ldots,L-1\}$ or the
    counterclockwise distance $b-a\in\{1,\ldots,L-1\}$; the zero distance is
    excluded by $j\ne i$. Thus every overlapping off-diagonal start lies among
    the $L-1$ clockwise starts or the $L-1$ counterclockwise starts.
\end{proof}

\begin{lemma}[Parent-Hamiltonian gap from finite-overlap Friedrichs data]
    \label{lem:parent_hamiltonian_finite_overlap_friedrichs_gap}
    \lean{MPSTensor.parentHamiltonianES_gap_bound_of_finite_overlap_friedrichs}
    \leanok
    \uses{lem:parent_hamiltonian_finite_overlap_friedrichs_quadratic,
        lem:parent_hamiltonian_es_quadratic_reduction}
    Fix $L>1$ and set $m=2(L-1)$. Suppose uniformly for every $N\ge2L$ that
    the hypotheses of the fixed-chain finite-overlap Friedrichs reduction hold
    with $\gamma=1/(4L)$ and this value of $m$. Then $1/(4L)>0$, and the
    explicit norm lower bound with constant $1/(4L)$ follows for vectors in
    $(\mathcal G_{N,L}^{\mathrm{ES}}(A))^\perp$.
\end{lemma}

\begin{proof}
    \leanok
    The fixed-chain finite-overlap reduction gives the required quadratic-form
    estimate for every admissible $N$. Lemma~\ref{lem:parent_hamiltonian_es_quadratic_reduction}
    converts that estimate to the norm lower bound.
\end{proof}

\begin{lemma}[Parent-Hamiltonian gap from cyclic-window Friedrichs data]
    \label{lem:parent_hamiltonian_cyclic_window_friedrichs_gap}
    \lean{MPSTensor.parentHamiltonianES_gap_bound_of_cyclic_window_friedrichs}
    \leanok
    \uses{lem:parent_hamiltonian_finite_overlap_friedrichs_gap,
        lem:cyclic_window_overlap_cardinality, lem:transported_es_local_projections,
        lem:cyclic_window_nonoverlap_positive}
    Fix $L>1$ and let $i\sim_L j$ mean that the cyclic supports of the
    length-$L$ windows at $i$ and $j$ meet. Suppose uniformly for every
    $N\ge2L$ that overlapping off-diagonal terms satisfy
    \begin{equation}\label{eq:ph_cyclic_window_overlap_estimate}
        \operatorname{Re}\langle h_{i,\mathrm{ES}}v,h_{j,\mathrm{ES}}v\rangle
        \ge -(1-\tfrac{1}{4L})\frac{1}{2(L-1)}
        \operatorname{Re}\langle h_{i,\mathrm{ES}}v,v\rangle.
    \end{equation}
    Then $1/(4L)>0$, and the explicit norm lower bound with constant $1/(4L)$
    follows for vectors in $(\mathcal G_{N,L}^{\mathrm{ES}}(A))^\perp$.
\end{lemma}

\begin{proof}
    \leanok
    \uses{lem:parent_hamiltonian_finite_overlap_friedrichs_gap,
        lem:cyclic_window_overlap_cardinality, lem:transported_es_local_projections,
        lem:cyclic_window_nonoverlap_positive}
    Lemma~\ref{lem:cyclic_window_overlap_cardinality} gives the row-cardinality
    bound $m=2(L-1)$ for the concrete cyclic-window overlap relation, and
    Lemma~\ref{lem:transported_es_local_projections} supplies the symmetric
    projection structure of each transported local term. If two cyclic supports
    do not meet, the corresponding windows are site-disjoint, so
    Lemma~\ref{lem:cyclic_window_nonoverlap_positive} gives the required
    non-negative ordered cross term. Apply
    Lemma~\ref{lem:parent_hamiltonian_finite_overlap_friedrichs_gap} with these
    facts and the overlapping-window estimate~(\ref{eq:ph_cyclic_window_overlap_estimate}).
\end{proof}

\begin{lemma}[Parent-Hamiltonian gap from overlapping cyclic Friedrichs data]
    \label{lem:parent_hamiltonian_cyclic_overlap_friedrichs_gap}
    \lean{MPSTensor.parentHamiltonianES_gap_bound_of_cyclic_window_overlap_friedrichs}
    \leanok
    \uses{lem:parent_hamiltonian_cyclic_window_friedrichs_gap}
    Fix $L>1$ and let $i\sim_L j$ mean that the cyclic supports of the
    length-$L$ windows at $i$ and $j$ meet. Suppose uniformly for every
    $N\ge2L$ that overlapping off-diagonal terms satisfy
    \begin{equation}\label{eq:ph_cyclic_overlap_friedrichs}
        \operatorname{Re}\langle h_{i,\mathrm{ES}}v,h_{j,\mathrm{ES}}v\rangle
        \ge -(1-\tfrac{1}{4L})\frac{1}{2(L-1)}
        \operatorname{Re}\langle h_{i,\mathrm{ES}}v,v\rangle.
    \end{equation}
    Then the explicit norm lower bound with constant $1/(4L)$ follows for
    vectors in $(\mathcal G_{N,L}^{\mathrm{ES}}(A))^\perp$.
\end{lemma}

\begin{proof}
    \leanok
    \uses{lem:parent_hamiltonian_cyclic_window_friedrichs_gap}
    Apply Lemma~\ref{lem:parent_hamiltonian_cyclic_window_friedrichs_gap}.
    Its non-overlap positivity and finite-row hypotheses are already internal to
    that cyclic-window reduction; the overlapping estimate~(\ref{eq:ph_cyclic_overlap_friedrichs}) supplies its
    remaining ordered cross-term condition.
\end{proof}

\begin{lemma}[Parent-Hamiltonian gap from overlapping cyclic norm-compression data]
    \label{lem:parent_hamiltonian_cyclic_overlap_norm_gap}
    \lean{MPSTensor.parentHamiltonianES_gap_bound_of_cyclic_window_overlap_norm_bound}
    \leanok
    \uses{lem:parent_hamiltonian_cyclic_overlap_friedrichs_gap,
        rem:projection_geometry_rowsum_identities}
    Fix $L>1$. Suppose uniformly for every $N\ge2L$ and every overlapping
    off-diagonal pair that
    \[
        \|h_{i,\mathrm{ES}}(h_{j,\mathrm{ES}}v)\|
        \le (1-\tfrac{1}{4L})\frac{1}{2(L-1)}
        \|h_{i,\mathrm{ES}}v\|.
    \]
    Then the explicit norm lower bound with constant $1/(4L)$ follows for
    vectors in $(\mathcal G_{N,L}^{\mathrm{ES}}(A))^\perp$.
\end{lemma}

\begin{proof}
    \leanok
    \uses{lem:parent_hamiltonian_cyclic_overlap_friedrichs_gap,
        rem:projection_geometry_rowsum_identities}
    The norm-compression condition is converted to the ordered Friedrichs lower
    bound by the projection-geometry Cauchy--Schwarz lemma. The preceding
    overlapping-cyclic Friedrichs reduction then gives the norm lower bound.
\end{proof}

\begin{lemma}[Parent-Hamiltonian gap from a cyclic overlap compression constant]
    \label{lem:parent_hamiltonian_cyclic_overlap_norm_constant_gap}
    \lean{MPSTensor.parentHamiltonianES_gap_bound_of_cyclic_window_overlap_norm_bound_of_le}
    \lean{MPSTensor.parentHamiltonianES_gap_bound_of_cyclic_window_overlap_norm_bound_of_lt}
    \leanok
    \uses{lem:parent_hamiltonian_finite_overlap_norm_quadratic,
        lem:cyclic_window_overlap_cardinality, lem:cyclic_window_nonoverlap_positive,
        lem:parent_hamiltonian_es_quadratic_reduction}
    Fix $L>1$, and let $m=2(L-1)$. Suppose uniformly for every $N\ge2L$ and every
    overlapping off-diagonal pair that
    \[
        \|h_{i,\mathrm{ES}}(h_{j,\mathrm{ES}}v)\|
        \le \eta\,\|h_{i,\mathrm{ES}}v\|.
    \]
    If $\gamma>0$, $\gamma\le1$, and $\eta\le(1-\gamma)/m$, then $0<\gamma$
    and, for every admissible $N$ and every
    $v\in(\mathcal G_{N,L}^{\mathrm{ES}}(A))^\perp$,
    \[
        \gamma\|v\|\le \|H_{N,L}^{\mathrm{ES}}(A)v\|.
    \]
    The assertion $0<\gamma$ repeats the hypothesis, in the same shape as the
    strict statement below. In particular, if
    $\eta\ge0$ and $\eta m<1$, then $0<1-\eta m$ and, for every admissible $N$
    and every $v\in(\mathcal G_{N,L}^{\mathrm{ES}}(A))^\perp$,
    \[
        (1-\eta m)\|v\|\le \|H_{N,L}^{\mathrm{ES}}(A)v\|.
    \]
\end{lemma}

\begin{proof}
    \leanok
    \uses{lem:parent_hamiltonian_finite_overlap_norm_quadratic,
        lem:cyclic_window_overlap_cardinality, lem:cyclic_window_nonoverlap_positive,
        lem:parent_hamiltonian_es_quadratic_reduction}
    The cyclic-window cardinality lemma gives at most $m=2(L-1)$ overlapping
    off-diagonal neighbours in each row, while the non-overlap positivity lemma
    handles disjoint supports. The finite-overlap norm-compression reduction gives
    the quadratic-form estimate with the chosen $\gamma$, and the positive
    quadratic-form criterion converts it to the norm lower bound. The strict form
    is the same statement with $\gamma=1-\eta m$.
\end{proof}

\begin{theorem}[Martingale criterion for spectral gap]\label{thm:martingale_criterion}
    \notready
    \uses{def:is_frustration_free, lem:quadratic_form_to_norm_bound}
    Consider a frustration-free Hamiltonian
    \[
        H = \sum_{i=1}^n h_i
    \]
    with projector interaction terms ($h_i^2 = h_i$).
    Suppose that for every unordered pair $\{i,j\}$ of overlapping terms there
    exist constants $c_{\{i,j\}} \ge 0$ and $\gamma > 0$ satisfying
    \begin{equation}\label{eq:ph_martingale_condition}
        h_i h_j + h_j h_i
        \ge -c_{\{i,j\}}(1-\gamma)(h_i + h_j),
    \end{equation}
    and suppose further that the coefficients attached to the terms overlapping
    $h_i$ satisfy
    \[
        \sum_{\substack{j\ne i\\ \{i,j\}\text{ overlapping}}}c_{\{i,j\}}\le1
    \]
    for every $i$.
    Then the smallest positive eigenvalue of $H$ satisfies
    $\lambda_{\min}^+(H) \ge \gamma$.
\end{theorem}

\begin{proof}
    \notready
    The relevant geometry is a family of translated length-$L$ projector
    windows on the ring; two terms $h_i$ and $h_j$ interact only when their
    supports overlap, equivalently when the cyclic distance $d_N(i,j)$ is
    smaller than $L$.
    From $h_i^2 = h_i$,
    \[
        H^2 = \sum_i h_i + \sum_{\substack{i < j \\ \text{overlapping}}}
        (h_i h_j + h_j h_i) + \sum_{\substack{i < j \\ \text{non-overlapping}}}
        (h_i h_j + h_j h_i).
    \]
    Non-overlapping projectors commute ($h_i h_j = h_j h_i$), so their
    product $h_i h_j$ is itself a projector and hence PSD; thus the last
    sum is $\ge 0$.
    For the overlapping sum, the martingale condition~(\ref{eq:ph_martingale_condition}) gives
    $h_i h_j + h_j h_i \ge -c_{\{i,j\}}(1-\gamma)(h_i + h_j)$.
    Summing over unordered pairs $\{i,j\}$:
    \[
        \sum_{\{i,j\}, \text{overlap}} (h_i h_j + h_j h_i)
        \ge -(1-\gamma) \sum_{\{i,j\}, \text{overlap}}
        c_{\{i,j\}}(h_i + h_j).
    \]
    Collecting terms for each $h_i$:
    \[
        \sum_{\{i,j\}, \text{overlap}} c_{\{i,j\}}(h_i + h_j)
        = \sum_i \Bigl(\sum_{\substack{j \ne i \\ \text{overlap}}}
        c_{\{i,j\}}\Bigr) h_i
        \le H,
    \]
    where the last inequality uses the row-sum hypothesis
    $\sum_{\substack{j\ne i\\ \{i,j\}\text{ overlapping}}}c_{\{i,j\}}\le1$.
    Combining, $H^2 \ge H - (1-\gamma)H = \gamma H$.
    For any eigenvector with $H\psi = \lambda\psi$ and $\lambda > 0$,
    this gives $\lambda^2 \ge \gamma\lambda$, hence
    $\lambda \ge \gamma$~\cite{Kastoryano2018Martingale}.
\end{proof}

\begin{lemma}[Martingale condition for MPS]\label{lem:martingale_mps}
    \notready
    \uses{thm:martingale_criterion, def:parent_interaction,
        def:local_term_parent, thm:ground_space_intersection, def:injective}
    Let $A$ be injective and fix an interaction range $L \ge 2$ as in
    Definition~\ref{def:parent_interaction}. Let $h_L(A) := \Pi_{G_L(A)^\perp}$
    be the $L$-site parent interaction, and for each chain length and site $i$
    let $h_i$ denote its translate as in Definition~\ref{def:local_term_parent}.
    Then for every pair of overlapping terms $h_i$ and $h_j$ with
    $0 < d_N(i,j) < L$, where $d_N(i,j) := \min(|i-j|,\, N-|i-j|)$ is the
    cyclic distance on the ring (i.e.\ whose $L$-site supports share at
    least one site under periodic boundary conditions),
    the martingale condition of Theorem~\ref{thm:martingale_criterion} holds with
    constants $c_{\{i,j\}}$ and $\gamma > 0$ depending only on $A$ (not on $N$).
    Cirac--Perez-Garcia--Schuch--Verstraete~\cite[Section~IV.C]{Cirac2021Matrix}
    obtain this uniformity by the martingale method, invoking finite-chain-state
    estimates to control the constants.  For the explicit cyclic-window
    constants displayed here, the remaining quantitative step is the
    overlapping-window Friedrichs-angle bound in
    Theorem~\ref{thm:friedrichs_gap_bound}; the preceding lemmas reduce the
    martingale inequality and the spectral gap to that estimate.
    (When $L = 1$ the interaction terms have disjoint supports, so there are no
    overlapping pairs and the spectral gap follows directly without the
    martingale argument.)
\end{lemma}

\begin{proof}
    \uses{thm:ground_space_intersection, def:injective}
    \notready
    Two terms $h_i, h_j$ with $0 < d_N(i,j) < L$ have supports overlapping
    in $L - d_N(i,j)$ sites (where $d_N$ is the cyclic distance on
    the $N$-site ring). Since $A$ is injective, the intersection property
    (Theorem~\ref{thm:ground_space_intersection}) gives the qualitative identity
    $\ker(h_i) \cap \ker(h_j) = G_{[i,j+L-1]}(A)$.
    Because the local Hilbert space is finite-dimensional, the subspaces
    $\ker(h_i)$ and $\ker(h_j)$ live in a fixed finite-dimensional space.
    The intersection identity shows that $\ker(h_i) \cap \ker(h_j)$ is
    exactly $G_{[i,j+L-1]}(A)$; in particular there is no non-trivial
    vector in one kernel that is orthogonal to the other yet lies in both.
    Hence the Friedrichs angle $\theta_{d_N(i,j)}$ between $\ker(h_i)$ and
    $\ker(h_j)$ is strictly positive (for the fixed injective tensor
    $A$). Because the angle depends only on the overlap pattern
    $d_N(i,j) \in \{1,\ldots,L-1\}$ and not on the absolute position,
    the number of distinct angles is at most $L-1$.
    Let $\alpha := \max_{1 \le s < L} \cos\theta_s < 1$ be the
    worst-case cosine of the Friedrichs angles.

    \emph{Row-sum bound.}
    Since $L \ge 2$, each site $i$ overlaps with at most $2(L-1)$
    neighbours (those $j$ with $0 < d_N(i,j) < L$).
    Set $c_{\{i,j\}} := \tfrac{1}{2(L-1)}$ for every unordered overlapping
    pair and $c_{\{i,j\}} := 0$ otherwise. These constants depend only on the
    cyclic distance $d_N(i,j)$.
    Then
    $\sum_{\substack{j\ne i\\ \{i,j\}\text{ overlapping}}}c_{\{i,j\}}
    \le 2(L-1)\cdot\tfrac{1}{2(L-1)} = 1$,
    satisfying the row-sum hypothesis of
    Theorem~\ref{thm:martingale_criterion}.

    \emph{Explicit cyclic-window constant.}
    With the row coefficient above and
    \[
        \gamma_L=\frac{1}{4L},
    \]
    the martingale condition is the fixed cyclic-window estimate
    \[
        h_i h_j + h_j h_i
        \ge
        -\left(1-\frac{1}{4L}\right)\frac{1}{2(L-1)}(h_i+h_j)
    \]
    for every overlapping off-diagonal pair. A sufficient form, used below, is
    the one-sided norm-compression estimate
    \[
        \|h_i h_j v\|
        \le
        \left(1-\frac{1}{4L}\right)\frac{1}{2(L-1)}\|h_i v\|.
    \]
    The qualitative Friedrichs-angle conclusion $\alpha<1$ does not by itself
    imply this numerical bound; the missing step is the uniform
    overlapping-window estimate with the displayed coefficient. Thus the
    remaining comparison is whether the principal-angle estimates cited in
    \cite{Fannes1992Finitely, Nachtergaele1996Spectral,
    Kastoryano2018Martingale} supply this cyclic-window inequality, with
    constants independent of $N$, under the convention used here.
    % Paper-gap note: docs/paper-gaps/cpgsv21_martingale_overlap.tex.
\end{proof}

\begin{theorem}[Explicit Friedrichs-angle gap bound]
    \label{thm:friedrichs_gap_bound}
    \uses{def:parent_ground_space_es, def:parent_hamiltonian_es,
        thm:parent_ground_space_es_kernel,
        rem:euclidean_local_projector_ingredients,
        lem:euclidean_local_projector_positive,
        lem:local_term_es_average,
        lem:transported_es_local_projections,
        lem:transported_es_nonoverlap_positive,
        lem:parent_hamiltonian_es_quadratic_reduction,
        lem:parent_hamiltonian_ordered_rowsum_gap,
        lem:parent_hamiltonian_finite_overlap_friedrichs_quadratic,
        lem:parent_hamiltonian_finite_overlap_friedrichs_gap,
        lem:cyclic_window_overlap_cardinality,
        lem:cyclic_window_nonoverlap_positive,
        lem:parent_hamiltonian_cyclic_window_friedrichs_gap,
        lem:parent_hamiltonian_cyclic_overlap_friedrichs_gap,
        lem:parent_hamiltonian_cyclic_overlap_norm_gap}
    For an injective tensor $A$ and a range $L > 1$, the theorem asserts
    the analytic estimate with the explicit constant
    \[
        \gamma_{L} := \frac{1}{4L} > 0.
    \]
    Concretely, it asserts that if $N \ge 2L$ and
    $v \in (\mathcal G_{N,L}^{\mathrm{ES}}(A))^{\perp}$, then
    \[
        \gamma_{L} \|v\| \le \|H_{N,L}^{\mathrm{ES}}(A) v\|.
    \]
\end{theorem}

\begin{proof}
    \notready
    The local averaging, positivity, and symmetric-projection statements are now
    established in Lemmas~\ref{lem:local_term_es_average},
    \ref{lem:transported_es_positive}, and~\ref{lem:transported_es_local_projections}.
    Lemma~\ref{lem:parent_hamiltonian_es_quadratic_reduction} reduces the present
    theorem to the remaining MPS-specific task: deriving the uniform
    quadratic-form estimate. Lemma~\ref{lem:parent_hamiltonian_ordered_rowsum_gap}
    establishes the finite-sum algebra from ordered cross-term row bounds, and
    Lemmas~\ref{lem:parent_hamiltonian_finite_overlap_friedrichs_quadratic}
    and~\ref{lem:parent_hamiltonian_finite_overlap_friedrichs_gap} state the
    common finite-overlap specialization with $m=2(L-1)$.
    Lemma~\ref{lem:transported_es_nonoverlap_positive} supplies the non-overlap
    positivity statement for site-disjoint windows, and
    Lemma~\ref{lem:cyclic_window_nonoverlap_positive} translates this to the
    concrete support-disjoint cyclic windows.
    Lemma~\ref{lem:cyclic_window_overlap_cardinality} gives the finite-overlap
    row-cardinality bound.
    Lemma~\ref{lem:parent_hamiltonian_cyclic_window_friedrichs_gap} combines
    these ingredients with the finite-overlap reduction, so the remaining
    analytic task is the Friedrichs-angle estimate for overlapping cyclic windows.
    Lemma~\ref{lem:parent_hamiltonian_cyclic_overlap_friedrichs_gap} states the
    same overlap-only formulation, and
    Lemma~\ref{lem:parent_hamiltonian_cyclic_overlap_norm_gap} states a sufficient
    norm-compression formulation.
\end{proof}

\begin{theorem}[Spectral gap for MPS parent Hamiltonians]\label{thm:parent_hamiltonian_gapped}
    \uses{thm:friedrichs_gap_bound, lem:quadratic_form_to_norm_bound}
    For an injective tensor $A$ and a range $L > 1$, there exists $\gamma > 0$ such that
    for every $N \ge 2L$ and every vector
    $v \in (\mathcal G_{N,L}^{\mathrm{ES}}(A))^{\perp}$ one has
    \[
        \gamma \|v\| \le \|H_{N,L}^{\mathrm{ES}}(A) v\|.
    \]
    By Lemma~\ref{lem:quadratic_form_to_norm_bound}, this implies the usual uniform lower
    bound on every nonzero eigenvalue of the parent Hamiltonian.
\end{theorem}

\begin{proof}
    \notready
    This follows immediately from Theorem~\ref{thm:friedrichs_gap_bound} by taking as
    $\gamma$ the explicit constant produced there, which yields the required existential
    spectral-gap constant.
\end{proof}

\section{Ground space of non-injective MPS}\label{sec:degenerate_ground_space}

For an MPS with more than one block in canonical form, after blocking to
block-injective canonical form, the physical--virtual correspondence is
restricted to block-diagonal boundary conditions.  In the notation of
\cite[lines~2113--2117]{Cirac2021Matrix}, set
\[
    \widetilde A^\omega:=(A_1^\omega,\ldots,A_g^\omega)
    \in\prod_{j=1}^g \MN{D_j}.
\]
Then block injectivity is the statement
\[
    \forall X=(X_j)_{j=1}^g\in\bigoplus_{j=1}^g \MN{D_j},
    \quad
    \exists (c_\omega(X))_\omega,
    \quad
    X=\sum_\omega c_\omega(X)\widetilde A^\omega.
\]
Equivalently, with the same coefficient family,
\[
    \forall j,\qquad
    X_j=\sum_\omega c_\omega(X)A_j^\omega .
\]
The parent-Hamiltonian theorem in \cite[lines~2121--2128]{Cirac2021Matrix} is
\[
    N\ge L_0+1
    \quad\Longrightarrow\quad
    \ker H_N=\spn\{V^{(N)}(A_j):j=1,\ldots,g\}.
\]
The statements below record the corresponding cyclic-chain ground-space
formulation; the passage to $\ker H_N$ is kept separate.

\begin{definition}[Chain ground space for the block-diagonal tensor]
    \label{def:parent_hamiltonian_ground_space_bnt}
    \notready
    \uses{def:chain_ground_space}
    For a basis of normal tensors $A=(A_j)_{j=1}^g$ and nonzero weights
    $\mu_j$, let
    \[
        B := \bigoplus_{j=1}^g \mu_j A_j .
    \]
    The periodic chain ground space considered in this section is
    \[
        \mathcal G_{N,L}(B).
    \]
\end{definition}

\begin{lemma}[Block-diagonal chain ground space]
    \label{lem:parent_hamiltonian_ground_space_eq}
    \notready
    \uses{def:parent_hamiltonian_ground_space_bnt}
    By definition, for $B=\bigoplus_{j=1}^g \mu_j A_j$ one has
    \[
        \mathcal G_{N,L}(B)
        =
        \mathcal G_{N,L}\!\left(\bigoplus_{j=1}^g \mu_j A_j\right).
    \]
\end{lemma}

\begin{definition}[Vector space generated by a basis of normal tensors]
    \label{def:bnt_span}
    \notready
    For a basis of normal tensors $A_1,\ldots,A_g$, the vector space generated
    by their length-$N$ matrix product vectors is
    \[
        \spn\{V^{(N)}(A_j): j=1,\ldots,g\}.
    \]
\end{definition}

\begin{definition}[Finite blockwise physical--virtual correspondence]
    \label{def:blockwise_product_word_span}
    \lean{MPSTensor.WordTupleSpanTop}
    \leanok
    \uses{rem:word_tuple_span_top}
    For a fixed blocked length $m$, the finite-length blockwise
    physical--virtual correspondence is
    \[
        \spn\{(A_1^\omega,\ldots,A_g^\omega): |\omega|=m\}
        =
        \prod_{j=1}^g \MN{D_j}.
    \]
    Equivalently,
    \[
        \forall X=(X_j)_j\in\bigoplus_{j=1}^g \MN{D_j},
        \quad
        \exists (c_\omega(X))_{|\omega|=m},
        \quad
        (X_j)_{j=1}^g
        =
        \sum_{|\omega|=m} c_\omega(X)(A_1^\omega,\ldots,A_g^\omega).
    \]
    This is the fixed-length form of the block-injective correspondence of
    \cite[lines~2113--2117]{Cirac2021Matrix}.
\end{definition}

\begin{lemma}[Each normal-tensor vector lies in the chain ground space]
    \label{lem:bnt_mpv_mem_ground_space}
    \notready
    \uses{def:parent_hamiltonian_ground_space_bnt, lem:mpv_mem_chain_ground_space}
    If $1 < L$ and $N \ge L + 1$, then for each normal tensor $A_j$ in the basis
    the vector $V^{(N)}(A_j)$ lies in
    $\mathcal G_{N,L}(\bigoplus_k \mu_k A_k)$.
\end{lemma}

\begin{proof}
    \notready
    Write $\Gamma_N^B$ for the length-$N$ boundary-to-state map of the tensor
    $B$, and let $\iota_j(Z)$ be the block-diagonal matrix with $Z$ in the
    $j$-th block and $0$ elsewhere. The direct-sum evaluation gives
    \[
        \Gamma_N^{\oplus_k\mu_k A_k}\!\left(\iota_j(\mu_j^{-N}X)\right)
        =
        \Gamma_N^{A_j}(X).
    \]
    Taking $X=I$ gives the vector $V^{(N)}(A_j)$, and
    Lemma~\ref{lem:mpv_mem_chain_ground_space} applied to the individual block
    gives its membership in $\mathcal G_{N,L}(\bigoplus_k\mu_k A_k)$.
\end{proof}

\begin{lemma}[Inclusion of a block ground space]
    \label{lem:chain_ground_space_block_le_assembled}
    \lean{MPSTensor.chainGroundSpace_block_le_toTensorFromBlocks}
    \leanok
    \uses{def:chain_ground_space, lem:local_ground_space_block_le_assembled}
    For each block index $j$ with $\mu_{j} \neq 0$, the chain ground space of the block
    $A_j$ embeds into the chain ground space of the direct-sum tensor:
    \[
        \mathcal G_{N,L}(A_{j})
        \subseteq
        \mathcal G_{N,L}\!\left(\bigoplus_k \mu_k A_k\right).
    \]
\end{lemma}

\begin{proof}
    \leanok
    The cyclic-window definition of the chain ground space reduces the inclusion to
    the local block-embedding statement that each ground vector of $A_j$ on the
    window of length $L$ lifts to the direct sum by the identity
    \[
        \Gamma_L^{\oplus_k\mu_k A_k}\!\left(\iota_j(\mu_j^{-L}X)\right)
        =
        \Gamma_L^{A_j}(X).
    \]
    Here $\iota_j$ inserts a matrix in the $j$-th diagonal block and has zero
    off that block. Thus $\iota_j(\mu_j^{-L}X)$ has exactly the local vector of
    the $j$-th block on that window.
\end{proof}

\begin{lemma}[Forward inclusion from the block ground spaces]
    \label{lem:isup_chain_ground_space_block_le_assembled}
    \lean{MPSTensor.iSup_chainGroundSpace_block_le_toTensorFromBlocks}
    \leanok
    \uses{def:chain_ground_space, lem:chain_ground_space_block_le_assembled}
    Assume $\mu_{j} \neq 0$ for every block index $j$. The linear sum of the
    blockwise periodic-chain ground spaces is contained in the direct-sum
    tensor's periodic-chain ground space:
    \[
        \sum_{j} \mathcal G_{N,L}(A_{j})
        \subseteq
        \mathcal G_{N,L}\!\left(\bigoplus_k \mu_k A_k\right),
    \]
    where $j$ ranges over all block indices.
\end{lemma}

\begin{proof}
    \leanok
    Take the linear sum of the inclusions
    \[
        \mathcal G_{N,L}(A_j)
        \subseteq
        \mathcal G_{N,L}\!\left(\bigoplus_k \mu_k A_k\right)
    \]
    over all block indices $j$.
\end{proof}

\begin{lemma}[Forward inclusion for the direct-sum tensor]
    \label{lem:isup_chain_ground_space_block_le_parent_ground_space}
    \notready
    \uses{def:parent_hamiltonian_ground_space_bnt,
        lem:isup_chain_ground_space_block_le_assembled}
    Let $A^i=\bigoplus_j\mu_j A_j^i$ be in block-injective canonical form,
    where $A_1,\ldots,A_g$ is a basis of normal tensors.  The linear sum of the
    blockwise periodic-chain ground spaces is contained in the periodic chain
    ground space of the block-diagonal tensor:
    \[
        \sum_{j} \mathcal G_{N,L}(A_{j})
        \subseteq
        \mathcal G_{N,L}\!\left(\bigoplus_k \mu_k A_k\right).
    \]
\end{lemma}

\begin{proof}
    \notready
    The block-injective canonical-form hypothesis supplies $\mu_{j} \neq 0$ for
    every $j$, and the block-diagonal chain ground space is
    $\mathcal G_{N,L}(\bigoplus_j \mu_j A_j)$ by definition.
\end{proof}

\begin{lemma}[Block-diagonal commutant reduction from span extension]
    \label{lem:block_diagonal_commutant_reduction_from_projection_span}
    \notready
    \uses{def:block_projection, def:block_diagonal_matrix}
    Let $S$ be a family of virtual matrices. If each sector projection satisfies
    \[
        P_j\in\spn S
    \]
    and a boundary matrix $X$ satisfies
    \[
        XM=MX\qquad (M\in S),
    \]
    then, for every sector $j$,
    \[
        XP_j=P_jX,
    \]
    and hence
    \[
        P_jXP_k=0\qquad (j\ne k).
    \]
    Thus $X$ is block diagonal with respect to the sector decomposition.
\end{lemma}

\begin{proof}
    \notready
    \uses{lem:block_diagonal_of_commutes_block_projection,
        lem:block_diagonal_iff_off_block_zero}
    The equations $XM=MX$ are linear in $M$, so they hold for every
    $M\in\spn S$. Substituting $M=P_j$ gives $XP_j=P_jX$. Multiplying this
    identity on the left by $P_j$ and on the right by $P_k$ with $j\ne k$
    gives $P_jXP_k=0$.
\end{proof}

\begin{lemma}[Block projections from product-algebra span]
    \label{lem:block_projection_span_from_product_algebra_span}
    \notready
    Let $T_a = (T_a(i))_i$ be a family of tuples spanning the full product algebra
    $\prod_i \MN{n_i}$, and let $c_i\neq0$ be nonzero scalars. Then, for every
    sector $k$, the sector projection $P_k$ lies in the span of the scaled
    dependent block-diagonal matrices
    \[
        \bigoplus_i c_i T_a(i).
    \]
\end{lemma}

\begin{proof}
    \notready
    Apply the linear map $(M_i)_i \mapsto \bigoplus_i c_i M_i$ to the full-span
    product-algebra element with $(c_k)^{-1} I$ in the $k$-th component and zero in
    every other component. Its image is exactly $P_k$.
\end{proof}

\begin{lemma}[Sector projections from blockwise word span]
    \label{lem:block_projection_span_from_product_word_span}
    \notready
    If the simultaneous length-$m$ block-word tuples
    \[
        \omega \longmapsto (A_j^{\omega})_j
    \]
    span the full product algebra $\prod_j \MN{D_j}$ and every $\mu_j$ is nonzero, then the
    pulled-back length-$m$ word products of $\bigoplus_j\mu_j A_j$ span every
    virtual sector projection $P_j$. Consequently, any boundary matrix commuting with all
    those direct-sum words has block-diagonal pulled-back form, and its off-block entries vanish.
\end{lemma}

\begin{proof}
    \notready
    \uses{lem:block_projection_span_from_product_algebra_span,
        lem:block_diagonal_commutant_reduction_from_projection_span}
    First embed the product algebra linearly as dependent block-diagonal matrices, scaling the
    $j$-th block by $\mu_j^m$. Since $\mu_j^m\neq0$, the product-algebra span contains the
    tuple with $(\mu_j^m)^{-1}I$ in one component and zero elsewhere; its scaled diagonal image is
    exactly $P_j$. The word-evaluation formula for the direct-sum tensor identifies these scaled
    diagonal matrices with the pulled-back direct-sum word products, and the commutant conclusion
    follows from Lemma~\ref{lem:block_diagonal_commutant_reduction_from_projection_span}.
\end{proof}

\begin{definition}[Pairwise block-separating equations]
    \label{def:pairwise_block_separating_words}
    \lean{MPSTensor.HasBlockSelectorOn}
    \lean{MPSTensor.HasPairBlockSeparatingWords}
    \leanok
    Fix a finite target set $T$ of block indices. A length-$S$ block-separating
    equation for $k$ against $T$ is a choice of coefficients $c_\omega$ such that
    \[
        \sum_{|\omega|=S} c_\omega A_\ell^\omega
        =
        \begin{cases}
            I, & \ell=k,\\
            0, & \ell\in T.
        \end{cases}
    \]
    The pairwise version requires this equation for each ordered pair of
    distinct blocks, with $T=\{j\}$.
\end{definition}

\begin{definition}[Pair trace-separation criterion]
    \label{def:pair_trace_separation_words}
    \lean{MPSTensor.pairAlgSpan}
    \lean{MPSTensor.pairWordTuple}
    \lean{MPSTensor.PairWordTupleSpanTop}
    \lean{MPSTensor.PairTraceSeparatingAt}
    \lean{MPSTensor.pairEvalWordTuple}
    \lean{MPSTensor.pairCumulativeSpan}
    \lean{MPSTensor.PairCumulativeWordTupleSpanTop}
    \lean{MPSTensor.PairTraceSeparatingUpTo}
    \lean{MPSTensor.PairTraceSeparatingAll}
    \lean{MPSTensor.pairAllWordsSpan}
    \lean{MPSTensor.PairAllWordsSpanTop}
    \leanok
    For two blocks $A$ and $B$, the length-$S$ pair word tuple is
    $w \mapsto (A^w,B^w)$. The pair product-span condition says that these
    tuples span $\MN{D_A}\times\MN{D_B}$. Equivalently, in the trace-dual
    formulation, the only pair of test matrices $(\Delta_A,\Delta_B)$ whose
    trace pairing with every length-$S$ tuple vanishes is $(0,0)$. The
    cumulative variant replaces one fixed length by all word lengths at most
    $S$, and the all-length variant allows every finite word.
\end{definition}

\begin{lemma}[Three-block trace reduction from a nonzero middle matrix]
    \label{lem:three_block_trace_reduction_nonzero_middle}
    \lean{MPSTensor.exists_right_trace_test_of_three_block_trace_relation_left}
    \lean{MPSTensor.exists_left_trace_test_of_three_block_trace_relation_right}
    \leanok
    \uses{def:l_blk_injective, lem:block_injective_two_sided_word_span}
    Let $A$ be $L$-block injective, and let $\Delta_A\neq 0$.
    Suppose that, for all words $u,v,t$ of length $L$,
    \[
        \tr\!(\Delta_A A_v A_t A_u)
        + \tr\!(\Delta_B B_v B_t B_u)=0.
    \]
    Then for every $Z\in M_{D_A}(\mathbb C)$ there is a matrix
    $W\in M_{D_B}(\mathbb C)$ such that, for all words $t$ of
    length $L$,
    \[
        \tr(ZA_t)+\tr(WB_t)=0.
    \]
    This is the trace-dual algebraic step in the two-block case of the
    direct-sum lemma of \cite{PerezGarcia2007Matrix}.
\end{lemma}

\begin{proof}\leanok\uses{lem:block_injective_two_sided_word_span}
    Since $\Delta_A\neq0$, Lemma~\ref{lem:block_injective_two_sided_word_span}
    writes $Z$ as a linear combination of matrices $A_u\Delta_A A_v$.
    For one such term, cyclicity of the trace rewrites the three-block relation
    as
    \[
        \tr(A_u\Delta_A A_vA_t)
        +\tr(B_u\Delta_B B_vB_t)=0.
    \]
    Taking the same linear combination of the matrices $B_u\Delta_BB_v$
    gives the required $W$, and linearity gives the general case.
\end{proof}

\begin{lemma}[Finite-chain image inclusion from a three-block relation]
    \label{lem:three_block_relation_ground_space_inclusion}
    \lean{MPSTensor.leftTraceWordMap}
    \lean{MPSTensor.groundSpaceMap_eq_leftTraceWordMap}
    \lean{MPSTensor.groundSpace_eq_leftTraceWordMap_range}
    \lean{MPSTensor.leftTraceWordMap_range_le_of_three_block_trace_relation_left}
    \lean{MPSTensor.leftTraceWordMap_range_eq_of_three_block_trace_relation}
    \lean{MPSTensor.groundSpace_le_of_three_block_trace_relation_left}
    \lean{MPSTensor.groundSpace_eq_of_three_block_trace_relation}
    \leanok
    \uses{lem:three_block_trace_reduction_nonzero_middle}
    Under the hypotheses of
    Lemma~\ref{lem:three_block_trace_reduction_nonzero_middle}, the length-$L$
    spaces \(\mathcal G_L^A\) and \(\mathcal G_L^B\) satisfy
    \[
        \mathcal G_L^A \subseteq \mathcal G_L^B.
    \]
    If the corresponding nonzero middle matrix exists on both sides, then
    $\mathcal G_L^A=\mathcal G_L^B$.
\end{lemma}

\begin{proof}\leanok
    \uses{lem:three_block_trace_reduction_nonzero_middle}
    Write a vector in $\mathcal G_L^A$ as
    $t\mapsto \tr(A_tZ)$. The preceding lemma supplies $W$ with
    $\tr(ZA_t)+\tr(WB_t)=0$ for every word $t$. Cyclicity of the
    trace gives
    \[
        \tr(A_tZ)=\tr(ZA_t)=-\tr(WB_t)=\tr(B_t(-W)),
    \]
    so the vector lies in $\mathcal G_L^B$. The equality statement follows
    by applying the same argument with the two blocks exchanged.
\end{proof}

\begin{lemma}[Dimension step for two-block local spaces]
    \label{lem:direct_sum_two_block_dimension_step}
    \lean{MPSTensor.leftTraceWordMap_injective_of_isNBlkInjective}
    \lean{MPSTensor.leftTraceWordMap_range_finrank_eq_of_isNBlkInjective}
    \lean{MPSTensor.leftTraceWordMap_range_eq_of_range_le_of_isNBlkInjective_of_dim_ge}
    \lean{MPSTensor.leftTraceWordMap_range_eq_of_three_block_trace_relation_left_of_dim_ge}
    \lean{MPSTensor.not_three_block_trace_relation_left_of_isNBlkInjective_of_dim_gt}
    \lean{MPSTensor.groundSpace_finrank_eq_of_isNBlkInjective}
    \lean{MPSTensor.bondDim_eq_of_groundSpace_eq_of_isNBlkInjective}
    \lean{MPSTensor.bondDim_eq_and_groundSpace_eq_of_groundSpace_le_of_isNBlkInjective_of_dim_ge}
    \lean{MPSTensor.bondDim_eq_and_groundSpace_eq_of_three_block_trace_relation_left_of_dim_ge}
    \lean{MPSTensor.groundSpace_inf_eq_bot_of_not_bondDim_eq_and_groundSpace_eq_of_dim_ge}
    \lean{MPSTensor.pairTraceSeparatingAt_threeBlock_of_isNBlkInjective_of_dim_gt}
    \lean{MPSTensor.pairTraceSeparatingAt_threeBlock_of_isNBlkInjective_of_dim_ne}
    \lean{MPSTensor.chainGroundSpace_eq_of_groundSpace_eq}
    \lean{MPSTensor.mpvSubmodule_ne_of_not_exists_mpv_eq_smul}
    \lean{MPSTensor.not_exists_mpv_eq_smul_of_not_exists_forall_mpv_eq_mul}
    \lean{MPSTensor.not_bondDim_eq_and_groundSpace_eq_of_mpvSubmodule_ne}
    \lean{MPSTensor.groundSpace_inf_eq_bot_of_exists_not_forall_mpv_eq_mul_of_dim_ge}
    \lean{MPSTensor.pairTraceSeparatingAt_of_groundSpace_inf_eq_bot_of_injective_groundSpaceMap}
    \lean{MPSTensor.pairTraceSeparatingAt_of_groundSpace_inf_eq_bot_of_isNBlkInjective}
    \lean{MPSTensor.pairTraceSeparatingAt_threeBlock_of_exists_not_forall_mpv_eq_mul_of_dim_ge}
    \lean{MPSTensor.groundSpace_inf_eq_bot_of_blocksNotGaugePhaseEquiv_same_dim_of_dim_ge}
    \lean{MPSTensor.pairTraceSeparatingAt_threeBlock_of_blocksNotGaugePhaseEquiv_same_dim_of_dim_ge}
    \lean{MPSTensor.forall_pairTraceSeparatingAt_threeBlock_of_blocksNotGaugePhaseEquiv}
    \lean{MPSTensor.exists_forall_pairTraceSeparatingAt_threeBlock_of_blocksNotGaugePhaseEquiv}
    \leanok
    \uses{lem:three_block_relation_ground_space_inclusion,
        thm:chain_ground_space_eq_mpv_blk}
    If $A$ and $B$ are length-$L$ block-injective, the map
    $Z\mapsto(t\mapsto\tr(A_tZ))$ has image dimension $D_A^2$, and
    similarly for $B$. Therefore an inclusion
    $\mathcal G_L^A\subseteq \mathcal G_L^B$, together with
    $D_B\le D_A$, forces $D_A=D_B$ and
    $\mathcal G_L^A=\mathcal G_L^B$. In particular, the three-block
    trace relation from
    Lemma~\ref{lem:three_block_relation_ground_space_inclusion} gives this
    equality whenever the ordering \(D_B\le D_A\) is available. If the
    inclusion comes from the strictly larger block, $D_B<D_A$, then the same
    dimension step already contradicts the existence of the three-block trace
    relation. More generally, if an external source hypothesis rules out the
    simultaneous conclusion $D_A=D_B$ and
    $\mathcal G_L^A=\mathcal G_L^B$, then the three-block local spaces
    $\mathcal G_{3L}^A$ and $\mathcal G_{3L}^B$ have zero intersection.
    With injectivity at the three-block length, the strict-size branch itself
    already gives homogeneous pair trace separation at length $3L$.
    One such source hypothesis is non-proportionality of the periodic MPS vectors
    at a sufficiently long chain length, equivalently distinctness of the
    corresponding spaces $\mathbb C\,V^{(N)}(A)$: equality
    \(\mathcal G_L^A=\mathcal G_L^B\) gives
    equality of the periodic-chain constraint spaces, and injective uniqueness
    identifies those spaces with the two one-dimensional MPS spans. For same-dimensional
    separated BNT blocks, the non-equivalence condition supplies such
    non-proportional MPS vectors at arbitrarily large chain lengths; the
    direct-sum injectivity hypotheses remain explicit. Zero intersection of
    the two three-block local spaces also gives homogeneous pair trace
    separation at that length once the boundary-to-chain maps at the
    three-block length are injective. Combining the equal-dimension and
    unequal-dimension branches over all ordered block pairs gives one common
    three-block trace-separation length, and hence the positive-length
    blockwise word span required by the block-separation argument. In the
    source theorem this span is obtained from finite-length injectivity for
    each block. The proved specialization records the
    stronger length-one span condition separately; under that extra hypothesis
    the same trace-separation argument starts at the corresponding one-block
    length and then passes to positive multiples. This trace separation is
    equivalently a full homogeneous pair span; concatenating words propagates
    fullness to later homogeneous lengths, so the block-separation conditions
    also give a uniform finite pair-span period window.
    Combining this window with the conditional period-window homogenization
    theorem gives the homogeneous blockwise word span without separately
    assuming homogeneous identity pairs. The same conclusion applies
    to the corresponding normal canonical-form hypotheses with a basis of
    normal tensors once one-site injectivity is supplied explicitly, since those
    hypotheses carry the irreducibility and normalization assumptions used in
    the direct-sum comparison.
\end{lemma}

\begin{proof}\leanok
    \uses{lem:three_block_relation_ground_space_inclusion}
    Block injectivity makes the boundary-to-chain map injective, so
    $\dim\mathcal G_L^A=D_A^2$ and $\dim\mathcal G_L^B=D_B^2$.
    Inclusion gives $D_A^2\le D_B^2$, while $D_B\le D_A$ gives the
    reverse inequality. Hence $D_A=D_B$, and equality of dimensions inside
    the inclusion gives \(\mathcal G_L^A=\mathcal G_L^B\). In the strict-size
    branch this contradicts $D_B<D_A$. For the conditional directness
    statement, a vector in the intersection of the two three-block local spaces
    has two boundary representations. If the first boundary matrix is zero,
    the vector is zero. Otherwise the difference of the two coefficient
    formulas is exactly the homogeneous three-block trace relation, so the
    dimension step gives the forbidden collapse. In the equal-size branch,
    equal local spaces impose the same cyclic local constraints on a
    periodic chain. The injective uniqueness theorem then identifies both
    cyclic ground spaces with the spans of their MPV states, contradicting
    non-proportionality of those states.
\end{proof}

\begin{lemma}[All-length pair trace separation has finite cumulative length]
    \label{lem:all_length_pair_trace_separation_finite_cutoff}
    \lean{MPSTensor.pairEvalWordTuple_mem_pairCumulativeSpan}
    \lean{MPSTensor.pairCumulativeSpan_mono}
    \lean{MPSTensor.PairTraceSeparatingUpTo.mono}
    \lean{MPSTensor.pairTraceSeparatingUpTo_of_pairTraceSeparatingAt}
    \lean{MPSTensor.pairTraceSeparatingAt_of_pairWordTupleSpanTop}
    \lean{MPSTensor.pairWordTupleSpanTop_iff_pairTraceSeparatingAt}
    \lean{MPSTensor.pairCumulativeWordTupleSpanTop_of_pairTraceSeparatingUpTo}
    \lean{MPSTensor.pairTraceSeparatingUpTo_of_pairCumulativeWordTupleSpanTop}
    \lean{MPSTensor.pairCumulativeWordTupleSpanTop_iff_pairTraceSeparatingUpTo}
    \lean{MPSTensor.pairCumulativeWordTupleSpanTop_of_pairWordTupleSpanTop}
    \lean{MPSTensor.pairCumulativeWordTupleSpanTop_of_pairTraceSeparatingAt}
    \lean{MPSTensor.pairAllWordsSpanTop_of_pairTraceSeparatingAll}
    \lean{MPSTensor.pairTraceSeparatingAll_of_pairAllWordsSpanTop}
    \lean{MPSTensor.pairAllWordsSpanTop_iff_pairTraceSeparatingAll}
    \lean{MPSTensor.exists_pairCumulativeWordTupleSpanTop_of_pairAllWordsSpanTop}
    \lean{MPSTensor.exists_pairTraceSeparatingUpTo_of_pairTraceSeparatingAll}
    \lean{MPSTensor.exists_pairTraceSeparatingUpTo_iff_pairTraceSeparatingAll}
    \lean{MPSTensor.exists_forall_pairTraceSeparatingUpTo_of_forall_pairTraceSeparatingAll}
    \leanok
    \uses{def:pair_trace_separation_words}
    If no nonzero trace functional vanishes on all finite pair words, then
    there is a finite $S$ such that vanishing on every word of length at most
    $S$ already forces the test pair to be zero. For a finite family of ordered
    block pairs, these bounds can be replaced by one common bound.
\end{lemma}

\begin{proof}
    \leanok
    The trace condition is dual to the span of all finite pair words. If that
    span is the whole product algebra, every vector lies in some cumulative
    span. Since the product algebra is finite-dimensional, the increasing
    sequence of cumulative spans stabilizes; at the stabilization index it is
    already the whole product algebra. Applying the finite trace-duality
    criterion gives the stated finite bound. This proves the cumulative
    finite-dimensional step; the remaining BNT step is to convert it to a
    single homogeneous length. Cumulative span-top alone is not identity
    padding, because the length-zero word already contributes the identity.
\end{proof}

\begin{lemma}[Pair algebra projections are full]
    \label{lem:pair_algebra_projection_fullness}
    \lean{MPSTensor.pairAlgSpan_map_fst_eq_top_of_isInjective}
    \lean{MPSTensor.pairAlgSpan_map_snd_eq_top_of_isInjective}
    \lean{MPSTensor.pairEvalWordTuple_mem_pairAlgSpan}
    \lean{MPSTensor.pair_mul_mem_pairAllWordsSpan}
    \lean{MPSTensor.pairAlgSpan_toSubmodule_le_pairAllWordsSpan}
    \lean{MPSTensor.pairAllWordsSpanTop_of_pairAlgSpan_eq_top}
    \lean{MPSTensor.pairTraceSeparatingAll_of_pairAlgSpan_eq_top}
    \lean{MPSTensor.gaugePhaseEquiv_of_pairAlgSpan_eq_graph_algEquiv}
    \lean{MPSTensor.not_pairAlgSpan_eq_graph_algEquiv_of_not_gaugePhaseEquiv}
    \lean{MPSTensor.subdirect_matrix_pair_eq_top_of_dim_ne}
    \lean{MPSTensor.pairAlgSpan_eq_top_of_injective_dim_ne}
    \lean{MPSTensor.matrixPairSubalgebra_eq_top_of_axes}
    \lean{MPSTensor.matrixPairSubalgebra_eq_top_of_leftAxisIdeal_eq_top}
    \lean{MPSTensor.matrixPairSubalgebra_eq_top_of_rightAxisIdeal_eq_top}
    \lean{MPSTensor.matrixPairSubalgebra_eq_graph_algEquiv_of_axisIdeals_eq_bot}
    \lean{MPSTensor.subdirect_matrix_pair_eq_top_or_eq_graph_algEquiv}
    \lean{MPSTensor.pairTraceSeparatingAll_of_injective_dim_ne}
    \leanok
    \uses{def:pair_trace_separation_words}
    If one of the two blocks is injective, then the corresponding coordinate
    projection of the pair algebra generated by the simultaneous one-letter
    pairs is the full matrix algebra. Moreover, every finite pair word lies in
    this pair algebra, and the all-word linear span is the submodule generated
    by the pair algebra. If the pair algebra is the graph of an algebra
    automorphism of the matrix algebra, Skolem--Noether turns this graph case
    into gauge-phase equivalence, so under non-gauge-equivalence the graph
    branch is impossible. For two blocks of unequal bond dimensions, the graph
    branch is also impossible because it would identify matrix algebras with
    different finite dimensions over $\C$.
\end{lemma}

\begin{proof}
    \leanok
    Projecting the pair generators to one coordinate gives exactly the
    one-block generator family. Injectivity says that the linear span of this
    family is the full matrix algebra, and the linear span is contained in the
    unital algebra it generates. The pair-word membership follows by induction
    on the word, using closure of the generated algebra under multiplication.
    Conversely, the all-word linear span is closed under componentwise
    multiplication, so it contains the algebra generated by the pair generators.
    In the graph case, each generator satisfies $B_i=\varphi(A_i)$; applying
    Skolem--Noether to $\varphi$ gives an ordinary gauge, hence a gauge-phase
    equivalence with scalar $1$.
\end{proof}

\begin{lemma}[Conditional padding from exact homogeneous identity witnesses]
    \label{lem:pair_trace_separation_homogenization_with_padding}
    \lean{MPSTensor.pairEvalWordTuple_mem_span_pairWordTuple_length}
    \lean{MPSTensor.pair_mul_mem_span_pairWordTuple_add}
    \lean{MPSTensor.pairIdentity_mem_pairWordTupleSpan_zero}
    \lean{MPSTensor.pairIdentity_mem_pairWordTupleSpan_mul}
    \lean{MPSTensor.pairIdentity_mem_pairWordTupleSpan_add}
    \lean{MPSTensor.pairIdentity_mem_pairWordTupleSpan_add_mul}
    \lean{MPSTensor.pairIdentity_mem_pairWordTupleSpan_eventually_of_period_window}
    \lean{MPSTensor.pairIdentity_mem_pairWordTupleSpan_eventually_of_pairWordTupleSpanTop_period_window}
    \lean{MPSTensor.pairIdentity_mem_pairWordTupleSpan_eventually_of_pairAllWordsSpanTop_period_window}
    \lean{MPSTensor.pairWordTupleSpanTop_of_pairCumulativeWordTupleSpanTop_of_identity_padding}
    \lean{MPSTensor.pairTraceSeparatingAt_of_pairTraceSeparatingUpTo_of_identity_padding}
    \lean{MPSTensor.exists_pairTraceSeparatingAt_of_pairTraceSeparatingUpTo_of_eventual_identity_padding}
    \lean{MPSTensor.exists_pairTraceSeparatingAt_of_pairAllWordsSpanTop_of_identity_padding}
    \lean{MPSTensor.exists_forall_pairTraceSeparatingAt_of_pairTraceSeparatingAll_of_identity_padding}
    \lean{MPSTensor.exists_forall_pairTraceSeparatingAt_of_pairTraceSeparatingAll_of_identity_period_windows}
    \lean{MPSTensor.exists_forall_pairTraceSeparatingAt_of_pairSpanTop_period_windows}
    \lean{MPSTensor.exists_pairTraceSeparatingAt_of_not_gaugePhaseEquiv_of_pairWordTupleSpanTop_period_window}
    \lean{MPSTensor.exists_pairTraceSeparatingAt_of_injective_dim_ne_of_pairWordTupleSpanTop_period_window}
    \lean{MPSTensor.pairAlgSpan_eq_top_of_injective_not_gaugePhaseEquiv}
    \lean{MPSTensor.pairTraceSeparatingAll_of_injective_not_gaugePhaseEquiv}
    \lean{MPSTensor.pairTraceSeparatingAll_of_injective_not_gaugePhaseEquiv_cast_left}
    \lean{MPSTensor.pairTraceSeparatingAll_and_exists_pairTraceSeparatingUpTo_of_injective_not_gaugePhaseEquiv}
    \lean{MPSTensor.exists_pairTraceSeparatingAt_of_pairTraceSeparatingUpTo_of_identity_period_window}
    \leanok
    \uses{def:pair_trace_separation_words, lem:all_length_pair_trace_separation_finite_cutoff}
    Homogeneous identity-padding witnesses multiply by concatenation of words.
    Therefore, if one has a positive period and a full residue window of
    homogeneous pair spans, then the simultaneous identity pair lies in all
    sufficiently large homogeneous pair spans. Combined with all-length pair
    trace separation, this period-window input gives one homogeneous
    pair trace-separation length. This does not prove the block-separation
    span condition of
    \cite{PerezGarcia2007Matrix} by itself: fullness of spans up to a
    bounded length is weaker than the exact homogeneous input used here.
\end{lemma}

\begin{proof}
    \leanok
    The length-zero word gives the identity only at length zero, so it is not a
    positive-length padding theorem. The useful padding facts are homogeneous:
    if the identity pair is in the span at lengths $m$ and $n$, componentwise
    multiplication puts it in the span at length $m+n$. A period witness and a
    complete residue window therefore give homogeneous padding at every
    sufficiently large length, which is the input needed by the
    all-length-to-fixed-length implication.
\end{proof}

\begin{lemma}[Pair trace separation gives pairwise block separation]
    \label{lem:pair_trace_separation_to_pairwise_block_separation}
    \lean{MPSTensor.pairWordTupleSpanTop_of_pairTraceSeparatingAt}
    \lean{MPSTensor.hasBlockSelectorOn_of_pairWordTupleSpanTop}
    \lean{MPSTensor.hasBlockSelectorOn_of_pairTraceSeparatingAt}
    \lean{MPSTensor.hasPairBlockSeparatingWords_of_forall_pairWordTupleSpanTop}
    \lean{MPSTensor.hasPairBlockSeparatingWords_of_forall_pairTraceSeparatingAt}
    \lean{MPSTensor.exists_hasPairBlockSeparatingWords_of_exists_forall_pairTraceSeparatingAt}
    \leanok
    \uses{def:pairwise_block_separating_words, def:pair_trace_separation_words}
    At a fixed length $S$, the pair trace-separation criterion is dual to full
    pair product-span. Therefore a pair trace-separation witness supplies a
    linear combination of word evaluations that is identity on the first block
    and zero on the second.
    A common witness for all ordered distinct pairs gives pairwise
    block-separating words at the same length.
\end{lemma}

\begin{proof}
    \leanok
    If the pair span were proper, a nonzero linear functional would vanish on
    it. Nondegeneracy of the matrix trace pairing represents that functional
    as $(\Delta_A,\Delta_B)$, contradicting pair trace separation. Once the
    pair span is the whole product algebra, the tuple $(I,0)$ is in the span.
    Using the same coefficients on the ambient block family gives the required
    pairwise block-separating equation.
\end{proof}

\begin{lemma}[Block-separating equations from pairwise separation]
    \label{lem:block_separating_equations_from_pairwise_separation}
    \lean{MPSTensor.pointwise_mul_mem_span_wordTuple_add}
    \lean{MPSTensor.hasBlockSelectorOn_empty}
    \lean{MPSTensor.HasBlockSelectorOn.mul}
    \lean{MPSTensor.hasBlockSelectorOn_finset_of_pairBlockSeparatingWords}
    \lean{MPSTensor.hasBlockSelectorWords_of_forall_hasBlockSelectorOn_univ_erase}
    \lean{MPSTensor.hasBlockSelectorWords_of_pairBlockSeparatingWords}
    \lean{MPSTensor.propBlockInjective_of_common_blockInjective_of_pairBlockSeparatingWords}
    \lean{MPSTensor.wordTupleSpanTop_of_common_blockInjective_of_pairBlockSeparatingWords}
    \lean{MPSTensor.hasBiCF_of_common_blockInjective_of_pairBlockSeparatingWords}
    \leanok
    \uses{def:pairwise_block_separating_words}
    If every ordered pair of distinct blocks admits a common-length pairwise
    block-separating equation, then for each $k$ there are coefficients
    $c_\omega^{(k)}$ with
    \[
        \sum_{|\omega|=(g-1)S} c_\omega^{(k)} A_j^\omega
        =
        \begin{cases}
            I, & j=k,\\
            0, & j\ne k.
        \end{cases}
    \]
    Hence common block injectivity plus pairwise separation implies the
    finite-length blockwise physical--virtual correspondence in the
    block-injectivity proposition in \cite[lines~340--345]{Cirac2016MPDO_arXiv}
    and the block-injective canonical-form trace-separation condition.
\end{lemma}

\begin{proof}
    \leanok
    The empty word gives the identity tuple, corresponding to an empty target set.
    The span of word tuples is closed under pointwise multiplication because the
    product of a length-$L$ word and a length-$S$ word is the concatenated
    length-$(L+S)$ word. Inducting over the finite set of blocks different from
    $k$ and multiplying the corresponding pairwise equations keeps the value $I$
    on $k$ and introduces one additional zero block at each step. Extracting
    coefficients from the final tuple-span membership gives the coefficient form
    of the displayed block-separating equation.
\end{proof}

\begin{lemma}[Homogeneous padding for blockwise product spans]
    \label{lem:homogeneous_word_tuple_span_padding}
    \lean{MPSTensor.wordTupleSpanTop_add_of_identity_mem_span_wordTuple}
    \lean{MPSTensor.wordTupleSpanTop_add_of_wordTupleSpanTop}
    \leanok
    \uses{def:blockwise_product_word_span}
    Suppose the length-\(L\) simultaneous block-word tuples span
    \(\prod_j M_{D_j}(\mathbb C)\), and suppose the simultaneous identity tuple
    \((I,\ldots,I)\) lies in the span of the length-\(S\) simultaneous
    block-word tuples. Then the length-\(L+S\) simultaneous block-word tuples
    span \(\prod_j M_{D_j}(\mathbb C)\). In particular, two full homogeneous
    spans at lengths \(L\) and \(S\) give a full homogeneous span at length
    \(L+S\).
\end{lemma}

\begin{proof}
    \leanok
    Let \(M=(M_j)_j\) be a target tuple.  Write
    \[
        M_j=\sum_{|u|=L} a_u A^j_u,
        \qquad
        I=\sum_{|v|=S} b_v A^j_v
    \]
    simultaneously for every block \(j\).  Multiplying the two expansions gives
    \[
        M_j
        =
        M_j I
        =
        \sum_{|u|=L}\sum_{|v|=S} a_u b_v A^j_uA^j_v
        =
        \sum_{|u|=L}\sum_{|v|=S} a_u b_v A^j_{uv}.
    \]
    The words \(uv\) have length \(L+S\), so the target tuple lies in the
    length-\(L+S\) span.
\end{proof}

\begin{lemma}[Period-window propagation for blockwise product spans]
    \label{lem:period_window_word_tuple_span_padding}
    \lean{MPSTensor.wordTupleSpanTop_add_mul_of_wordTupleSpanTop}
    \lean{MPSTensor.wordTupleSpanTop_eventually_of_wordTupleSpanTop_period_window}
    \leanok
    \uses{def:blockwise_product_word_span,
        lem:homogeneous_word_tuple_span_padding}
    Let \(p>0\). Suppose the length-\(p\) simultaneous block-word tuples span
    \(\prod_j M_{D_j}(\mathbb C)\), and suppose that for every
    \(0\le r<p\), the length-\(s+r\) simultaneous block-word tuples span
    \(\prod_j M_{D_j}(\mathbb C)\). Then there is a length \(N\) such that,
    for every \(n\ge N\), the length-\(n\) simultaneous block-word tuples span
    \(\prod_j M_{D_j}(\mathbb C)\). More concretely, one may take \(N=s\):
    writing
    \[
        n=s+r+qp,\qquad 0\le r<p,
    \]
    the span at length \(s+r\) propagates to the span at length \(n\) by
    multiplying \(q\) times by the period-\(p\) span.
\end{lemma}

\begin{proof}
    \leanok
    The Euclidean division of \(n-s\) by \(p\) gives
    \(n=s+r+qp\) with \(0\le r<p\).  The assumed residue-window span gives the
    full product algebra at length \(s+r\). Applying homogeneous padding with
    the length-\(p\) span once gives the full product algebra at length
    \(s+r+p\); iterating gives the full product algebra at length \(s+r+qp\).
\end{proof}

\begin{lemma}[Blockwise word span from block-separating equations]
    \label{lem:product_word_span_from_bnt_block_separation}
    \lean{MPSTensor.wordTupleSpanTop_of_common_blockInjective_of_blockSelectorWords}
    \lean{MPSTensor.wordTupleSpanTop_of_common_blockInjective_of_pairBlockSeparatingWords}
    \leanok
    \uses{def:blockwise_product_word_span,
        lem:block_separating_equations_from_pairwise_separation,
        lem:homogeneous_word_tuple_span_padding}
    Let $J$ be a finite set of $r$ blocks, and let $(A_j)_{j\in J}$ be a
    finite family of tensors with common block injectivity at length $L$.
    Suppose there are length-$S$ block-separating equations: for each block
    $k$,
    \[
        \sum_{|\omega|=S} c_\omega^{(k)} A_j^\omega
        =
        \begin{cases}
            I, & j=k,\\
            0, & j\ne k.
        \end{cases}
    \]
    Then the simultaneous length-$(L+S)$ block-word evaluations span
    $\prod_j \MN{D_j}$. If the block-separating equations are first obtained
    from pairwise block-separating equations at length $S$, the same conclusion
    holds at length $L+(r-1)S$.
\end{lemma}

\begin{proof}
    \leanok
    \uses{lem:block_separating_equations_from_pairwise_separation}
    Block injectivity writes each target block matrix as a length-$L$ linear
    combination of word evaluations. For a target tuple $(M_j)_j$, write $M_k$
    in this form in block $k$, multiply by the block-separating equation for
    $k$, and sum over $k$. Thus, if
    $M_k=\sum_{|u|=L}a_u^{(k)}A_k^u$ and
    $\sum_{|v|=S}b_v^{(k)}A_j^v=\delta_{j,k}I$, then for every block $j$,
    \[
        M_j=
        \sum_{k\in J}\sum_{|u|=L}\sum_{|v|=S}
        a_u^{(k)}b_v^{(k)} A_j^u A_j^v .
    \]
    The right side is a linear combination of length-$(L+S)$ word tuples, so
    these tuples span the full product algebra.
    If one starts with pairwise separation instead, first use
    Lemma~\ref{lem:block_separating_equations_from_pairwise_separation} to
    multiply the pairwise equations into the full equations.
\end{proof}

\begin{lemma}[BNT block separation gives a finite blockwise product span]
    \label{lem:bnt_direct_sum_block_separation_word_span}
    \lean{MPSTensor.hasPairBlockSeparatingWords_threeBlock_of_blocksNotGaugePhaseEquiv_c1}
    \lean{MPSTensor.propBlockInjective_of_blocksNotGaugePhaseEquiv_directSum_c1}
    \lean{MPSTensor.wordTupleSpanTop_of_blocksNotGaugePhaseEquiv_directSum_selectors_c1}
    \lean{MPSTensor.exists_wordTupleSpanTop_of_blocksNotGaugePhaseEquiv_directSum_selectors_c1}
    \leanok
    \uses{lem:direct_sum_two_block_dimension_step,
        lem:pair_trace_separation_to_pairwise_block_separation,
        lem:block_separating_equations_from_pairwise_separation,
        lem:product_word_span_from_bnt_block_separation}
    Let \(A^1,\ldots,A^r\) be separated BNT blocks satisfying the
    irreducibility, left-canonicality, and normalized self-overlap hypotheses
    used in the direct-sum comparison. Fix \(L_0>0\). Assume that, for every
    block \(j\), the map
    \[
        \Gamma_{L_0}^{A^j}:M_{D_j}(\mathbb C)\to
        (\mathbb C^d)^{\otimes L_0},
        \qquad
        X\longmapsto
        \sum_{|u|=L_0}\tr(XA^j_u)\ket{u}
    \]
    is injective, and that the homogeneous word spans at the two source
    lengths are full:
    \[
        \operatorname{span}\{A^j_u:|u|=L_0+1\}=M_{D_j}(\mathbb C),
        \qquad
        \operatorname{span}\{A^j_v:|v|=3(L_0+1)\}=M_{D_j}(\mathbb C).
    \]
    Then for every ordered pair \(k\ne j\) there are coefficients
    \(c^{k,j}_\omega\) such that
    \[
        \sum_{|\omega|=3(L_0+1)} c^{k,j}_\omega A^k_\omega=I,
        \qquad
        \sum_{|\omega|=3(L_0+1)} c^{k,j}_\omega A^j_\omega=0.
    \]
    Consequently the simultaneous block-word tuples at length
    \[
        (L_0+1)+(r-1)3(L_0+1)
    \]
    span the full product algebra
    \[
        \prod_{j=1}^r M_{D_j}(\mathbb C).
    \]
\end{lemma}

\begin{proof}
    \leanok
    Lemma~\ref{lem:direct_sum_two_block_dimension_step} gives homogeneous
    trace separation at length \(3(L_0+1)\) for every ordered pair of distinct
    BNT blocks. Lemma~\ref{lem:pair_trace_separation_to_pairwise_block_separation}
    converts these trace-separation equations into the displayed pairwise
    block-separating equations. Multiplying the equations for the \(r-1\)
    blocks different from \(k\) gives a block-separating equation for \(k\), and
    Lemma~\ref{lem:product_word_span_from_bnt_block_separation} then supplies the
    displayed product span.
\end{proof}

\begin{lemma}[Blockwise boundary compatibility from traces]
    \label{lem:blockwise_boundary_compatibility_from_traces}
    \lean{MPSTensor.pgvwc07_blockwise_compatibility_of_trace_decomposition}
    \leanok
    Let \(A^1,\ldots,A^g\) be block tensors with the following common
    word-span property: for a fixed \(m\), the simultaneous tuples
    \[
        w\longmapsto
        \left(A^1_w,\ldots,A^g_w\right),
        \qquad |w|=m,
    \]
    span the full product algebra
    \(\prod_j M_{D_j}(\mathbb C)\).
    Suppose that matrices \(C^j_a,D^j_b\in M_{D_j}(\mathbb C)\) satisfy,
    for all physical indices \(a,b\) and all words \(w\) of length \(m\),
    \[
        \sum_j \operatorname{tr}\!\left(A^j_b C^j_a A^j_w\right)
        =
        \sum_j \operatorname{tr}\!\left(D^j_b A^j_a A^j_w\right).
    \]
    Then
    \[
        A^j_bC^j_a=D^j_bA^j_a
    \]
    for every block \(j\) and all \(a,b\).
\end{lemma}

\begin{proof}
    \leanok
    For fixed \(a,b\), set
    \[
        \Delta_j=A^j_bC^j_a-D^j_bA^j_a .
    \]
    The assumed trace equality says that
    \[
        \sum_j \operatorname{tr}\!\left(\Delta_jA^j_w\right)=0
    \]
    for every word \(w\) of length \(m\).  Since the tuples
    \((A^1_w,\ldots,A^g_w)\) span the product algebra, the product trace
    pairing is nondegenerate and hence every \(\Delta_j\) is zero.
\end{proof}

\begin{lemma}[Normalized boundary matrix identities]
    \label{lem:normalized_boundary_matrix_identities}
    \lean{MPSTensor.pgvwc07_boundary_matrix_identities_of_compatibility}
    \leanok
    Let $(A_a)_{a\in I}$, $(C_a)_{a\in I}$, and $(D_b)_{b\in I}$ be
    matrices in $M_D(\mathbb C)$, indexed by a finite physical alphabet. If
    \[
        \sum_a A_a A_a^\dagger=\Id,
        \qquad
        A_b C_a=D_b A_a
    \]
    for all $a,b\in I$, and if
    \[
        E=\sum_a C_a A_a^\dagger,
    \]
    then
    \[
        D_b=A_bE,
        \qquad
        A_bC_a=A_bEA_a
    \]
    for all $a,b\in I$.
\end{lemma}

\begin{proof}
    \leanok
    Multiply $D_b$ by the normalization:
    \[
        D_b
        =
        D_b\sum_a A_aA_a^\dagger
        =
        \sum_a D_bA_aA_a^\dagger
        =
        \sum_a A_bC_aA_a^\dagger
        =
        A_bE .
    \]
    Substituting this in $A_bC_a=D_bA_a$ gives
    $A_bC_a=A_bEA_a$.
\end{proof}

\begin{lemma}[Left-boundary summands are ground-space vectors]
    \label{lem:left_boundary_summands_mem_block_ground_spaces}
    \lean{MPSTensor.pgvwc07_sum_leftBoundaryComponents_mem_iSup_groundSpace}
    \leanok
    Let \(A^1,\ldots,A^g\) be block tensors, and let
    \(C^j_a,E^j\in M_{D_j}(\mathbb C)\) satisfy
    \[
        A^j_b C^j_a=A^j_bE^jA^j_a
    \]
    for every block \(j\) and all physical indices \(a,b\).
    For a word \(\sigma=(\sigma_1,\ldots,\sigma_{n+2})\), define
    \[
        \alpha_j(\sigma)
        =
        \operatorname{tr}\!\left(
            A^j_{\sigma_{n+2}} C^j_{\sigma_1}
            A^j_{\sigma_2}\cdots A^j_{\sigma_{n+1}}
        \right).
    \]
    Then
    \[
        \sum_j \alpha_j
        \in
        \bigvee_j G_{n+2}(A^j).
    \]
\end{lemma}

\begin{proof}
    \leanok
    Fix \(j\).  By the assumed boundary identity and cyclicity of the
    trace,
    \[
        \operatorname{tr}\!\left(
            A^j_{\sigma_{n+2}} C^j_{\sigma_1}
            A^j_{\sigma_2}\cdots A^j_{\sigma_{n+1}}
        \right)
        =
        \operatorname{tr}\!\left(
            A^j_{\sigma_1}A^j_{\sigma_2}\cdots
            A^j_{\sigma_{n+2}}E^j
        \right).
    \]
    Thus \(\alpha_j\) is the image of \(E^j\) under the parametrization
    of \(G_{n+2}(A^j)\).  Hence each \(\alpha_j\) lies in
    \(G_{n+2}(A^j)\), and the finite sum lies in the supremum of these
    subspaces.
\end{proof}

\begin{lemma}[Left-boundary trace decompositions lie in the block ground spaces]
    \label{lem:left_boundary_trace_decomposition_mem_block_ground_spaces}
    \lean{MPSTensor.pgvwc07_mem_iSup_groundSpace_of_trace_decomposition}
    \leanok
    Let \(A^1,\ldots,A^g\) be block tensors with common word-span separation at
    length \(n\), and suppose that
    \[
        \sum_a A^j_a A^{j\,\dagger}_a=\Id
    \]
    for every block \(j\).  Suppose a vector \(\psi\in(\mathbb C^d)^{\otimes(n+2)}\)
    has the left-boundary trace decomposition
    \[
        \psi=\sum_j\alpha_j,
        \qquad
        \alpha_j(i_1,\ldots,i_{n+2})
        =
        \operatorname{tr}\!\left(
            A^j_{i_{n+2}} C^j_{i_1}
            A^j_{i_2}\cdots A^j_{i_{n+1}}
        \right),
    \]
    and suppose that the two trace decompositions agree for every middle word:
    \[
        \sum_j
        \operatorname{tr}\!\left(
            A^j_b C^j_a A^j_w
        \right)
        =
        \sum_j
        \operatorname{tr}\!\left(
            D^j_b A^j_a A^j_w
        \right).
    \]
    Then
    \[
        \psi\in\bigvee_j G_{n+2}(A^j).
    \]
\end{lemma}

\begin{proof}
    \leanok
    The trace-decomposition equality and the common word-span hypothesis give
    \[
        A^j_bC^j_a=D^j_bA^j_a
    \]
    for every \(j,a,b\).  By the normalization,
    \[
        D^j_b
        =
        D^j_b\sum_a A^j_aA^{j\,\dagger}_a
        =
        \sum_a A^j_bC^j_aA^{j\,\dagger}_a
        =
        A^j_bE^j,
        \qquad
        E^j:=\sum_a C^j_aA^{j\,\dagger}_a.
    \]
    Hence \(A^j_bC^j_a=A^j_bE^jA^j_a\).  Lemma
    \ref{lem:left_boundary_summands_mem_block_ground_spaces} then gives
    \[
        \psi=\sum_j\alpha_j\in\bigvee_j G_{n+2}(A^j).
    \]
\end{proof}

\begin{lemma}[One-step block intersection from block restrictions]
    \label{lem:block_restrictions_mem_block_ground_spaces}
    \lean{MPSTensor.pgvwc07_mem_iSup_groundSpace_of_iSup_restrictions}
    \leanok
    \uses{lem:left_boundary_trace_decomposition_mem_block_ground_spaces}
    Let \(A^1,\ldots,A^g\) be block tensors with common word-span separation at
    length \(n\), and suppose
    \[
        \sum_a A^j_aA^{j\,\dagger}_a=\Id
    \]
    for every block \(j\).  If \(\psi\in(\mathbb C^d)^{\otimes(n+2)}\) satisfies
    \[
        \psi(a,-)\in\bigvee_jG_{n+1}(A^j),
        \qquad
        \psi(-,b)\in\bigvee_jG_{n+1}(A^j)
    \]
    for every physical index \(a,b\), then
    \[
        \psi\in\bigvee_jG_{n+2}(A^j).
    \]
\end{lemma}

\begin{proof}
    \leanok
    Choose finite decompositions of the fixed-boundary vectors:
    \[
        \psi(a,i_2,\ldots,i_{n+2})
        =
        \sum_j
        \operatorname{tr}\!\left(
            A^j_{i_2}\cdots A^j_{i_{n+2}}C^j_a
        \right),
    \]
    and
    \[
        \psi(i_1,\ldots,i_{n+1},b)
        =
        \sum_j
        \operatorname{tr}\!\left(
            A^j_{i_1}\cdots A^j_{i_{n+1}}D^j_b
        \right).
    \]
    Evaluating these two decompositions on
    \((a,w,b)\), where \(w\) has length \(n\), gives
    \[
        \sum_j\operatorname{tr}\!\left(A^j_bC^j_aA^j_w\right)
        =
        \sum_j\operatorname{tr}\!\left(D^j_bA^j_aA^j_w\right).
    \]
    The first decomposition also gives
    \[
        \psi=\sum_j\alpha_j,\qquad
        \alpha_j(i_1,\ldots,i_{n+2})
        =
        \operatorname{tr}\!\left(
            A^j_{i_{n+2}}C^j_{i_1}
            A^j_{i_2}\cdots A^j_{i_{n+1}}
        \right).
    \]
    Lemma~\ref{lem:left_boundary_trace_decomposition_mem_block_ground_spaces}
    therefore gives
    \[
        \psi\in\bigvee_jG_{n+2}(A^j).
    \]
\end{proof}

\begin{lemma}[One-step block intersection characterization]
    \label{lem:block_restrictions_iff_block_ground_spaces}
    \lean{MPSTensor.pgvwc07_mem_iSup_groundSpace_iff_iSup_restrictions}
    \leanok
    \uses{lem:block_restrictions_mem_block_ground_spaces}
    Under the hypotheses of
    Lemma~\ref{lem:block_restrictions_mem_block_ground_spaces}, an
    \((n+2)\)-site vector satisfies
    \[
        \psi\in\bigvee_jG_{n+2}(A^j)
    \]
    if and only if
    \[
        \psi(a,-)\in\bigvee_jG_{n+1}(A^j),
        \qquad
        \psi(-,b)\in\bigvee_jG_{n+1}(A^j)
    \]
    for every physical index \(a,b\).
\end{lemma}

\begin{proof}
    \leanok
    The reverse implication is
    Lemma~\ref{lem:block_restrictions_mem_block_ground_spaces}.  For the
    forward implication, write
    \[
        \psi=\sum_j\gamma_j,\qquad \gamma_j\in G_{n+2}(A^j).
    \]
    Each \(\gamma_j\) has fixed-boundary restrictions in \(G_{n+1}(A^j)\):
    \[
        \gamma_j(a,-)\in G_{n+1}(A^j),
        \qquad
        \gamma_j(-,b)\in G_{n+1}(A^j).
    \]
    Taking the finite sum over \(j\) gives the two restrictions of \(\psi\) in
    \(\bigvee_jG_{n+1}(A^j)\).
\end{proof}

\begin{lemma}[Subspace form of the one-step block intersection]
    \label{lem:block_restriction_intersection_subspace}
    \lean{MPSTensor.pgvwc07_iSup_groundSpace_eq_restriction_intersection}
    \leanok
    \uses{lem:block_restrictions_iff_block_ground_spaces}
    Under the hypotheses of
    Lemma~\ref{lem:block_restrictions_mem_block_ground_spaces}, set
    \[
        S_n:=\bigvee_jG_{n+1}(A^j).
    \]
    Then
    \[
        \left\{\psi:\forall b,\ \psi(-,b)\in S_n\right\}
        \cap
        \left\{\psi:\forall a,\ \psi(a,-)\in S_n\right\}
        =
        \bigvee_jG_{n+2}(A^j).
    \]
\end{lemma}

\begin{proof}
    \leanok
    Lemma~\ref{lem:block_restrictions_iff_block_ground_spaces} gives the
    displayed equivalence for each vector \(\psi\).  Interpreting the two sides
    of that equivalence as subspaces gives the stated equality.
\end{proof}

\begin{lemma}[Period-window form of the block intersection identity]
    \label{lem:period_window_block_restriction_intersection}
    \lean{MPSTensor.pgvwc07_iSup_restriction_intersection_eventually_of_period_window}
    \leanok
    \uses{lem:period_window_word_tuple_span_padding,
        lem:block_restriction_intersection_subspace}
    Let \(p>0\). Suppose the length-\(p\) simultaneous block-word tuples span
    \(\prod_j M_{D_j}(\mathbb C)\), and suppose that for every \(0\le r<p\),
    the length-\(s+r\) simultaneous block-word tuples span
    \(\prod_j M_{D_j}(\mathbb C)\). Assume also the normalization
    \[
        \sum_a A^j_aA^{j\,\dagger}_a=I
    \]
    for every block \(j\). Then there is a length \(N\) such that, for every
    \(n\ge N\), if
    \[
        S_n:=\bigvee_jG_{n+1}(A^j),
    \]
    then
    \[
        \left\{\psi:\forall b,\ \psi(-,b)\in S_n\right\}
        \cap
        \left\{\psi:\forall a,\ \psi(a,-)\in S_n\right\}
        =
        \bigvee_jG_{n+2}(A^j).
    \]
\end{lemma}

\begin{proof}
    \leanok
    Lemma~\ref{lem:period_window_word_tuple_span_padding} gives a length \(N\)
    such that the simultaneous block-word tuples span
    \(\prod_j M_{D_j}(\mathbb C)\) at every length \(n\ge N\). For each such
    \(n\), Lemma~\ref{lem:block_restriction_intersection_subspace} applies and
    gives the displayed subspace equality.
\end{proof}

\begin{lemma}[BNT block separation gives one block-intersection step]
    \label{lem:bnt_direct_sum_block_separation_block_intersection}
    \lean{MPSTensor.pgvwc07_iSup_restriction_intersection_of_bnt_directSum_selectors_c1}
    \leanok
    \uses{lem:bnt_direct_sum_block_separation_word_span,
        lem:block_restriction_intersection_subspace}
    Under the hypotheses of
    Lemma~\ref{lem:bnt_direct_sum_block_separation_word_span}, assume in addition the
    normalization
    \[
        \sum_a A^j_aA^{j\,\dagger}_a=I
    \]
    for every block \(j\). Let
    \[
        n=(L_0+1)+(r-1)3(L_0+1),
        \qquad
        S_n=\bigvee_jG_{n+1}(A^j).
    \]
    Then
    \[
        \left\{\psi:\forall b,\ \psi(-,b)\in S_n\right\}
        \cap
        \left\{\psi:\forall a,\ \psi(a,-)\in S_n\right\}
        =
        \bigvee_jG_{n+2}(A^j).
    \]
\end{lemma}

\begin{proof}
    \leanok
    Lemma~\ref{lem:bnt_direct_sum_block_separation_word_span} gives
    \[
        \operatorname{span}\{(A^1_w,\ldots,A^r_w): |w|=n\}
        =
        \prod_jM_{D_j}(\mathbb C).
    \]
    Substituting this span into
    Lemma~\ref{lem:block_restriction_intersection_subspace}, together with the
    normalization, gives the stated subspace equality.
\end{proof}

\begin{lemma}[BNT block separation with a block-injectivity window]
    \label{lem:bnt_direct_sum_block_separation_period_window_block_intersection}
    \lean{MPSTensor.pgvwc07_iSup_restriction_intersection_eventually_of_bnt_directSum_period_window}
    \leanok
    \uses{lem:direct_sum_two_block_dimension_step,
        lem:pair_trace_separation_to_pairwise_block_separation,
        lem:block_separating_equations_from_pairwise_separation,
        lem:period_window_block_restriction_intersection}
    Let \(A^1,\ldots,A^r\) be separated BNT blocks satisfying the
    irreducibility, left-canonicality, and normalized self-overlap hypotheses
    used in the direct-sum comparison. Assume \(L>1\), and assume that, for
    every block \(j\),
    \[
        \operatorname{span}\{A^j_a:a=0,\ldots,d-1\}=M_{D_j}(\mathbb C),
        \qquad
        \operatorname{span}\{A^j_u:|u|=L\}=M_{D_j}(\mathbb C),
        \qquad
        \operatorname{span}\{A^j_v:|v|=L+(L+L)\}=M_{D_j}(\mathbb C).
    \]
    Set
    \[
        S=L+(L+L),
        \qquad
        q=(r-1)S.
    \]
    Suppose that \(p>0\) and let \(s_0\) be a starting length. Assume that, for
    every block \(j\),
    \[
        \operatorname{span}\{A^j_u:|u|=p\}=M_{D_j}(\mathbb C),
    \]
    and that, for every \(0\le s<p+q\) and every block \(j\),
    \[
        \operatorname{span}\{A^j_v:|v|=s_0+s\}=M_{D_j}(\mathbb C).
    \]
    Assume also the normalization
    \[
        \sum_a A^j_aA^{j\,\dagger}_a=I
    \]
    for every block \(j\). Then there is \(N\) such that, for every \(n\ge N\),
    if
    \[
        S_n:=\bigvee_jG_{n+1}(A^j),
    \]
    then
    \[
        \left\{\psi:\forall b,\ \psi(-,b)\in S_n\right\}
        \cap
        \left\{\psi:\forall a,\ \psi(a,-)\in S_n\right\}
        =
        \bigvee_jG_{n+2}(A^j).
    \]
\end{lemma}

\begin{proof}
    \leanok
    The direct-sum dimension step gives, for each ordered pair of distinct
    blocks, length-\(S\) equations which are the identity on the first block and
    zero on the second. Multiplying these equations over the \(r-1\) other
    blocks gives a block-separating equation of total length \(q\). Combining
    this equation with block injectivity at length \(p\) gives the full product
    span at length \(p+q\). Combining it with the assumed block-injectivity
    window gives the full product span at lengths \(s_0+q+s\), for
    \(0\le s<p+q\). The period-window form of the block-intersection identity
    then gives the displayed equality for all sufficiently large \(n\).
\end{proof}

\begin{lemma}[Normalized block-separating equations give all large product spans]
    \label{lem:unital_block_separating_product_span}
    \lean{MPSTensor.wordTupleSpanTop_of_ge_of_common_blockInjective_of_unital_of_pairBlockSeparatingWords}
    \lean{MPSTensor.wordTupleSpanTop_of_ge_of_bnt_directSum_unital_c1}
    \leanok
    \uses{def:pairwise_block_separating_words,
        lem:bnt_direct_sum_block_separation_word_span,
        thm:wordspan_propagation_unital}
    Let \(A^1,\ldots,A^r\) be blocks satisfying the normalization
    \[
        \sum_a A^j_aA^{j\,\dagger}_a=I
    \]
    and assume that every block is injective at length \(L\):
    \[
        \operatorname{span}\{A^j_u:|u|=L\}=M_{D_j}(\mathbb C).
    \]
    If there are block-separating equations of length \(S\), then, for every
    \(n\ge L+(r-1)S\),
    \[
        \operatorname{span}\{(A^1_w,\ldots,A^r_w):|w|=n\}
        =
        \prod_j M_{D_j}(\mathbb C).
    \]
    In particular, for separated BNT blocks satisfying the same normalization,
    if \(L_0>0\) and \(\Gamma_{L_0}^{A^j}\) is injective for every block \(j\),
    the same conclusion holds with
    \[
        S=3(L_0+1),
        \qquad
        n\ge (L_0+1)+(r-1)3(L_0+1).
    \]
\end{lemma}

\begin{proof}
    \leanok
    By Theorem~\ref{thm:wordspan_propagation_unital}, the normalization
    propagates the length-\(L\) span equality to every larger homogeneous length:
    \[
        m\ge L
        \quad\Longrightarrow\quad
        \operatorname{span}\{A^j_u:|u|=m\}=M_{D_j}(\mathbb C).
    \]
    For \(n\ge L+(r-1)S\), take \(m=n-(r-1)S\). Combining the length-\(m\)
    block span with the \(r-1\) block-separating equations gives the displayed
    product span at length \(n\). In the BNT case, the injectivity of
    \(\Gamma_{L_0}^{A^j}\) gives the full span at length \(L_0\), and the
    normalization propagates it to the lengths \(L_0+1\) and \(3(L_0+1)\).
    Lemma~\ref{lem:bnt_direct_sum_block_separation_word_span} then supplies the
    block-separating equations with \(S=3(L_0+1)\).
\end{proof}

\begin{lemma}[Product span makes the sum of local spaces direct]
    \label{lem:product_span_block_image_direct_sum}
    \lean{MPSTensor.groundSpace_iSupIndep_of_wordTupleSpanTop}
    \lean{MPSTensor.groundSpace_iSupIndep_of_ge_of_bnt_directSum_unital_c1}
    \leanok
    \uses{def:ground_space_parent, lem:unital_block_separating_product_span}
    Suppose
    \[
        \operatorname{span}\{(A^1_w,\ldots,A^r_w):|w|=n\}
        =
        \prod_jM_{D_j}(\mathbb C).
    \]
    Then the sum of the local spaces \(G_n(A^j)\) is direct: if
    \[
        \phi_j\in G_n(A^j),
        \qquad
        \sum_j\phi_j=0,
    \]
    then \(\phi_j=0\) for every \(j\).  In particular, for the normalized BNT
    hypotheses above, this directness holds for every
    \[
        n\ge (L_0+1)+(r-1)3(L_0+1).
    \]
\end{lemma}

\begin{proof}
    \leanok
    Write
    \[
        \phi_j(\sigma)=\tr(A^j_\sigma X_j).
    \]
    Evaluating \(\sum_j\phi_j=0\) on each word \(w\) of length \(n\) gives
    \[
        \sum_j\tr(A^j_wX_j)=0.
    \]
    Since \(\tr(A^j_wX_j)=\tr(X_jA^j_w)\), the product-span equation gives, for
    every \((M_1,\ldots,M_r)\in\prod_jM_{D_j}(\mathbb C)\),
    \[
        \sum_j\tr(X_jM_j)=0.
    \]
    Taking only the \(j\)-th component nonzero and using nondegeneracy of the
    trace pairing yields \(X_j=0\), hence \(\phi_j=0\). The BNT statement follows
    by substituting Lemma~\ref{lem:unital_block_separating_product_span}.
\end{proof}

\begin{lemma}[Normalized BNT blocks give the large-length block intersection]
    \label{lem:bnt_direct_sum_unital_block_intersection_ge}
    \lean{MPSTensor.pgvwc07_directSum_restriction_intersection_of_wordTupleSpanTop}
    \lean{MPSTensor.pgvwc07_iSup_restriction_intersection_of_ge_of_bnt_directSum_unital_c1}
    \lean{MPSTensor.pgvwc07_directSum_restriction_intersection_of_ge_of_bnt_directSum_unital_c1}
    \leanok
    \uses{lem:unital_block_separating_product_span,
        lem:product_span_block_image_direct_sum,
        lem:block_restriction_intersection_subspace}
    Let \(A^1,\ldots,A^r\) be separated BNT blocks satisfying the irreducibility,
    left-canonicality, and normalized self-overlap hypotheses used in the
    direct-sum comparison. Suppose \(L_0>0\), and for every block \(j\) the map
    \[
        \Gamma_{L_0}^{A^j}:M_{D_j}(\mathbb C)\to
        (\mathbb C^d)^{\otimes L_0}
    \]
    is injective. Assume also the normalization
    \[
        \sum_a A^j_aA^{j\,\dagger}_a=I .
    \]
    Set
    \[
        N=(L_0+1)+(r-1)3(L_0+1).
    \]
    For every \(n\ge N\), set
    \[
        S_n:=\bigvee_jG_{n+1}(A^j),
    \]
    and the two consecutive sums are internal direct sums:
    \[
        \bigvee_jG_{n+1}(A^j)
        =
        \bigoplus_jG_{n+1}(A^j),
        \qquad
        \bigvee_jG_{n+2}(A^j)
        =
        \bigoplus_jG_{n+2}(A^j).
    \]
    With this notation,
    \[
        \left\{\psi:\forall b,\ \psi(-,b)\in S_n\right\}
        \cap
        \left\{\psi:\forall a,\ \psi(a,-)\in S_n\right\}
        =
        \bigvee_jG_{n+2}(A^j).
    \]
\end{lemma}

\begin{proof}
    \leanok
    Lemma~\ref{lem:unital_block_separating_product_span} gives
    \[
        \operatorname{span}\{(A^1_w,\ldots,A^r_w):|w|=n\}
        =
        \prod_jM_{D_j}(\mathbb C)
    \]
    for every \(n\ge N\). Substituting this product-span equation into
    Lemma~\ref{lem:block_restriction_intersection_subspace}, together with the
    normalization equation, gives the displayed intersection identity.
    Lemma~\ref{lem:product_span_block_image_direct_sum} at lengths \(n+1\) and
    \(n+2\) gives the two directness statements.
\end{proof}

\begin{lemma}[Normalized BNT blocks give an eventual block intersection]
    \label{lem:bnt_direct_sum_unital_block_intersection}
    \lean{MPSTensor.pgvwc07_iSup_restriction_intersection_eventually_of_bnt_directSum_unital_c1}
    \leanok
    \uses{lem:bnt_direct_sum_unital_block_intersection_ge}
    Let \(A^1,\ldots,A^r\) be separated BNT blocks satisfying the irreducibility,
    left-canonicality, and normalized self-overlap hypotheses used in the
    direct-sum comparison. Suppose \(L_0>0\), and for every block \(j\) the map
    \[
        \Gamma_{L_0}^{A^j}:M_{D_j}(\mathbb C)\to
        (\mathbb C^d)^{\otimes L_0}
    \]
    is injective. Assume also the normalization
    \[
        \sum_a A^j_aA^{j\,\dagger}_a=I .
    \]
    Then there is a length \(N\) such that, for every \(n\ge N\), with
    \[
        S_n:=\bigvee_jG_{n+1}(A^j),
    \]
    one has
    \[
        \left\{\psi:\forall b,\ \psi(-,b)\in S_n\right\}
        \cap
        \left\{\psi:\forall a,\ \psi(a,-)\in S_n\right\}
        =
        \bigvee_jG_{n+2}(A^j).
    \]
\end{lemma}

\begin{proof}
    \leanok
    Apply Lemma~\ref{lem:bnt_direct_sum_unital_block_intersection_ge} and take
    \[
        N=(L_0+1)+(r-1)3(L_0+1).
    \]
\end{proof}

\begin{lemma}[Block-diagonal boundary-condition formula]
    \label{lem:block_diagonal_boundary_condition_sum}
    \lean{MPSTensor.BlockSumGroundSpace.groundSpaceMap_toTensorFromBlocks_eq_sum_blockDiagonal}
    \leanok
    \uses{def:ground_space_parent}
    Let
    \[
        B=\bigoplus_{j=1}^g\mu_jA_j.
    \]
    For block boundary matrices \(X_j\),
    \[
        \Gamma_L^B\!\left(\bigoplus_{j=1}^g X_j\right)
        =
        \sum_{j=1}^g\Gamma_L^{A_j}(\mu_j^LX_j).
    \]
\end{lemma}

\begin{proof}
    \leanok
    This is the diagonal-block trace formula specialized to a block-diagonal
    boundary matrix.  The \(j\)-th diagonal block of
    \(\bigoplus_kX_k\) is \(X_j\), giving the displayed sum.
\end{proof}

\begin{lemma}[One block in a block-diagonal local ground space]
    \label{lem:local_ground_space_block_le_assembled}
    \lean{MPSTensor.groundSpace_block_le_toTensorFromBlocks}
    \leanok
    \uses{def:ground_space_parent}
    Let
    \[
        B=\bigoplus_{k=1}^g\mu_kA_k.
    \]
    If \(\mu_j\ne0\), then, for every length \(L\),
    \[
        G_L(A_j)\subseteq G_L(B).
    \]
\end{lemma}

\begin{proof}
    \leanok
    A vector in \(G_L(A_j)\) has the form \(\Gamma_L^{A_j}(X)\).  Put
    \(\mu_j^{-L}X\) in the \(j\)-th diagonal block of a boundary matrix for
    \(B\) and put zero in every other block.  Applying \(\Gamma_L^B\) gives
    back \(\Gamma_L^{A_j}(X)\).
\end{proof}

\begin{lemma}[Local ground space of a block-diagonal tensor]
    \label{lem:block_sum_local_ground_space}
    \lean{MPSTensor.groundSpace_toTensorFromBlocks_eq_iSup}
    \leanok
    \uses{def:ground_space_parent, lem:block_diagonal_boundary_condition_sum,
        lem:local_ground_space_block_le_assembled}
    Let
    \[
        B=\bigoplus_{j=1}^g\mu_jA_j,
        \qquad
        \mu_j\ne0\quad(1\le j\le g).
    \]
    Then, for every length \(L\),
    \[
        G_L(B)=\bigvee_{j=1}^gG_L(A_j).
    \]
\end{lemma}

\begin{proof}
    \leanok
    A boundary matrix \(X\) for \(B\) has diagonal blocks \(X_j\), and
    block-diagonality gives
    \[
        \Gamma_L^B(X)
        =
        \sum_{j=1}^g\Gamma_L^{A_j}(\mu_j^LX_j).
    \]
    This proves \(G_L(B)\subseteq\bigvee_jG_L(A_j)\). The reverse inclusion is
    Lemma~\ref{lem:local_ground_space_block_le_assembled} for each block \(j\).
    The displayed equality follows.
\end{proof}

\begin{lemma}[Periodic constraints for a block-diagonal tensor]
    \label{lem:block_diagonal_periodic_constraints_propagated_sum}
    \lean{MPSTensor.chainGroundSpace_toTensorFromBlocks_le_iSup_groundSpace}
    \leanok
    \uses{lem:block_sum_local_ground_space,
        lem:periodic_constraints_propagated_subspace_sequence}
    Let
    \[
        B=\bigoplus_{j=1}^g\mu_jA_j,
        \qquad \mu_j\ne0\quad(1\le j\le g).
    \]
    Set
    \[
        S_M:=\bigvee_{j=1}^gG_M(A_j).
    \]
    If, for every \(M\ge L\),
    \[
        \mathbb C^d\otimes S_M\cap S_M\otimes\mathbb C^d=S_{M+1},
    \]
    then
    \[
        \mathcal G_{N,L}(B)\subseteq S_N.
    \]
\end{lemma}

\begin{proof}
    \leanok
    Lemma~\ref{lem:block_sum_local_ground_space} gives \(G_L(B)=S_L\).
    Apply Lemma~\ref{lem:periodic_constraints_propagated_subspace_sequence}
    to the sequence \(S_M\).
\end{proof}

\begin{lemma}[PGVWC propagation for the block-diagonal tensor]
    \label{lem:bnt_block_diagonal_periodic_constraints_propagated_sum}
    \lean{MPSTensor.chainGroundSpace_toTensorFromBlocks_le_iSup_groundSpace_of_ge_of_bnt_directSum_unital_c1}
    \leanok
    \uses{lem:bnt_direct_sum_unital_block_intersection_ge,
        lem:block_diagonal_periodic_constraints_propagated_sum}
    Let \(A_1,\ldots,A_g\) be separated normalized BNT blocks. Suppose
    \(L_0>0\), and for every block \(j\) the map
    \[
        \Gamma_{L_0}^{A^j}:M_{D_j}(\mathbb C)\to
        (\mathbb C^d)^{\otimes L_0}
    \]
    is injective. Assume also the normalization
    \[
        \sum_a A^j_aA^{j\,\dagger}_a=I .
    \]
    For nonzero scalars \(\mu_j\), set
    \[
        B=\bigoplus_{j=1}^g\mu_jA_j,
        \qquad
        S_M:=\bigvee_{j=1}^gG_M(A_j).
    \]
    If
    \[
        L\ge (L_0+1)+(g-1)3(L_0+1)+1,
    \]
    then, for every \(N\ge L\),
    \[
        \mathcal G_{N,L}(B)\subseteq S_N.
    \]
\end{lemma}

\begin{proof}
    \leanok
    Put
    \[
        T=(L_0+1)+(g-1)3(L_0+1).
    \]
    For \(M\ge L\), one has \(M-1\ge T\). Applying
    Lemma~\ref{lem:bnt_direct_sum_unital_block_intersection_ge} at
    \(n=M-1\) gives
    \[
        \mathbb C^d\otimes S_M\cap S_M\otimes\mathbb C^d=S_{M+1}
    \]
    for every \(M\ge L\). Lemma~\ref{lem:block_diagonal_periodic_constraints_propagated_sum}
    then propagates the local constraint from \(L\) to \(N\).
\end{proof}

\begin{lemma}[PGVWC propagation with a direct local-space sum]
    \label{lem:bnt_block_diagonal_periodic_constraints_internal_sum}
    \lean{MPSTensor.chainGroundSpace_toTensorFromBlocks_le_iSup_and_iSupIndep_of_bnt_unital_c1}
    \leanok
    \uses{lem:bnt_block_diagonal_periodic_constraints_propagated_sum,
        lem:product_span_block_image_direct_sum}
    With the hypotheses and notation of
    Lemma~\ref{lem:bnt_block_diagonal_periodic_constraints_propagated_sum}, set
    \[
        S_N:=\bigvee_{j=1}^gG_N(A_j).
    \]
    Then, for every $N\ge L$,
    \[
        \mathcal G_{N,L}(B)\subseteq S_N,
    \]
    and the sum defining \(S_N\) is internal: if
    \[
        \phi_j\in G_N(A_j),
        \qquad
        \sum_{j=1}^g\phi_j=0,
    \]
    then $\phi_j=0$ for every $j$.
\end{lemma}

\begin{proof}
    \leanok
    The inclusion is
    Lemma~\ref{lem:bnt_block_diagonal_periodic_constraints_propagated_sum}. Since
    $N\ge L$ and
    \[
        L\ge (L_0+1)+(g-1)3(L_0+1)+1,
    \]
    the product-span directness range of
    Lemma~\ref{lem:product_span_block_image_direct_sum} applies at length
    $N$:
    \[
        \operatorname{span}\{(A^1_w,\ldots,A^g_w): |w|=N\}
        =
        \prod_j M_{D_j}(\mathbb C)
        \quad\Longrightarrow\quad
        \ker\!\left(
            \bigoplus_{j=1}^g G_N(A_j)
            \longrightarrow (\mathbb C^d)^{\otimes N},
            \quad
            (\phi_j)_j\longmapsto \sum_j\phi_j
        \right)=0.
    \]
\end{proof}

\begin{lemma}[Block-diagonal chain vectors decompose in the local block spaces]
    \label{lem:bnt_block_diagonal_open_boundary_decomposition}
    \lean{MPSTensor.exists_unique_sum_groundSpace_of_chainGroundSpace_toTensorFromBlocks_of_bnt_unital_c1}
    \leanok
    \uses{lem:bnt_block_diagonal_periodic_constraints_propagated_sum}
    With the hypotheses and notation of
    Lemma~\ref{lem:bnt_block_diagonal_periodic_constraints_propagated_sum}, every
    vector
    \[
        \psi\in
        \mathcal G_{N,L}\!\left(\bigoplus_{j=1}^g\mu_jA_j\right)
    \]
    has a unique family \((\psi_j)_{j=1}^g\) such that
    \[
        \psi=\sum_{j=1}^g\psi_j,
        \qquad
        \psi_j\in G_N(A_j).
    \]
\end{lemma}

\begin{proof}
    \leanok
    \uses{lem:bnt_block_diagonal_periodic_constraints_internal_sum}
    Lemma~\ref{lem:bnt_block_diagonal_periodic_constraints_internal_sum} gives
    \[
        \mathcal G_{N,L}\!\left(\bigoplus_{j=1}^g\mu_jA_j\right)
        \subseteq
        \bigvee_{j=1}^g G_N(A_j),
    \]
    and the spaces \(G_N(A_j)\) are independent: for every \(j\) and every
    \(x_j\in G_N(A_j)\),
    \[
        \sum_j x_j=0
        \quad\Longrightarrow\quad
        x_j=0.
    \]
    Thus every
    \(\psi\in\mathcal G_{N,L}(\bigoplus_j\mu_jA_j)\) has a unique
    decomposition \(\psi=\sum_j\psi_j\) with \(\psi_j\in G_N(A_j)\).
    The remaining boundary step in \cite[Theorem~2blocks.2]{PerezGarcia2007Matrix}
    is to prove
    \[
        \psi_j\in\mathcal G_{N,L}(A_j)
    \]
    for each \(j\), by closing the periodic boundary with block-diagonal
    boundary conditions.
\end{proof}

\begin{lemma}[Block-diagonal periodic chain inclusions]
    \label{lem:bnt_block_diagonal_periodic_chain_inclusions}
    \lean{MPSTensor.chainGroundSpace_toTensorFromBlocks_two_inclusions_and_iSupIndep_of_bnt_unital_c1}
    \leanok
    \uses{lem:isup_chain_ground_space_block_le_assembled,
        lem:bnt_block_diagonal_periodic_constraints_internal_sum}
    With the hypotheses and notation of
    Lemma~\ref{lem:bnt_block_diagonal_periodic_constraints_propagated_sum}, one has
    \[
        \bigvee_{j=1}^g\mathcal G_{N,L}(A_j)
        \subseteq
        \mathcal G_{N,L}\!\left(\bigoplus_{j=1}^g\mu_jA_j\right)
        \subseteq
        \bigvee_{j=1}^gG_N(A_j),
    \]
    and the sum defining the right-hand side is internal.
\end{lemma}

\begin{proof}
    \leanok
    The first inclusion is
    Lemma~\ref{lem:isup_chain_ground_space_block_le_assembled}. The second
    inclusion and internal directness are
    Lemma~\ref{lem:bnt_block_diagonal_periodic_constraints_internal_sum}. The
    claimed statement is their conjunction.
    % The remaining PGVWC07 boundary-closing step with block-diagonal boundary
    % conditions replaces
    % $\bigvee_jG_N(A_j)$ by $\sum_j\mathcal G_{N,L}(A_j)$.
\end{proof}

\begin{lemma}[Boundary closing gives the block-diagonal chain equality]
    \label{lem:block_diagonal_boundary_closing_equality}
    \lean{MPSTensor.chainGroundSpace_toTensorFromBlocks_eq_iSup_chainGroundSpace_of_boundary_closing}
    \leanok
    \uses{lem:isup_chain_ground_space_block_le_assembled}
    Let
    \[
        B=\bigoplus_{j=1}^g\mu_jA_j,
        \qquad \mu_j\ne0\quad(1\le j\le g).
    \]
    If closing the periodic boundary with block-diagonal boundary conditions gives
    \[
        \mathcal G_{N,L}(B)
        \subseteq
        \bigvee_{j=1}^g\mathcal G_{N,L}(A_j),
    \]
    then
    \[
        \mathcal G_{N,L}(B)
        =
        \bigvee_{j=1}^g\mathcal G_{N,L}(A_j).
    \]
\end{lemma}

\begin{proof}
    \leanok
    The reverse inclusion is
    Lemma~\ref{lem:isup_chain_ground_space_block_le_assembled}; combine the two
    inclusions.
\end{proof}

\begin{lemma}[Boundary decomposition gives the block-diagonal chain inclusion]
    \label{lem:block_diagonal_boundary_decomposition_inclusion}
    \lean{MPSTensor.chainGroundSpace_toTensorFromBlocks_le_iSup_of_boundary_decomposition}
    \leanok
    Let
    \[
        B=\bigoplus_{j=1}^g\mu_jA_j.
    \]
    Suppose that every vector $\psi\in\mathcal G_{N,L}(B)$ can be written as
    \[
        \psi=\sum_{j=1}^g\psi_j,
        \qquad
        \psi_j\in\mathcal G_{N,L}(A_j).
    \]
    Then
    \[
        \mathcal G_{N,L}(B)
        \subseteq
        \bigvee_{j=1}^g\mathcal G_{N,L}(A_j).
    \]
\end{lemma}

\begin{proof}
    \leanok
    For $\psi\in\mathcal G_{N,L}(B)$, choose the displayed decomposition. Each
    summand lies in the corresponding subspace
    $\mathcal G_{N,L}(A_j)$, hence their finite sum lies in
    $\bigvee_j\mathcal G_{N,L}(A_j)$.
\end{proof}

\begin{lemma}[Boundary decomposition in the finite injectivity range]
    \label{lem:bnt_block_diagonal_boundary_decomposition_equality}
    \lean{MPSTensor.chainGroundSpace_toTensorFromBlocks_eq_iSup_and_iSupIndep_of_bnt_c1_boundary_decomposition}
    \leanok
    \uses{lem:block_diagonal_boundary_decomposition_inclusion,
        lem:isup_chain_ground_space_block_le_assembled,
        lem:bnt_block_diagonal_periodic_constraints_internal_sum}
    Let $A_1,\ldots,A_g$ be separated normalized blocks in the basis of normal
    tensors, and let
    \[
        B=\bigoplus_{j=1}^g\mu_jA_j,
        \qquad \mu_j\ne0.
    \]
    Assume each block is injective at length $L_0$, assume
    $\sum_aA^j_aA^{j\dagger}_a=I$ for every $j$, and suppose
    \[
        L\ge (L_0+1)+(g-1)3(L_0+1)+1,
        \qquad
        N\ge L.
    \]
    If every vector $\psi\in\mathcal G_{N,L}(B)$ can be written as
    \[
        \psi=\sum_{j=1}^g\psi_j,
        \qquad
        \psi_j\in\mathcal G_{N,L}(A_j),
    \]
    then
    \[
        \mathcal G_{N,L}(B)
        =
        \bigvee_{j=1}^g\mathcal G_{N,L}(A_j),
    \]
    and the sum $\bigvee_jG_N(A_j)$ is internal.
\end{lemma}

\begin{proof}
    \leanok
    Lemma~\ref{lem:block_diagonal_boundary_decomposition_inclusion} gives
    \[
        \mathcal G_{N,L}(B)
        \subseteq
        \bigvee_{j=1}^g\mathcal G_{N,L}(A_j).
    \]
    The blockwise inclusion gives the reverse containment, hence equality. The
    internality of $\bigvee_jG_N(A_j)$ is the directness conclusion of
    Lemma~\ref{lem:bnt_block_diagonal_periodic_constraints_internal_sum}.
\end{proof}

\begin{theorem}[Block-diagonal intersection identity]
    \label{thm:block_diagonal_ground_space_intersection}
    \notready
    \uses{def:ground_space_parent, def:chain_ground_space, def:l_blk_injective,
        def:mpv, lem:blockwise_boundary_compatibility_from_traces,
        lem:normalized_boundary_matrix_identities,
        lem:left_boundary_summands_mem_block_ground_spaces,
        lem:left_boundary_trace_decomposition_mem_block_ground_spaces,
        lem:block_restrictions_mem_block_ground_spaces,
        lem:block_restrictions_iff_block_ground_spaces,
        lem:block_restriction_intersection_subspace,
        lem:period_window_block_restriction_intersection,
        lem:bnt_direct_sum_block_separation_block_intersection,
        lem:bnt_direct_sum_block_separation_period_window_block_intersection,
        lem:unital_block_separating_product_span,
        lem:product_span_block_image_direct_sum,
        lem:bnt_direct_sum_unital_block_intersection_ge,
        lem:bnt_direct_sum_unital_block_intersection,
        lem:block_sum_local_ground_space,
        lem:block_diagonal_periodic_constraints_propagated_sum,
        lem:bnt_block_diagonal_periodic_constraints_propagated_sum,
        lem:bnt_block_diagonal_periodic_chain_inclusions,
        lem:bnt_block_diagonal_open_boundary_decomposition,
        lem:block_diagonal_boundary_closing_equality,
        lem:block_diagonal_boundary_decomposition_inclusion,
        lem:bnt_block_diagonal_boundary_decomposition_equality,
        thm:wordspan_propagation_unital}
    Let $A_1,\ldots,A_g$ be the blocks in a block-injective canonical form, and
    assume that their MPV families are pairwise different:
    \[
        j\ne k
        \quad\Longrightarrow\quad
        \exists n\ge1,\; V^{(n)}(A_j)\ne V^{(n)}(A_k).
    \]
    Let $L_0$ be a common injectivity length. In the multi-block range used in
    \cite[Theorem~2blocks.2]{PerezGarcia2007Matrix}, namely
    \[
        g\ge2,
        \qquad
        L_0\ge1,
        \qquad
        L\ge 3(g-1)(L_0+1)+1,
    \]
    the block-diagonal intersection property says that, for every $m\ge L$,
    \[
        \mathbb C^d\otimes
        \left(\bigoplus_{j=1}^g G_m(A_j)\right)
        \cap
        \left(\bigoplus_{j=1}^g G_m(A_j)\right)\otimes\mathbb C^d
        =
        \bigoplus_{j=1}^g G_{m+1}(A_j).
    \]
    For nonzero weights $\mu_j$ and every $N\ge L+L_0$, the corresponding
    periodic chain ground spaces satisfy
    \[
        \mathcal G_{N,L}\!\left(\bigoplus_{j=1}^g \mu_j A_j\right)
        =
        \sum_{j=1}^g \mathcal G_{N,L}(A_j).
    \]
\end{theorem}
% This source-level structural theorem is intentionally separated from the
% following conditional reductions, which assume the needed block-diagonal
% boundary representation separately.

\begin{proof}
    \notready
    The source proof takes a vector $\psi$ in the left-hand side and writes its
    two boundary descriptions as
    \[
        \psi=\sum_{j=1}^g \alpha_j=\sum_{j=1}^g \beta_j,
        \qquad
        \alpha_j\in\mathbb C^d\otimes G_m(A_j),
        \qquad
        \beta_j\in G_m(A_j)\otimes\mathbb C^d .
    \]
    In the source notation, the two components have boundary matrices
    $C^j_{i_1}$ and $D^j_{i_{m+1}}$: the first family leaves the initial
    physical index open, and the second family leaves the final physical index
    open.
    \[
        \alpha_j
        =
        \sum_{i_1,\ldots,i_{m+1}}
        \tr\!\left(
            A^j_{i_{m+1}} C^j_{i_1}
            A^j_{i_2}\cdots A^j_{i_m}
        \right)
        |i_1\cdots i_{m+1}\rangle,
    \]
    \[
        \beta_j
        =
        \sum_{i_1,\ldots,i_{m+1}}
        \tr\!\left(
            D^j_{i_{m+1}}
            A^j_{i_1} A^j_{i_2}\cdots A^j_{i_m}
        \right)
        |i_1\cdots i_{m+1}\rangle.
    \]
    The direct-sum separation lemma used in
    \cite[Theorem~2blocks.2]{PerezGarcia2007Matrix}, together with injectivity
    of each block, gives the blockwise matrix identity
    \[
        A^j_{i_{m+1}} C^j_{i_1}
        =
        D^j_{i_{m+1}} A^j_{i_1}
    \]
    for every $j$, $i_1$, and $i_{m+1}$. For fixed $j$, set
    \[
        E^j:=\sum_a C^j_a A^{j\,\dagger}_a .
    \]
    In the normalization used in the source proof,
    $\sum_a A^j_a A^{j\,\dagger}_a=\Id$, so
    Lemma~\ref{lem:normalized_boundary_matrix_identities} gives
    \[
        D^j_b
        =
        \sum_a D^j_b A^j_a A^{j\,\dagger}_a
        =
        \sum_a A^j_b C^j_a A^{j\,\dagger}_a
        =
        A^j_b E^j .
    \]
    Hence
    \[
        A^j_b C^j_a
        =
        A^j_b E^j A^j_a
    \]
    for every $a,b$, and cyclicity of the trace gives
    \[
        \alpha_j
        =
        \sum_{i_1,\ldots,i_{m+1}}
        \tr\!\left(A^j_{i_1}\cdots A^j_{i_m}
                  A^j_{i_{m+1}}E^j\right)
        |i_1\cdots i_{m+1}\rangle
        =
        \Gamma_{m+1}^{A_j}(E^j)
        \in G_{m+1}(A_j).
    \]
    Since $\psi=\sum_j\alpha_j$, this gives
    \[
        \psi\in\bigoplus_j G_{m+1}(A_j).
    \]
    For one block the preceding implication is
    \[
        \mathbb C^d\otimes G_m(A_j)
        \cap
        G_m(A_j)\otimes\mathbb C^d
        \subseteq
        G_{m+1}(A_j).
    \]
    The reverse inclusion is blockwise:
    \[
        G_{m+1}(A_j)
        \subseteq
        \mathbb C^d\otimes G_m(A_j)
        \cap
        G_m(A_j)\otimes\mathbb C^d .
    \]
    Thus
    \[
        \mathbb C^d\otimes G_m(A_j)
        \cap
        G_m(A_j)\otimes\mathbb C^d
        =
        G_{m+1}(A_j).
    \]
    With
    \[
        S_m:=\bigoplus_{j=1}^g G_m(A_j),
    \]
    the preceding calculation is the open-segment recursion
    \[
        \mathbb C^d\otimes S_m
        \cap
        S_m\otimes\mathbb C^d
        =
        S_{m+1}
    \]
    for every $m\ge L$. For $B=\bigoplus_j\mu_j A_j$ with every
    $\mu_j\ne0$, Lemma~\ref{lem:block_sum_local_ground_space} gives
    \[
        G_L(B)=\bigvee_jG_L(A_j).
    \]
    In the range where the sums \(\bigvee_jG_m(A_j)\) are direct, this identifies
    $G_L(B)$ with the subspace denoted in the source by
    $S=\bigoplus_j\mathcal G_L^{A^j}$.
    The passage from the open-segment recursion to the periodic chain is the
    separate closure step in \cite[Theorem~2blocks.2]{PerezGarcia2007Matrix}. It is
    not the inclusion $G_N(A_j)\subseteq\mathcal G_{N,L}(A_j)$, which fails
    for a general open-boundary matrix crossing the periodic boundary. Under
    $N\ge L+L_0$, the required boundary-condition comparison is
    \[
        \mathcal G_{N,L}(B)
        =
        \sum_{j=1}^g \mathcal G_{N,L}(A_j).
    \]
\end{proof}

\begin{lemma}[Conditional reduction to block-diagonal boundary conditions]
    \label{lem:conditional_boundary_matrix_block_split_projection_span}
    \notready
    \uses{lem:block_diagonal_commutant_reduction_from_projection_span,
        lem:eval_word_block_diagonal_tensor, lem:mpv_mem_chain_ground_space,
        lem:isup_chain_ground_space_block_le_assembled,
        thm:chain_ground_space_eq_mpv_blk, lem:center_scalar}
    Let $A^i=\bigoplus_j\mu_j A_j^i$, and assume $1<L$, $N \ge L+1$, and
    $0<m$. Suppose that every virtual sector projection $P_j$ lies in the
    pulled-back span of the length-$m$ words $A^\omega$. If a boundary matrix
    $X$ commutes with all those words, then the commutant reduction gives
    \[
        X=\bigoplus_j X_j.
    \]
    Therefore the vector obtained by applying the ground-space map to $X$ lies in
    $\sum_j \mathcal G_{N,L}(A_j)$. Consequently, if every vector in
    $\mathcal G_{N,L}\!\left(\bigoplus_j\mu_j A_j\right)$ admits such a
    commuting boundary matrix representation, as in the direct-sum analogue of
    Theorem~\ref{thm:boundary_matrix_commutes}, then
    \[
        \mathcal G_{N,L}\!\left(\bigoplus_j\mu_j A_j\right)
        \subseteq \sum_j \mathcal G_{N,L}(A_j),
    \]
    and together with the already-proved forward inclusion this gives equality.
    Composing the reverse inclusion with blockwise injective uniqueness gives
    the conditional containment in the vector space generated by the basis of
    normal tensors.  This is a conditional form of the block-diagonal
    boundary-condition step in \cite[lines~2126--2128]{Cirac2021Matrix}.
\end{lemma}
% The ground-space theorem below uses the conservative length range of
% the theorem titled "Degeneracy of the ground space v2" in
% \cite{PerezGarcia2007Matrix}.

\begin{proof}
    \notready
    \uses{lem:block_diagonal_commutant_reduction_from_projection_span,
        lem:eval_word_block_diagonal_tensor, lem:mpv_mem_chain_ground_space,
        lem:isup_chain_ground_space_block_le_assembled,
        thm:chain_ground_space_eq_mpv_blk, lem:center_scalar}
    The commutant reduction makes the pulled-back boundary matrix block
    diagonal. Restricting the same word-commutation identities to a diagonal
    block and using the injectivity of that block shows, by the scalar-commutant
    lemma, that each diagonal block is scalar. Write $\Gamma_N^B$ for the
    length-$N$ boundary-to-state map of the tensor $B$. For a block-diagonal
    boundary matrix,
    \[
        \Gamma_N^{\oplus_j\mu_j A_j}\!\left(\bigoplus_j X_j\right)
        =
        \sum_j \mu_j^N \Gamma_N^{A_j}(X_j).
    \]
    If $X_j=\lambda_j I$, this is a finite sum of scalar multiples of the
    vectors $V^{(N)}(A_j)$. The reverse inclusion and equality statements are
    reformulations of this calculation, combined with the forward inclusion and
    the blockwise uniqueness theorem.
\end{proof}

\begin{lemma}[Block-diagonal boundary reduction from blockwise word span]
    \label{lem:conditional_block_split_from_product_word_span}
    \notready
    \uses{def:blockwise_product_word_span}
    Let $A^i=\bigoplus_j\mu_j A_j^i$ be in block-injective canonical form, and
    assume $1<L$, $N \ge L+1$, and $0<m$. If
    \[
        \spn\{(A_1^\omega,\ldots,A_g^\omega): |\omega|=m\}
        =
        \prod_j \MN{D_j}
    \]
    and
    every vector in
    $\mathcal G_{N,L}\!\left(\bigoplus_j \mu_j A_j\right)$ comes with a boundary
    matrix representation commuting with all length-$m$ direct-sum words, then
    \[
        \mathcal G_{N,L}\!\left(\bigoplus_j \mu_j A_j\right)
        \subseteq \sum_j \mathcal G_{N,L}(A_j).
    \]
    Combined with the forward inclusion, this gives the corresponding conditional
    equality. In the parent-ground-space formulation, the same hypotheses imply
    containment in $\spn\{V^{(N)}(A_j): j=1,\ldots,g\}$. The blockwise word span
    and boundary representation are the displayed hypotheses.
\end{lemma}

\begin{proof}
    \notready
    \uses{lem:block_projection_span_from_product_word_span,
        lem:conditional_boundary_matrix_block_split_projection_span}
    The blockwise word-span hypothesis gives
    \[
        P_j\in
        \spn\left\{
        \bigl(\bigoplus_r\mu_r A_r\bigr)^\omega:|\omega|=m
        \right\}.
    \]
    If $X$ commutes with all displayed word products, then
    $XP_j=P_jX$ for every $j$, hence $X=\bigoplus_jX_j$. The direct-sum
    boundary formula
    \[
        \Gamma_N^{\oplus_j\mu_j A_j}\!\left(\bigoplus_jX_j\right)
        =
        \sum_j\mu_j^N\Gamma_N^{A_j}(X_j)
    \]
    gives the reverse inclusion, and the blockwise uniqueness theorem gives the
    stated span containment.
\end{proof}

\begin{lemma}[Block-diagonal boundary reduction from block-separating equations]
    \label{lem:conditional_block_split_from_bnt_block_separation}
    \notready
    \uses{def:bnt, rem:word_tuple_span_top,
        def:parent_hamiltonian_ground_space_bnt, def:bnt_span}
    Let $A^i=\bigoplus_j\mu_j A_j^i$ be in block-injective canonical form, with
    $A_1,\ldots,A_g$ a basis of normal tensors.  Assume $1 < L$ and
    $N \ge L + 1$, and suppose
    that for each block $k$ there are coefficients $c_\omega^{(k)}$ such that
    \[
        \sum_{|\omega|=S} c_\omega^{(k)} A_j^\omega
        =
        \begin{cases}
            I, & j=k,\\
            0, & j\ne k.
        \end{cases}
    \]
    If every vector in
    $\mathcal G_{N,L}\!\left(\bigoplus_j \mu_j A_j\right)$ has a boundary matrix
    representation commuting with all length-$(1 + S)$ direct-sum words, then
    \[
        \mathcal G_{N,L}\!\left(\bigoplus_j \mu_j A_j\right)
        \subseteq \sum_j \mathcal G_{N,L}(A_j).
    \]
    Together with the forward inclusion, this gives the corresponding
    conditional equality. In the parent-ground-space formulation, the same
    hypotheses imply containment in $\spn\{V^{(N)}(A_j): j=1,\ldots,g\}$. This
    is the block-injective ground-space argument of
    \cite[lines~2120--2128]{Cirac2021Matrix}, where the re-growing step is
    restricted to block-diagonal boundary conditions.
\end{lemma}

\begin{proof}
    \notready
    \uses{lem:one_block_injective_of_injective,
        lem:product_word_span_from_bnt_block_separation,
        lem:conditional_block_split_from_product_word_span}
    Finite-length injectivity gives a block length at which each block has the
    required physical--virtual correspondence. In the current length-one
    specialization, Lemma~\ref{lem:one_block_injective_of_injective} supplies
    the corresponding one-block injectivity. The displayed separation equations imply
    \[
        \spn\{(A_1^\omega,\ldots,A_g^\omega):|\omega|=1+S\}
        =
        \prod_j\MN{D_j}.
    \]
    This is Lemma~\ref{lem:product_word_span_from_bnt_block_separation} with $L=1$.
    Therefore Lemma~\ref{lem:conditional_block_split_from_product_word_span}
    applies with $m=1+S$. Its conclusion is the displayed containment, and the
    equality follows after adding the forward inclusion.
\end{proof}

\begin{theorem}[Containment in the space generated by a basis of normal tensors]
    \label{thm:parent_ground_space_le_bnt_span}
    \notready
    \uses{def:parent_hamiltonian_ground_space_bnt, def:bnt_span,
        thm:block_diagonal_ground_space_intersection,
        thm:chain_ground_space_eq_mpv_normal}
    If $A^i=\bigoplus_j\mu_j A_j^i$ is in block-injective canonical form, where
    $A_1,\ldots,A_g$ is a basis of normal tensors, assume $g\ge2$ and
    $\mu_j\ne0$ for every $j$. Let $L_0\ge1$ be a common injectivity length for
    the blocks, and suppose
    \[
        L_0<L,
        \qquad
        L\ge 3(g-1)(L_0+1)+1,
        \qquad
        N\ge L+L_0,
    \]
    then
    \[
        \mathcal G_{N,L}\!\left(\bigoplus_j \mu_j A_j\right)
        \subseteq
        \spn\{V^{(N)}(A_j): j=1,\ldots,g\}.
    \]
\end{theorem}

\begin{proof}
    \notready
    The length hypotheses are the conservative range used in
    \cite[Theorem~2blocks.2]{PerezGarcia2007Matrix}. The shorter range stated in
    \cite[lines~2126--2128]{Cirac2021Matrix} requires the strengthened
    block-diagonal re-growing step. Since $g\ge2$ and $L_0\ge1$,
    \[
        L\ge 3(g-1)(L_0+1)+1 \ge 3L_0+4 > L_0,
        \qquad
        N\ge L+L_0\ge L+1,
    \]
    Because $A_1,\ldots,A_g$ is a basis of normal tensors, the block MPV
    families are non-redundant:
    \[
        j\ne k
        \quad\Longrightarrow\quad
        \exists n\ge1,\; V^{(n)}(A_j)\ne V^{(n)}(A_k).
    \]
    Thus the ambient hypotheses of the block-diagonal intersection theorem
    are available. Theorem~\ref{thm:block_diagonal_ground_space_intersection}
    gives the periodic decomposition
    \[
        \mathcal G_{N,L}\!\left(\bigoplus_j \mu_j A_j\right)
        =
        \sum_j \mathcal G_{N,L}(A_j).
    \]
    Theorem~\ref{thm:chain_ground_space_eq_mpv_normal} gives
    \[
        \mathcal G_{N,L}(A_j)
        =
        \spn\{V^{(N)}(A_j)\}
    \]
    for every $j$. Substituting these one-block identities into the preceding
    sum gives the displayed containment.
\end{proof}

\begin{theorem}[Chain ground space generated by a basis of normal tensors]
    \label{thm:gs_eq_bnt_span}
    \notready
    \uses{thm:parent_ground_space_le_bnt_span, lem:bnt_mpv_mem_ground_space,
        def:parent_hamiltonian_ground_space_bnt, def:bnt_span}
    If $A^i=\bigoplus_j\mu_j A_j^i$ is in block-injective canonical form, where
    $A_1,\ldots,A_g$ is a basis of normal tensors, assume $g\ge2$ and
    $\mu_j\ne0$ for every $j$. Let $L_0\ge1$ be a common injectivity length for
    the blocks, and suppose
    \[
        L_0<L,
        \qquad
        L\ge 3(g-1)(L_0+1)+1,
        \qquad
        N\ge L+L_0,
    \]
    then
    \[
        \mathcal G_{N,L}\!\left(\bigoplus_j \mu_j A_j\right)
        =
        \spn\{V^{(N)}(A_j): j=1,\ldots,g\}.
    \]
\end{theorem}

\begin{proof}
    \notready
    We prove the two inclusions separately. The inclusion
    \[
        \mathcal G_{N,L}\!\left(\bigoplus_j \mu_j A_j\right)
        \subseteq
        \spn\{V^{(N)}(A_j): j=1,\ldots,g\}
    \]
    is exactly Theorem~\ref{thm:parent_ground_space_le_bnt_span}. For the reverse inclusion,
    the nonzero-weight hypothesis and Lemma~\ref{lem:bnt_mpv_mem_ground_space}
    show that each vector $V^{(N)}(A_j)$ lies in
    $\mathcal G_{N,L}(\bigoplus_j \mu_j A_j)$. Since the chain ground space is a
    linear subspace, it contains every linear combination
    of such vectors, and hence contains
    $\spn\{V^{(N)}(A_j): j=1,\ldots,g\}$. Thus
    \[
        \spn\{V^{(N)}(A_j): j=1,\ldots,g\}
        \subseteq
        \mathcal G_{N,L}\!\left(\bigoplus_j \mu_j A_j\right),
    \]
    and the two spaces are equal.
\end{proof}

\noindent\emph{Remark.} A separate kernel-identification statement for these direct-sum
block families would convert the chain-ground-space statement into the
corresponding statement for $\ker(H_N)$.
